\documentclass[fontset=sikou,punct=kaiming]{ctexbook}
\usepackage{xeCJKfntef}
\setcounter{secnumdepth}{4}
\ctexset{
	subsubsection=
	{
		numbering = true,
		name = {,、},
		number = \chinese{subsubsection},
		aftername = {\hspace{0pt}}
	},
}
\usepackage[a4paper,inner=1.5cm,textwidth=46\ccwd,top=2cm,bottom=2cm]{geometry}
\setlength{\headheight}{13pt}
\usepackage{fancyhdr}
\usepackage{lastpage}
\usepackage{calc,xcolor}
\newcommand{\progressbar}[2]{%
  \textcolor{black}{\rule{0.6\textwidth * \real{#1} / \real{#2}}{1.5ex}}%
	\textcolor{black!30}{\rule{0.6\textwidth - 0.6\textwidth * \real{#1} / \real{#2}}{1.5ex}}%
}

\pagestyle{fancy}
\fancyhf{}
\fancyhead[C]{\sffamily \booktitle}
\fancyhead[ER]{\sffamily \leftmark}
\fancyhead[OL]{\sffamily \rightmark}
\fancyhead[EL,OR]{\sffamily\thepage}

\fancypagestyle{plain}{%
	\fancyhf{}
	\cfoot{\sffamily\thepage}
	\renewcommand\headrulewidth{0pt}
	\renewcommand\footrulewidth{0pt}
}
\makeatletter
\fancypagestyle{myfancy}{
	\pagestyle{fancy}
	\fancyhead[C]{\sffamily \booktitle}
	\fancyhead[ER]{\sffamily \leftmark}
	\fancyhead[OL]{\sffamily \rightmark}
	\fancyhead[EL,OR]{\sffamily\thepage}
	\fancyfoot[EC]{%
		\@ifundefined{lastpage@lastpage}{%
		\progressbar{\thepage}{100}%
		}{%
		\progressbar{\thepage}{\lastpage@lastpage}%
		}%
	}
	\fancyfoot[OC]{\sffamily\fbox{仅供学习使用，严禁商业使用。此版本为第~\version{} 版，若想查看是否为最新版，去看 QQ: 2098881095 的空间}}
}
\makeatother
\usepackage[toc]{multitoc}
\usepackage{amsmath,amsthm}
\usepackage[partial=upright]{unicode-math}
\usepackage{physics,siunitx}
\setmainfont{XITS}
\setmathfont[StylisticSet = 8]{XITS Math}
\usepackage{sourcesanspro,sourcecodepro}
\newtheoremstyle{titi}
	{5pt}%
	{5pt}%
	{}%
	{0pt}%
	{\bfseries}%
	{}%
	{0pt}%
	{\thmname{#1}\thmnumber{#2}\thmnote{（#3）}}%
\theoremstyle{titi}
\newtheorem{ti}{}[subsubsection]
\renewcommand{\theti}{\arabic{ti}．}
\newtheorem*{solution}{解}
% \usepackage{ntheorem}
% {
% 	\theoremstyle{nonumberplain}
% 	\theoremheaderfont{\bfseries}
% 	\theorembodyfont{\normalfont}
% 	\newtheorem{solution}{解}
% }
% {
% 	\theoremstyle{change}
% 	\theoremheaderfont{\bfseries}
% 	\theorembodyfont{\normalfont}
% 	\newtheorem{ti}{}[subsubsection]
% }
% \renewcommand{\theti}{\arabic{ti}.}
\usepackage{tabularx,multirow}
\newcolumntype{Z}{>{\centering\arraybackslash}X}
\usepackage[hidelinks,bookmarksopen=true,bookmarksnumbered=true]{hyperref}
\usepackage{bookmark}
\newcommand{\ee}{\symup e}
\newcommand{\ii}{\symup i}
\DeclareMathOperator{\Ln}{Ln}
\newcommand{\LL}{\symscr{L}}
\newcommand{\EE}{\symbb{E}\,}
\newcommand{\bA}{\symbfit A}
\newcommand{\bB}{\symbfit B}
\newcommand{\bC}{\symbfit C}
\newcommand{\bD}{\symbfit D}
\newcommand{\bE}{\symbfit E}
\newcommand{\bX}{\symbfit X}
\newcommand{\astt}{*}
\newcommand{\balpha}{\symbfit \alpha}
\newcommand{\bbeta}{\symbfit \beta}
\newcommand{\TT}{\mathrm{T}}
\renewcommand{\leq}{\leqslant}
\renewcommand{\geq}{\geqslant}
\DeclareMathOperator{\Cov}{Cov}
\DeclareMathOperator{\Prj}{Prj}
% \makeatletter
% \let\mycases={\cases}
% \newenvironment{mycases}{%
%   \matrix@check\mycases\let\@ifnextchar\new@ifnextchar%
%   \left\lbrace%
%   \def\arraystretch{1.2}%
%   \array{@{}l}%
% }{%
%   \endarray\right.%
% }
% \makeatother
\def\theenumi{\arabic{enumi}}
\def\labelenumi{(\theenumi)}
\def\kuo{\mbox{(\hspace{\ccwd})}}
\def\kuoA{(\nobreak\makebox[\ccwd]{A}\nobreak)}
\def\kuoB{(\nobreak\makebox[\ccwd]{B}\nobreak)}
\def\kuoC{(\nobreak\makebox[\ccwd]{C}\nobreak)}
\def\kuoD{(\nobreak\makebox[\ccwd]{D}\nobreak)}
\def\hua{\CJKunderline*[hidden = true]{瞻彼阕者虚室生白}}
\def\huaa#1{\underline{ #1 }}
\long\def\guanggao{
	\vspace*{\fill}
	\begin{center}\thispagestyle{empty}
		广告位招租
	\end{center}
	\vspace*{\fill}
}
\setlength{\lineskip}{2.5pt}
\setlength{\lineskiplimit}{2.5pt}
\renewcommand\arraystretch{1.2}

\usepackage{calc,ifthen}

\newlength{\choicelengtha}
\newlength{\choicelengthb}
\newlength{\choicelengthc}
\newlength{\choicelengthd}
\newlength{\choicelengthe}
\newlength{\maxlength}

\makeatletter
\newcommand{\fourch}[4]{
  \par
  \settowidth{\choicelengtha}{A．#1}
  \settowidth{\choicelengthb}{B．#2}
  \settowidth{\choicelengthc}{C．#3}
  \settowidth{\choicelengthd}{D．#4}
  \ifthenelse{\lengthtest{\choicelengtha>\choicelengthb}}{\setlength{\maxlength}{\choicelengtha}}{\setlength{\maxlength}{\choicelengthb}}
  \ifthenelse{\lengthtest{\choicelengthc>\maxlength}}{\setlength{\maxlength}{\choicelengthc}}{}
  \ifthenelse{\lengthtest{\choicelengthd>\maxlength}}{\setlength{\maxlength}{\choicelengthd}}{}
  \ifthenelse{\lengthtest{\maxlength>0.8\linewidth}}
  {%
    \noindent%
    \begin{tabular}{@{}p{\linewidth}@{}}
      \setlength\tabcolsep{0pt}
      \@hangfrom{\textsf A．}#1 \\
      \@hangfrom{\textsf B．}#2 \\
      \@hangfrom{\textsf C．}#3 \\
      \@hangfrom{\textsf D．}#4 \\
    \end{tabular}
  }%
  {%
    \ifthenelse{\lengthtest{\maxlength>0.22\linewidth}}
    {%
      \noindent%
      \begin{tabular}{@{}p{0.48\linewidth}@{\hspace*{0.04\linewidth}}p{0.48\linewidth}@{}}
        \setlength\tabcolsep{0pt}
        \@hangfrom{\textsf A．}#1 & \@hangfrom{\textsf B．}#2 \\
        \@hangfrom{\textsf C．}#3 & \@hangfrom{\textsf D．}#4 \\
      \end{tabular}
    }%
    {%
      \noindent%
      \begin{tabular}{@{}*{3}{p{0.22\linewidth}@{\hspace*{0.04\linewidth}}}p{0.22\linewidth}@{}}
        \setlength\tabcolsep{0pt}
        \@hangfrom{\textsf A．}#1  & \@hangfrom{\textsf B．}#2 & \@hangfrom{\textsf C．}#3 & \@hangfrom{\textsf D．}#4 \\
      \end{tabular}
    }%
  }%
  \unskip\unskip
}
\makeatother
\def\onech{\fourch}
\def\twoch{\fourch}
\def\booktitle{JXUST 工科数学复习宝典}
\def\version{1.04}

\title{\booktitle（第~\version{} 版）}
\author{电气 173 何骏炜}
\begin{document}
	\frontmatter
	\pdfbookmark{标题页}{title}
	\maketitle
	\chapter{声明}
本汇总不得用于商业用途, 最新版下载地址： \href{https://github.com/sikouhjw/Mathematical-rearrangement}{Github} (用电脑点击链接即可), 不保证题目、答案的正确性, 如有错误可通过QQ群(见图~\ref{fig:1}~)\footnote{991832226}或者邮箱\footnote{489765924@qq.com}联系我们
\begin{figure}[htbp]
	\centering
	\includegraphics[width=0.3\textwidth]{2weima.png}
	\caption{二维码}\label{fig:1}
\end{figure}

点击 \href{https://github.com/sikouhjw/Mathematical-rearrangement}{Github} 后，找到 $\mathrm{gongkeshuxuechongpai.pdf}$ 后点击，点击 $\mathrm{Download}$ 即可

	\tableofcontents
	\bookmark[dest=\HyperLocalCurrentHref, level=chapter]{目录}
	\newpage
	\guanggao
	\mainmatter
	\pagestyle{myfancy}
	\chapter{《高等数学》试卷汇总}
\section{《高等数学（一）》期中}
\subsection{2018-2019 A7}
\subsubsection{选择题}
\begin{ti}
	$\lim_{x \to 0} x \cos \frac{1}{x} = $\kuo
	\fourch{$-1$}{$1$}{$0$}{不存在}
\end{ti}

\begin{ti}
	$\lim_{x \to x_{0}} f(x_{0}) = 0$ 及\kuo, 则 $\lim_{x \to x_{0}} f(x) g(x) = 0$
	\twoch{$g(x)$ 为任意函数时}{当 $g(x)$ 为有界函数时}{仅当 $\lim_{x \to x_{0}} g(x) = 0$ 时}{仅当 $\lim_{x \to x_{0}} g(x)$ 存在时}
\end{ti}

\begin{ti}
	函数 $f(x)$ 在点 $x_{0}$ 处有定义, 是 $f(x)$ 在该点处连续的\kuo
	\fourch{必要条件}{充分条件}{充要条件}{无关条件}
\end{ti}

\begin{ti}
	已知 $y = \sin x$, 则 $y^{(10)} = $\kuo
	\fourch{$- \sin x$}{$\sin x$}{$- \cos x$}{$\cos x$}
\end{ti}

\begin{ti}
	函数 $y = x + \arctan x$ 在 $(-\infty,+\infty)$ 上\kuo
	\fourch{不连续}{连续不可导}{单调递减}{单调递增}
\end{ti}

\begin{ti}
	设 $f(x) = \begin{cases}
		\frac{2}{3}x^{3}, & x \leqslant 1 \\
		x^{2}, & x > 1
	\end{cases}$, 则 $f(x)$ 在点 $x = 1$ 处的\kuo
	\twoch{左、右导数都存在}{左导数存在, 右导数不存在}{左导数不存在, 右导数存在}{左、右导数均不存在}
\end{ti}

\begin{ti}
	$f'(x_{0}) = 0, f''(x_{0}) > 0$ 是函数 $f(x)$ 在点 $x = x_{0}$ 处取得极小值的一个\kuo
	\twoch{充要条件}{充分不必要条件}{必要不充分条件}{既不充分也不必要条件}
\end{ti}

\begin{ti}
	$f(x)$ 在 $(-\infty,+\infty)$ 内可导, 且 $\forall x_{1},x_{2}$, 当 $x_{1} > x_{2}$ 时 $f(x_{1}) > f(x_{2})$, 则\kuo
	\fourch{任意 $x$, $f'(x) > 0$}{任意 $x$, $f'(-x) \leqslant 0$}{$f(-x)$ 单调递增}{$-f(-x)$ 单调递增}
\end{ti}

\begin{ti}
	指出曲线 $y = \frac{x}{3 - x}$ 的渐近线\kuo
	\twoch{没有水平渐近线, 也没有垂直渐近线}{有垂直渐近线, 没有水平渐近线}{既有垂直渐近线, 又有水平渐近线}{没有垂直渐近线, 只有水平渐近线}
\end{ti}

\subsubsection{填空题}
\begin{ti}
	$\lim_{x \to 0} \frac{2x - \sin x}{x^{4} - x} = $ \hua
\end{ti}

\begin{ti}
	当 $x \to 0$, $\sin 3x$ 是 $\ee^{x} - 1$ 的 \hua 阶无穷小
\end{ti}

\begin{ti}
	函数 $f(x) = \begin{cases}
		\sqrt{1 + x} - 1, & x \ne 0 (x \geqslant -1) \\
		k, & x = 0
	\end{cases}$ 在 $x = 0$ 处连续, 则 $k = $ \hua 
\end{ti}

\begin{ti}
	设 $f(x) = \frac{3}{5 - x} + \frac{x^{2}}{5}$, 则 $f'(0) = $ \hua
\end{ti}

\begin{ti}
	曲线 $\begin{cases}
		x = 2 \ee^{t} \\
		y = \ee^{-t}
	\end{cases}$ 在 $t = 0$ 处切线方程为\hua
\end{ti}

\begin{ti}
	设 $f(x)$ 一阶可导, $y = f(1 + \sin x)$, 则 $y' = $ \hua
\end{ti}

\begin{ti}
	若 $f(x)$ 在 $[a,b]$ 上连续, 在 $(a,b)$ 内可导, 则至少存在一点 $\xi \in (a,b)$ 使得 $f(b) - f(a) = $ \hua
\end{ti}

\begin{ti}
	函数 $f(x) = x + 2 \cos x$ 在 $\bigl[ 0, \frac{\uppi}{2} \bigr]$ 上的最大值为 \hua
\end{ti}

\begin{ti}
	曲线 $y = x^{3} - 3 x^{2} + 5$ 的拐点为 \hua
\end{ti}

\subsubsection{大题}
\begin{ti}
	求 $\lim_{x \to \infty} \frac{4 x^{4} - 3 x^{2} + 1}{2 x^{4} + 5 x^{2} - 6}$
\end{ti}

\begin{ti}
	求 $\lim_{x \to 0} \bigl( 1 + 2x^{2} \bigr)^{x - 2}$
\end{ti}

\begin{ti}
	设 $y = \ee^{\arctan \sqrt{x}}$, 求 $y'$ 及 $\dd{y}$
\end{ti}

\begin{ti}
	求由方程 $x + y - \ee^{2x} + \ee^{y} = 0$ 确定的隐函数导数 $\frac{\dd{y}}{\dd{x}}$
\end{ti}

\begin{ti}
	设 $f(x) = \begin{cases}
		2x + 1, & x < 0, \\
		x^{2} + 1, & x \geqslant 0,
	\end{cases}$ 讨论 $f(x)$ 在 $x = 0$ 处的连续性与可导性
\end{ti}

\begin{ti}
	求 $\lim_{x \to 1} \frac{x^{3} - 1 + \ln x}{\ee^{x} - \ee}$
\end{ti}

\begin{ti}
	讨论函数 $y = x \ee^{x}$ 的图形性态并绘图
\end{ti}

\section{《高等数学（一）》期末}
\subsection{2018-2019 A15}
\subsubsection{选择题}
\begin{ti}
	设 $f(x)=
	\begin{cases}
	\frac{x^2-3x+2}{x-2}, & x\ne2\\
	a, & x=2
	\end{cases}
	$ 为连续函数, 则 $a=$ \kuo
	\fourch{任意}{0}{2}{1}
\end{ti}

\begin{ti}
	设 $f(x)=
	\begin{cases}
	x^2\sin\frac{1}{x}, & x\ne0\\
	0, & x=0
	\end{cases}
	$, 则 $f'(0)=$ \kuo
	\fourch{1}{$-1$}{0}{不存在}
\end{ti}

\begin{ti}
	函数 $f(x)$ 有连续二阶导数, 且 $f(0)=0$, $f'(0)=1$, $f''(0)=-2$, 则 $\lim_{x\to0}\frac{f(x)-x}{x^2}=$ \kuo
	\fourch{不存在}{0}{$-1$}{$-2$}
\end{ti}

\begin{ti}
	$\int \sec x\tan x\dd{x}=$ \kuo
	\fourch{$\tan x+C$}{$\sec x+C$}{$\arctan x+C$}{$\arccot x+C$}
\end{ti}

\begin{ti}
	设 $f(x)$ 的一个原函数是 $\ee^{-2x}$, 则 $f(x)$ 等于\kuo
	\fourch{$-2 \ee^{-2 x}$}{$-4 \ee^{-2x}$}{$\ee^{-2x}$}{$4 \ee^{-2x}$}
\end{ti}

\begin{ti}
	$\int\cos\ee^x\dd{\ee^x}=$ \kuo
	\fourch{$\sin\ee^x+C$}{$-\sin\ee^x+C$}{$\arcsin\ee^x+C$}{$-\arcsin\ee^x+C$}
\end{ti}

\begin{ti}
	若 $\int_{0}^{1}(x+k)\dd{x}=2$, 则 $k=$ \kuo
	\fourch{0}{1}{$-1$}{$\frac{3}{2}$}
\end{ti}

\begin{ti}
	下列反常积分中收敛的是\kuo
	\fourch{$\int_{1}^{+\infty}\frac{1}{\sqrt{x}}\dd{x}$}{$\int_{1}^{+\infty}\frac{1}{1+x^2}\dd{x}$}{$\int_{1}^{+\infty}\frac{1}{1+x}\dd{x}$}{$\int_{2}^{+\infty}\frac{1}{x\ln x}\dd{x}$}
\end{ti}

\subsubsection{填空题}
\begin{ti}
	$\lim_{x\to0}\frac{\sin^2x}{x}=$ \hua
\end{ti}

\begin{ti}
	$y=x^n$ ($n$ 是正整数)的 $n$ 阶导数是 \hua
\end{ti}

\begin{ti}
	当 $a = $ \hua{} 时函数 $f(x)=a\sin x+\frac13\sin3x$ 在 $x=\frac{\uppi}{3}$ 处取极大值
\end{ti}

\begin{ti}
	$\int\sec x(\sec x-\tan x)\dd{x}=$ \hua
\end{ti}

\begin{ti}
	经过点 $(1,2)$ 且其切线的斜率为 $2x$ 的曲线方程是 \hua
\end{ti}

\begin{ti}
	通过定积分的几何意义知 $\int_{-1}^{1}(1+x)\sqrt{1-x^2}\dd{x}=$ \hua
\end{ti}

\begin{ti}
	比较大小 $\int_{1}^{3}x^2\dd{x}$ \hua{} $\int_{1}^{3}x^3\dd{x}$
\end{ti}

\begin{ti}
	由 $y=x^2+1,y=0,x=1,x=0$ 所围平面图形绕 $x$ 轴旋转一周所得立体的体积用定积分表示为 \hua
\end{ti}

\subsubsection{大题}
\begin{ti}
	求 $\lim_{x\to0}\left(1+\frac{1}{2x}\right)^{x+3}$
\end{ti}

\begin{ti}
	计算 $\lim_{x\to0}\frac{ x-\int_0^x\ee^{-t^2}\dd{t}}{x^3}$
\end{ti}

\begin{ti}
	设 $y=\sin\left(x^2\right)$, 求 $\frac{\dd{y}}{\dd{x}},\frac{\dd^2{y}}{\dd{x^2}}$
\end{ti}

\begin{ti}
	计算不定积分 $\int\sin^3x\dd{x}$
\end{ti}

\begin{ti}
	计算不定积分 $\int x\ln x\dd{x}$
\end{ti}

\begin{ti}
	计算广义积分 $\int_{2}^{+\infty}\frac{1}{x(\ln x)^2}\dd{x}$
\end{ti}

\begin{ti}
	计算 $\int_{1}^{5}\frac{\sqrt{u-1}}{u}\dd{u}$
\end{ti}

\begin{ti}
	已知
	\[
	I_n=\int_{0}^{\frac{\uppi}{2}}\sin^nt\dd{t}=
	\begin{cases}
	\frac{n-1}{n}\cdot\frac{n-3}{n-2}\cdot\frac{n-5}{n-4}\cdots\frac{2}{3}, & n\text{为大于}1\text{的奇数}\\
	\frac{n-1}{n}\cdot\frac{n-3}{n-2}\cdots\frac{1}{2}\cdot\frac{\uppi}{2}, & n\text{为正偶数}
	\end{cases},
	\]
	求星形线 $
	\begin{cases}
	x=a\cos^3t\\
	y=a\sin^3t
	\end{cases}
	$ 围成的平面图形的面积, 其中 $a>0$
\end{ti}

\begin{ti}
	设 $f(x),g(x)$ 均在 $[a,b]$ 上连续, 证明在 $(a,b)$ 内至少存在一个 $\xi$, 使得
	\[
		f(\xi)\int_{\xi}^{b}g(x)\dd{x}=g(\xi)\int_{a}^{\xi}f(x)\dd{x}
	\]
	（提示：辅助函数 $F(x) = \int_{a}^{x}f(t)\dd{t}\int_{x}^{b}g(t)\dd{t}$）
\end{ti}

\subsection{2019-2020 B12}
\subsubsection{选择题}
缺

\subsubsection{填空题}
\begin{ti}
	设 $f(x) = \ln(1 + x)$, 则 $\dd{f(x)} \bigr|_{\vartriangle x = 0.01}^{x = 2} = $ \hua
\end{ti}

\begin{ti}
	$\lim_{x \to 0} \frac{1 - \cos 2x}{x \sin x} = $ \hua
\end{ti}

\begin{ti}
	$\int_{0}^{1} x^{3} \dd{x} = $ \hua
\end{ti}

\begin{ti}
	函数 $f(x) = 4 + 8x^{3} - 3x^{4}$ 在 $x = $ \hua{} 取最大值
\end{ti}

\begin{ti}
	$\int_{-\frac{\uppi}{2}}^{\frac{\uppi}{2}} (x|\cos x| + \cos x) \dd{x} = $ \hua
\end{ti}

\begin{ti}
	用参数方程表示的曲线 $x = x(t), y = y(t)$ 在 $t \in [a,b]$ 对应的一段长度用定积分表示为 \hua
\end{ti}

\begin{ti}
	如果 $\ee^{-x}$ 是函数 $f(x)$ 的一个原函数, 则 $\int f(x) \dd{x} = $ \hua
\end{ti}

\begin{ti}
	$\int \frac{\cos 2x}{\cos x - \sin x} \dd{x} = $ \hua
\end{ti}

\subsubsection{综合题}
\begin{ti}
	计算 $\lim_{x \to \infty} \frac{\int_{0}^{x} \ee^{t^{2}} \dd{t}}{\ee^{x^{2}}}$
\end{ti}

\begin{ti}
	设 $\begin{cases}
		x = \ee^{t},\\
		y = t \ee^{t},
	\end{cases}$ 求 $\frac{\dd^{2}y}{\dd{x^{2}}}$
\end{ti}

\begin{ti}
	求 $\lim_{x \to 0} \frac{\tan x - \sin x}{\sin^{3}x}$
\end{ti}

\begin{ti}
	计算 $\int_{0}^{1} x^{2} \sqrt{1 - x^{2}} \dd{x}$
\end{ti}

\begin{ti}
	计算 $\int_{\ee}^{\ee^{2}} \frac{\dd{x}}{x \ln x}$
\end{ti}

\begin{ti}
	计算不定积分 $\int x \cos x \dd{x}$
\end{ti}

\begin{ti}
	计算 $\int \frac{\ee^{x}}{1 + \ee^{x}} \dd{x}$
\end{ti}

\begin{ti}
	求证：在 $[-1,1]$ 上 $\arcsin x + \arccos x = \frac{\uppi}{2}$
\end{ti}

\begin{ti}
	求由曲线 $y = x^{2}$ 及直线 $y = x$ 围成的平面图形的面积
\end{ti}

\subsection{2019-2020 B20}
\subsubsection{选择题}
\begin{ti}[3 分]
	函数 $f(x)$ 在点 $x_{0}$ 处有定义, 是 $f(x)$ 在该点处连续的 \kuo
	\fourch{充分条件}{必要条件}{充要条件}{无关的条件}
\end{ti}

\begin{ti}[3 分]
	若 $f(x) = \begin{cases}
		\ee^{ax}, & x < 0 \\
		b + \sin 2x, & x \geqslant 0
	\end{cases}$ 在 $x = 0$ 处可导, 则 $a,b$ 的值应为 \kuo
	\fourch{$a = -2, b = 1$}{$a = 1, b = 2$}{$a = 2, b = 1$}{$a = 2, b = -1$}
\end{ti}

\begin{ti}[3 分]
	函数 $y = \frac{1}{x}$ 在 $(0,1)$ 内的最小值是 \kuo
	\fourch{不存在}{0}{1}{任何小于 $1$ 的数}
\end{ti}

\begin{ti}[3 分]
	$\int a^{x} \dd{x} = $ \kuo
	\fourch{$\frac{1}{a} a^{x + 1} + C$}{$a^{x} + C$}{$a^{x} \ln a + C$}{$\frac{a^{x}}{\ln a} + C$}
\end{ti}

\begin{ti}[3 分]
	设 $f(x)$ 的一个原函数为 $F(x)$, 则 $\int f(2x) \dd{x} = $ \kuo
	\fourch{$2 F \bigl( \frac{x}{2} \bigr) + C$}{$F \bigl( \frac{x}{2} \bigr) + C$}{$F(2x) + C$}{$\frac{1}{2} F(2x) + C$}
\end{ti}

\begin{ti}[3 分]
	一物体由静止开始运动, 位移函数为 $s(t)$, $s(0) = 0$, 经 $t$ 秒后的速度是 $3t^{2} (\si{m/s})$, 问：在 \SI{3}{s} 后物体离开出发点的距离是多少？\kuo
	\fourch{$81$ 米}{$27$ 米}{$9$ 米}{$3$ 米}
\end{ti}

\begin{ti}[3 分]
	以下等式成立的是 \kuo
	\twoch{$\int_{-2}^{-1} \frac{1}{x} \dd{x} = \ln |-1| - \ln|-2|$}{$\int_{-1}^{1} \frac{1}{x} \dd{x} = \ln |1| - \ln|-1|$}{$\int_{-\frac{\uppi}{2}}^{\frac{\uppi}{2}} \cot x \dd{x} = \ln \bigl|\sin \frac{\uppi}{2}\bigr| - \ln \bigl|\sin\bigl( -\frac{\uppi}{2} \bigr)\bigr|$}{$\int_{-1}^{1} \bigl( x^{2} + \frac{1}{x} \bigr) \dd{x} = \int_{-1}^{1} x^{2} \dd{x}$}
\end{ti}

\begin{ti}
	设 $f(x) = \int_{0}^{x} \sin t^{2} \dd{t}, g(x) = x^{2}$, 则当 $x \to 0$ 时, $f(x)$ 是 $g(x)$ 的 \kuo
	\twoch{同阶但不等价的无穷小}{等价无穷小}{高阶无穷小}{低阶无穷小}
\end{ti}

\subsubsection{填空题}
\begin{ti}[3 分]
	设 $y = f(\sin x)$, $f$ 可微, 则 $\dd{y} = $ \hua
\end{ti}

\begin{ti}[3 分]
	$\lim_{n \to \infty} \frac{(n + 1)(3n + 2)(n + 3)}{4 n^{3}} = $ \hua
\end{ti}

\begin{ti}[3 分]
	$\int \csc x \cot x \dd{x} = $ \hua
\end{ti}

\begin{ti}[3 分]
	曲线 $y = x^{2} + \frac{1}{x - 1}$ 的铅直渐近线是 \hua
\end{ti}

\begin{ti}[3 分]
	$\int \bigl( 3 \ee^{x} + \frac{2}{x} \bigr) \dd{x} = $ \hua
\end{ti}

\begin{ti}[3 分]
	用参数方程表示的曲线 $x = x(t), y = y(t)$ 在 $t \in [a,b]$ 对应的一段长度用定积分表示为 \hua
\end{ti}

\begin{ti}[3 分]
	设函数 $f(x)$ 连续, 且 $f(x) = 2x - \int_{0}^{2} f(t) \dd{t}$, 则 $f(x) = $ \hua
\end{ti}

\begin{ti}[3 分]
	通过定积分的几何意义知 $\int_{-1}^{1} (1 + x) \sqrt{1 - x^{2}} \dd{x} = $ \hua
\end{ti}

\subsubsection{综合题}
\begin{ti}[6 分]
	计算 $\lim_{x \to 0} \frac{\int_{0}^{x} t(t + \sin t) \dd{t}}{x^{3}}$
\end{ti}

\begin{ti}[5 分]
	计算不定积分 $\int \frac{x^{2}}{x^{2} + 1} \dd{x}$
\end{ti}

\begin{ti}[6 分]
	求 $\lim_{x \to 0} \bigl( 1 + \frac{x}{2} \bigr)^{\frac{3}{x}}$
\end{ti}

\begin{ti}[6 分]
	计算 $\int_{\frac{1}{\uppi}}^{\frac{2}{\uppi}} \frac{1}{x^{2}} \cos \frac{1}{x} \dd{x}$
\end{ti}

\begin{ti}[6 分]
	$y = x^{3} + \sin x + \ln x + \ee^{x}$, 求 $y'$
\end{ti}

\begin{ti}[6 分]
	计算 $\int_{0}^{1} x \ee^{3x} \dd{x}$
\end{ti}

\begin{ti}[6 分]
	求曲线 $y = x^{2}$ 与 $x = y^{2}$ 所围成图形的面积
\end{ti}

\begin{ti}[5 分]
	计算 $\int \frac{\dd{x}}{x^{2} \sqrt{1 - x^{2}}}$
\end{ti}

\begin{ti}[6 分]
	求 $y = x^{3} + 3x^{2} - 9x + 14$ 在 $[-4,4]$ 上的最大值与最小值
\end{ti}

\subsection{2020-2021 A13}
\subsubsection{选择题（每小题3分，共24分）}
\begin{ti}
	下列等式中成立的是 \kuo
	\fourch{$\lim_{n\to\infty} \bigl( 1 + \frac{2}{n} \bigr)^n = \ee$}
	{$\lim_{n\to\infty} \bigl( 1 + \frac{1}{2n} \bigr)^n = \ee$}
	{$\lim_{n\to\infty} \bigl( 1 + \frac{1}{n} \bigr)^{2n} = \ee$}
	{$\lim_{n\to\infty} \bigl( 1 + \frac{1}{n} \bigr)^{n+2} = \ee$}
\end{ti}

\begin{ti}
	设 $y = f(\ln x)$, 则 $\dd{y} = $ \kuo
	\fourch{$f'(\ln x) \dd{x}$}{$f'(\ln x) \dd{\ln x}$}{$f'\bigl( \frac{1}{x} \bigr) \dd{x}$}{$f'(\ln x) \dd{\frac{1}{x}}$}
\end{ti}

\begin{ti}
	方程 $x^5 + x - 1 = 0$ 是 \kuo
	\fourch{有三个不同的实根}{有且仅有两个不同的实根}{没有实根}{有且仅有一个实根}
\end{ti}

\begin{ti}
	若 $\int f(x) \dd{x} = x^2 + C$, 则 $\int x f\bigl(x^2\bigr) \dd{x} = $ \kuo
	\fourch{$2x^4 + C$}{$2x^2 + C$}{$\frac{1}{2} x^4 + C$}{$\frac{1}{2} x^2 + C$}
\end{ti}

\begin{ti}
	设 $I = \int \sin x \cos x \dd{x}$, 则 $I = $ \kuo
	\fourch{$\frac{1}{2} \sin^2 x + C$}{$\frac{1}{2} \cos^2 x + C$}{$\frac{1}{4} \cos 2x + C$}{$\frac{1}{4} \sin 2x + C$}
\end{ti}

\begin{ti}
	若当 $a \leqslant x \leqslant b$, 有 $f(x) \leqslant g(x)$, 则 \kuo
	\fourch{$\int f(x) \dd{x} \leqslant \int g(x) \dd{x}$}
	{$\int_a^b f(x) \dd{x} < \int_a^b g(x) \dd{x}$}
	{$\int_a^b f(x) \dd{x} \leqslant \int_a^b g(x) \dd{x}$}
	{$\int f(x) \dd{x} < \int g(x) \dd{x}$}
\end{ti}

\begin{ti}
	以下等式成立的是 \kuo
	\fourch{$\int_{-1}^1 \cot x \dd{x} = \ln  \lvert\sin 1\rvert - \ln \lvert\sin(-1)\rvert$}
	{$\int_{-3}^{-2} \frac{1}{x} \dd{x} = \ln\lvert-2\rvert - \ln\lvert-3\rvert$}
	{$\int_{-1}^1 \bigl( x^2 + \frac{1}{x^3} \bigr) \dd{x} = \int_{-1}^1 x^2 \dd{x}$}
	{$\int_{-1}^1 \frac{1}{x^2} \dd{x} = 2 \int_0^1 \frac{1}{x^2} \dd{x}$}
\end{ti}

\begin{ti}
	星形线 $\left\lbrace\begin{array}{@{}l}
		x = a \cos^3 t \\
		y = a \sin^3 t
	\end{array}\right.$ (其中 $a>0$) 的周长为 \kuo
	\fourch
	{$4 \int_0^{\frac{\uppi}{2}} \sqrt{ \bigl[ \dd{\bigl( a \cos^3 t \bigr)} \bigr]^2 + \bigl[ \dd{\bigl( a \sin^3 t \bigr)} \bigr]^2 }$}
	{$4\int_0^{\frac{\uppi}{2}} \uppi a^2 \sin^6 t \dd{\bigl( a \cos^3 t \bigr)}$}
	{$4 \int_0^{a} \sqrt{ \bigl[ \dd{\bigl( a \cos^3 t \bigr)} \bigr]^2 + \bigl[ \dd{\bigl( a \sin^3 t \bigr)} \bigr]^2 }$}
	{$4 \int_0^{\frac{\uppi}{2}} a \sin^3t \dd{\bigl( a \cos^3 t \bigr)}$}
\end{ti}

\subsubsection{填空题（每空3分，共24分）}
\begin{ti}
	$\lim_{x \to 0} \frac{\ln(1+2x)}{\sin 5x} = $ \hua
\end{ti}

\begin{ti}
	设 $y=3^x$, 则 $y' = $ \hua
\end{ti}

\begin{ti}
	函数 $y = 2x^3 - 3x^2$ 在 $x = $ \hua{} 处有极大值
\end{ti}

\begin{ti}
	$\int \frac{\cos 2x}{\cos^2 x \sin^2 x} \dd{x} = $ \hua
\end{ti}

\begin{ti}
	函数 $f(x)$ 在 $(-\infty,+\infty)$ 上连续, 则 $\dd\bigl[ \int f(x) \dd{x} \bigr]$ 等于 \hua
\end{ti}

\begin{ti}
	由几何意义知定积分 $\int_{-1}^1 \bigl( x^4 \sin^3 x + 3 \sqrt{1-x^2} \bigr) \dd{x} = $ \hua
\end{ti}

\begin{ti}
	设函数 $f(x)$ 连续, 且 $f(x) = 3x^2 - \int_0^1 f(x) \dd{x}$, 则 $f(x) = $ \hua
\end{ti}

\begin{ti}
	由曲线 $y=x^2$ 及直线 $y=x$ 围成的平面图形绕 $x$ 轴旋转一周所得的旋转体的体积用定积分表示为 \hua
\end{ti}

\subsubsection{综合题（请写出求解或证明过程，9 小题，共 52 分）}
\begin{ti}[6分]
	求 $\lim_{x \to 0} \frac{4x^3 - 2x^2 + x}{3x^2 + 2x}$
\end{ti}

\begin{ti}[6分]
	计算 $\lim_{x \to a} \frac{1}{x-a} \int_a^x f(t) \dd{t}$, 其中 $f(x)$ 连续
\end{ti}

\begin{ti}[6分]
	设 $\left\lbrace\begin{array}{@{}l}
		x = 1+t^2 \\
		y = 3t^3+4
	\end{array}\right.$, 求 $\frac{\dd[2]{y}}{\dd{x^2}}$
\end{ti}

\begin{ti}[5分]
	求不定积分 $\int \frac{1}{x} \ln x \dd{x}$
\end{ti}

\begin{ti}[5分]
	计算不定积分 $\int x \arctan x \dd{x}$
\end{ti}

\begin{ti}[6分]
	计算 $\int_0^{+\infty} \frac{\dd{x}}{\sqrt{x} + x\sqrt{x}}$
\end{ti}

\begin{ti}[6分]
	$\int_{-\frac{\uppi}{4}}^{\frac{\uppi}{2}} \sqrt{\cos x - \cos^3 x} \dd{x}$
\end{ti}

\begin{ti}[6分]
	求由曲线 $y=\ee^x$ 与直线 $y=\ee x$、$y$ 轴围成的图形的面积
\end{ti}

\begin{ti}[6分]
	证明恒等式：$\arctan x + \arccot x = \frac{\uppi}{2}$ $(-\infty<x<\infty)$
\end{ti}

\section{《高等数学（二）》期中}
\subsection{2017-2018}
\subsubsection{选择题}
\begin{ti}
	下列方程中, 设 $y_{1},y_{2}$ 是它的解, 可以推知 $y_{1} + y_{2}$ 也是它的解的方程是 \kuo
	\twoch{$y' + p(x)y + q(x) = 0$}{$y'' + p(x)y' + q(x)y = 0$}{$y'' + p(x)y' + q(x)y = f(x)$}{$y'' + p(x)y' + q(x) = 0$}
\end{ti}

\begin{ti}
	微分方程 $\frac{\dd{y}}{\dd{x}} = \frac{y}{x} \ln \frac{y}{x}$ 是 \kuo
	\twoch{可分离变量的微分方程}{齐次微分方程}{一阶线性非齐次微分方程}{二阶微分方程}
\end{ti}

\begin{ti}
	二阶齐次线性微分方程 $y'' + y = 0$ 的通解为 $y = $ \kuo
	\twoch{$C_{1} \cos x$}{$C_{2} \sin x$}{$(C_{1} + C_{2})(\cos x + \sin x)$}{$C_{1} \cos x + C_{2} \sin x$}
\end{ti}

\begin{ti}
	两向量 $\vec{a} \cdot \vec{b}$ 平行的充要条件是 \kuo
	\fourch{$\vec{a} \cdot \vec{b} = \vec{0}$}{$\vec{a} \cdot \vec{b} = 0$}{$\vec{a} \times \vec{b} = \vec{0}$}{$\vec{a} \times \vec{b} = 0$}
\end{ti}

\begin{ti}
	设 $\vec{a} = (2,4,-1), \vec{b} = (0,-2,2)$, 则同时与 $\vec{a},\vec{b}$ 垂直的单位向量 $\vec{n} = $ \kuo
	\twoch{$6 \vec{i} - 4 \vec{j} - 4 \vec{k}$}{$-6 \vec{i} + 4 \vec{j} + 4 \vec{k}$}{$\frac{1}{| \vec{a} + \vec{b} |} \bigl( 6 \vec{i} - 4 \vec{j} - 4 \vec{k} \bigr)$}{$\frac{\pm 1}{| \vec{a} \times \vec{b} |} \bigl( 6 \vec{i} - 4 \vec{j} - 4 \vec{k} \bigr)$}
\end{ti}

\begin{ti}
	若平面 $Ax + By + Cz + D = 0$ 过 $x$ 轴, 则 \kuo
	\fourch{$A = D = 0$}{$B = 0, C \ne 0$}{$B \ne 0, C = 0$}{$B = C = 0$}
\end{ti}

\begin{ti}
	设 $f(x,y) = x^{2} + (y - 3) \arctan \frac{x}{y}$, 则 $f_{x}(2,3) = $ \kuo
	\fourch{1}{2}{3}{4}
\end{ti}

\begin{ti}
	设 $z = x^{2}y + 3xy^{2} + x$, 则 $\frac{\partial^{2}z}{\partial x \partial y} = $ \kuo
	\fourch{$x + 3y$}{$2x + 3y$}{$2x + 6y$}{$2x + 6y + 1$}
\end{ti}

\begin{ti}
	若 $z = \ee^{-x + y}$, 则 $\dd{z}\bigr|_{(1,1)} = $ \kuo
	\fourch{$- \dd{x} + \dd{y}$}{$\dd{x} - \dd{y}$}{$\dd{x} + \dd{y}$}{$- \dd{x} - \dd{y}$}
\end{ti}

\subsubsection{填空题}
\begin{ti}
	微分方程 $(y'')^{3} + y' \sin x = 0$ 的阶数为 \hua
\end{ti}

\begin{ti}
	方程 $y \dd{x} = x \dd{y}$ 的通解为 $y = $ \hua
\end{ti}

\begin{ti}
	微分方程 $y^{(4)} \sin x - y' = \ee^{x}$ 的通解中所含相互独立的任意常数的个数为 \hua
\end{ti}

\begin{ti}
	设向量 $\vec{a}$ 的方向角为 $\alpha,\beta,\gamma$, 满足 $\cos \alpha = 1$ 时, 向量 $\vec{\alpha}$ 垂直于 \hua{} 坐标面
\end{ti}

\begin{ti}
	过点 $(1,0,1)$ 及以 $(2,2,4)$ 为方向向量的直线的参数方程是 \hua
\end{ti}

\begin{ti}
	圆锥面 $z = \sqrt{x^{2} + y^{2}}$ 与平面 $z = 2$ 所围立体在 $xoy$ 平面上投影曲线方程是 \hua
\end{ti}

\begin{ti}
	函数 $z = \ln(x^{2} + y^{2} - 4)$ 的定义域 $D$ 是 \hua
\end{ti}

\begin{ti}
	设 $z = 3 x^{2}y^{2} + \ee^{x^{2}y}$, 则 $\frac{\partial z}{\partial x} = $ \hua
\end{ti}

\begin{ti}
	设 $z = \ln(2x + y)$, 则 $\frac{\partial^{2}z}{\partial x \partial y} = $ \hua
\end{ti}

\subsubsection{综合题}
\begin{ti}
	求方程 $\frac{\dd{y}}{\dd{x}} + \frac{1 - 2x}{x^{2}} y - 1 = 0$ 的通解
\end{ti}

\begin{ti}
	求方程 $y'' - 2y' - 3y = 3x + 1$ 的通解
\end{ti}

\begin{ti}
	设 $z = \ee^{u + v}$, 而 $u = xy,v = x - y$, 求 $\frac{\partial z}{\partial x},\frac{\partial z}{\partial y}$
\end{ti}

\begin{ti}
	求曲面 $2x^{2} + 3y^{2} + z^{2} = 6$ 在点 $(1,1,1)$ 处的切平面和法线方程
\end{ti}

\begin{ti}
	求函数 $f(x,y) = x + y - x^{2} - y^{2}$ 在曲线 $x^{2} + y^{2} = 1$ 上的最大值和最小值
\end{ti}

\subsection{2018-2019 B10}
\subsubsection{选择题}
\begin{ti}
	微分方程 $(y')^3+3\sqrt{y''}+x^4y'''=\sin x$ 的阶数是\kuo
	\fourch{1}{4}{2}{3}
\end{ti}

\begin{ti}
	设 $f(x,y)=x-y-\sqrt{x^2+y^2}$ , 则 $f_{x}(3,4)=$\kuo
	\fourch{$\frac{3}{5}$}{$\frac{2}{5}$}{$-\frac{2}{5}$}{$\frac{1}{5}$}
\end{ti}

\begin{ti}
	微分方程 $y'=\frac{y}{x}$ 的一个特解是\kuo
	\fourch{$y=2x$}{$\ee^y=x$}{$y=x^2$}{$y=\ln x$}
\end{ti}

\begin{ti}
	若 $z=\ln\sqrt{1+x^2+y^2}$ , 则 $\left.\dd{z}\right|_{(1,1)}=$\kuo
	\fourch{$\frac{\dd{x}+\dd{y}}{3}$}{$\frac{\dd{x}+\dd{y}}{2}$}{$\frac{\dd{x}+\dd{y}}{1}$}{$3(\dd{x}+\dd{y})$}
\end{ti}

\begin{ti}
	设直线 $L:\begin{cases}
	x+3y+2z+1=0\\
	2x-y-10z+3=0
	\end{cases}$ , 平面 $\eta:4x-2y+z-2=0$ , 则\kuo
	\fourch{$L$ 在 $\eta$ 上}{$L$ 平行于 $\eta$}{$L$ 垂直于 $\eta$}{$L$ 与 $\eta$ 斜交}
\end{ti}

\begin{ti}
	方程 $y'+3xy=6x^2y$ 是\kuo
	\twoch{二阶微分方程}{非线性微分方程}{一阶线性非齐次微分方程}{可分离变量的微分方程}
\end{ti}

\begin{ti}
	曲面 $\frac{x^2}{9}-\frac{y^2}{4}+\frac{z^2}{4}=1$ 与平面 $x=y$ 的交线是\kuo
	\fourch{两条直线}{双曲线}{椭圆}{抛物线}
\end{ti}

\begin{ti}
	设 $z=\ee^{x^2y}$ , 则 $\frac{\partial^2z}{\partial x\partial y}=$\kuo
	\twoch{$2y\left(1+x^3\right)\ee^{x^2y}$}{$\ee^{x^2y}$}{$2x\left(1+x^2y\right)\ee^{x^2y}$}{$2x\ee^{x^2y}$}
\end{ti}

\begin{ti}
	下列结论正确的是\kuo
	\twoch{$\vec{a}\times\left(\vec{b}-\vec{c}\right)=\vec{a}\times\vec{b}-\vec{a}\times\vec{c}$}{若 $\vec{a}\times\vec{b}=\vec{a}\times\vec{c}$ 且 $\vec{a}\ne\vec{0}$ , 则 $\vec{b}=\vec{c}$}{$\vec{a}\times\vec{b}=\vec{b}\times\vec{a}$}{若 $\left|\vec{a}\right|=1,\left|\vec{b}\right|=1$ , 则 $\left|\vec{a}\times\vec{b}\right|=1$}
\end{ti}


\subsubsection{填空题}
\begin{ti}
	平面过点 $(2,0,0),(0,1,0),(0,0,0.5)$ , 则该平面的方程是 \hua{}
\end{ti}

\begin{ti}
	设 $y_1$ 是 $y''+p(x)y'+q(x)y=f(x)$ 的解, $y_2$ 是 $y''+p(x)y'+q(x)y=f(x)$ 的解, 则 $y_1+y_2$ 是 \hua{}方程的解
\end{ti}

\begin{ti}
	设 $z=y\arctan x$ , 则 $\left.\mathrm{grad}\,z\right|_{(1,2)}=$ \hua{}
\end{ti}

\begin{ti}
	过点 $P(0,2,4)$ 且与两平面 $x+2z=1$ 和 $y-2z=2$ 平行的直线方程是 \hua{}
\end{ti}

\begin{ti}
	设 $f(x,y)=\arcsin\frac{y}{x}$ , 则 $f_y(1,0)=$ \hua{}
\end{ti}

\begin{ti}
	$y=\ee^x$ 是微分方程 $y''+py'+6y=0$ 的一个特解, 则 $p=$ \hua{}
\end{ti}

\begin{ti}
	已知平面 $\eta_1:A_1x+B_1y+C_1z+D_1=0$ 与平面 $\eta_2:A_2x+B_2y+C_2z+D_2=0$, 则 $\eta_1\perp\eta_2$ 的充要条件是 \hua{}
\end{ti}

\begin{ti}
	微分方程 $y''+2y'+5y=0$ 的通解为 $y=$ \hua{}
\end{ti}

\begin{ti}
	设 $z=\ee^{xy}+\cos\left(x^2+y\right)$, 则 $\frac{\partial z}{\partial y}=$ \hua{}
\end{ti}

\subsubsection{大题}
\begin{ti}
	求方程 $\frac{\dd{z}}{\dd{x}}=-z+4x$ 的通解
\end{ti}

\begin{ti}
	求曲线 $2z+1=\ln(xy)+\ee^z$ 在点 $M_{0}(1,1,0)$ 处的切平面和法线方程
\end{ti}

\begin{ti}
	设由方程组 $\begin{cases}
	x+y+z=0\\
	x^2+y^2+z^2=1
	\end{cases}$
	确定了隐函数 $x=x(z),y=y(z)$ , 求 $\frac{\dd{x}}{\dd{z}},\frac{\dd{y}}{\dd{z}}$
\end{ti}

\begin{ti}
	求方程 $y''+6y'+13y=\ee^t$ 的通解
\end{ti}

\begin{ti}
	设 $z=x^2y+\sin x+\varphi(xy+1)$ , 且 $\varphi(u)$ 具有一阶连续导数, 求 $\frac{\partial z}{\partial x},\frac{\partial z}{\partial y}$
\end{ti}


\subsection{2018-2019 B10 答案}
\subsubsection{选择题}
\begin{ti}
	微分方程 $(y')^3+3\sqrt{y''}+x^4y'''=\sin x$ 的阶数是\kuoD
	\fourch{1}{4}{2}{3}
\end{ti}

\begin{ti}
	设 $f(x,y)=x-y-\sqrt{x^2+y^2}$ , 则 $f_{x}(3,4)=$\kuoB
	\fourch{$\frac{3}{5}$}{$\frac{2}{5}$}{$-\frac{2}{5}$}{$\frac{1}{5}$}
\end{ti}

\begin{ti}
	微分方程 $y'=\frac{y}{x}$ 的一个特解是\kuoA
	\fourch{$y=2x$}{$\ee^y=x$}{$y=x^2$}{$y=\ln x$}
\end{ti}

\begin{ti}
	若 $z=\ln\sqrt{1+x^2+y^2}$ , 则 $\left.\dd{z}\right|_{(1,1)}=$\kuoA
	\fourch{$\frac{\dd{x}+\dd{y}}{3}$}{$\frac{\dd{x}+\dd{y}}{2}$}{$\frac{\dd{x}+\dd{y}}{1}$}{$3(\dd{x}+\dd{y})$}
\end{ti}

\begin{ti}
	设直线 $L:\begin{cases}
	x+3y+2z+1=0\\
	2x-y-10z+3=0
	\end{cases}$ , 平面 $\eta:4x-2y+z-2=0$ , 则\kuoC
	\fourch{$L$ 在 $\eta$ 上}{$L$ 平行于 $\eta$}{$L$ 垂直于 $\eta$}{$L$ 与 $\eta$ 斜交}
\end{ti}

\begin{ti}
	方程 $y'+3xy=6x^2y$ 是\kuoD
	\twoch{二阶微分方程}{非线性微分方程}{一阶线性非齐次微分方程}{可分离变量的微分方程}
\end{ti}

\begin{ti}
	曲面 $\frac{x^2}{9}-\frac{y^2}{4}+\frac{z^2}{4}=1$ 与平面 $x=y$ 的交线是\kuoB
	\fourch{两条直线}{双曲线}{椭圆}{抛物线}
\end{ti}

\begin{ti}
	设 $z=\ee^{x^2y}$ , 则 $\frac{\partial^2z}{\partial x\partial y}=$\kuoC
	\twoch{$2y\left(1+x^3\right)\ee^{x^2y}$}{$\ee^{x^2y}$}{$2x\left(1+x^2y\right)\ee^{x^2y}$}{$2x\ee^{x^2y}$}
\end{ti}

\begin{ti}
	下列结论正确的是\kuoA
	\twoch{$\vec{a}\times\left(\vec{b}-\vec{c}\right)=\vec{a}\times\vec{b}-\vec{a}\times\vec{c}$}{若 $\vec{a}\times\vec{b}=\vec{a}\times\vec{c}$ 且 $\vec{a}\ne\vec{0}$ , 则 $\vec{b}=\vec{c}$}{$\vec{a}\times\vec{b}=\vec{b}\times\vec{a}$}{若 $\left|\vec{a}\right|=1,\left|\vec{b}\right|=1$ , 则 $\left|\vec{a}\times\vec{b}\right|=1$}
\end{ti}


\subsubsection{填空题}
\begin{ti}
	平面过点 $(2,0,0),(0,1,0),(0,0,0.5)$ , 则该平面的方程是 \huaa{$\frac{x}{2}+y+2z=1$}
\end{ti}

\begin{ti}
	设 $y_1$ 是 $y''+p(x)y'+q(x)y=f(x)$ 的解, $y_2$ 是 $y''+p(x)y'+q(x)y=f(x)$ 的解, 则 $y_1+y_2$ 是 \huaa{$y''+p(x)y'+$} \huaa{$q(x)y=2f(x)$} 方程的解
\end{ti}

\begin{ti}
	设 $z=y\arctan x$ , 则 $\left.\mathrm{grad}\,z\right|_{(1,2)}=$ \huaa{$\dd{x}+\frac{\uppi}{4}\dd{y}$}
\end{ti}

\begin{ti}
	过点 $P(0,2,4)$ 且与两平面 $x+2z=1$ 和 $y-2z=2$ 平行的直线方程是 \huaa{$\frac{x}{-2}=\frac{y-2}{2}=\frac{z-4}{1}$}
\end{ti}

\begin{ti}
	设 $f(x,y)=\arcsin\frac{y}{x}$ , 则 $f_y(1,0)=$ \huaa{$1$}
\end{ti}

\begin{ti}
	$y=\ee^x$ 是微分方程 $y''+py'+6y=0$ 的一个特解, 则 $p=$ \huaa{$-7$}
\end{ti}

\begin{ti}
	已知平面 $\eta_1:A_1x+B_1y+C_1z+D_1=0$ 与平面 $\eta_2:A_2x+B_2y+C_2z+D_2=0$, 则 $\eta_1\perp\eta_2$ 的充要条件是 \huaa{$A_1A_2+B_1B_2+C_1C_2=0$}
\end{ti}

\begin{ti}
	微分方程 $y''+2y'+5y=0$ 的通解为 $y=$ \huaa{$C_1\ee^{-x}\sin(2x)+C_2\ee^{-x}\cos(2x)$}
\end{ti}

\begin{ti}
	设 $z=\ee^{xy}+\cos\left(x^2+y\right)$, 则 $\frac{\partial z}{\partial y}=$ \huaa{$x\ee^{xy}-\sin\left(x^2+y\right)$}
\end{ti}

\subsubsection{大题}
\begin{ti}
	求方程 $\frac{\dd{z}}{\dd{x}}=-z+4x$ 的通解
	\begin{solution}
		运用一阶线性非齐次微分方程公式, 得
		\begin{align*}
			z&=\ee^{-\int\dd{x}}\left( \int 4x\ee^{\int \dd{x}}\dd{x}+C\right) =\ee^{-x}\left( \int 4x\ee^{x}\dd{x}+C\right) \\
			&=\ee^{-x}\left( 4(x-1)\ee^{x}+C\right) =4(x-1)+C\ee^{-x}
		\end{align*}
	\end{solution}
\end{ti}

\begin{ti}
	求曲线 $2z+1=\ln(xy)+\ee^z$ 在点 $M_{0}(1,1,0)$ 处的切平面和法线方程
\end{ti}

\begin{ti}
	设由方程组 $\begin{cases}
		x+y+z=0\\
		x^2+y^2+z^2=1
		\end{cases}$
	确定了隐函数 $x=x(z),y=y(z)$ , 求 $\frac{\dd{x}}{\dd{z}},\frac{\dd{y}}{\dd{z}}$
	\begin{solution}
		对方程组 $\begin{cases}
		x+y+z=0\\
		x^2+y^2+z^2=1
		\end{cases}$ 两式求微分, 得
		\begin{equation*}
			\begin{cases}
			\dd{x}+\dd{y}+\dd{z}=0\\
			2x\dd{x}+2y\dd{y}+2z\dd{z}=0
			\end{cases}
		\end{equation*}
		解得
		\begin{equation*}
			\begin{cases}
			\frac{\dd{x}}{\dd{z}}=-\frac{x+2z}{2x+z}\\
			\frac{\dd{y}}{\dd{z}}=-\frac{y+2x}{2y+z}
			\end{cases}
		\end{equation*}
	\end{solution}
\end{ti}

\begin{ti}
	求方程 $y''+6y'+13y=\ee^t$ 的通解
	\begin{solution}
		方程 $y''+6y'+13y=\ee^t$ 对应的齐次方程 $y''+6y'+13y=0$ 的特征方程为 $r^2+6r+13=0$ , 解得 $r=-3\pm2\ii$, 那么齐次方程的通解为 $C_1\ee^{-3t}\sin(2t)+C_2\ee^{-3t}\cos(2t)$
		
		设特解为 $a\ee^{t}$, 代入方程 $y''+6y'+13y=\ee^t$ 后解得 $a=\frac{1}{20}$
		
		综上, 方程 $y''+6y'+13y=\ee^t$ 的通解为 $C_1\ee^{-3t}\sin(2t)+C_2\ee^{-3t}\cos(2t)+\frac{\ee^x}{20}$
	\end{solution}
\end{ti}

\begin{ti}
	设 $z=x^2y+\sin x+\varphi(xy+1)$ , 且 $\varphi(u)$ 具有一阶连续导数, 求 $\frac{\partial z}{\partial x},\frac{\partial z}{\partial y}$
	\begin{solution}
		$\frac{\partial z}{\partial x}=2xy+\cos x+y\varphi'(xy+1)$ , $\frac{\partial z}{\partial y}=x^2+x\varphi'(xy+1)$
	\end{solution}
\end{ti}

\subsection{2019-2020}
线上考试
\subsubsection{选择题}
\begin{ti}
	$y = \ee^{2x}$ 是微分方程 $y'' + py' + 6y = 0$ 的一个特解, 则此方程的通解是 \kuo
	\twoch{$y=C_1 \ee^{2x} + C_2 \ee^{-3x}$}{$y = (C_1 + xC_2) \ee^{2x}$}{$y = C_1 \ee^{2x} + C_2 \ee^{3x}$}{$y = \ee^{2x} (C_1 \sin 3x + C_2 \cos 3x)$}
\end{ti}

\begin{ti}
	非齐次线性微分方程 $x'' - 2x' + x = (t+2) \ee^{2t}$ 的一个待定的解形式 $x^\astt = $ \kuo
	\fourch{$t(At+B)\ee^{2t}$}{$(At+B)\ee^{2t}$}{$At^2 \ee^{2t}$}{$t^2 (At+B)\ee^{2t}$}
\end{ti}

\begin{ti}
	方程 $y' + 3xy = 6x^2$ 是 \kuo
	\twoch{可分离变量的微分方程}{齐次微分方程}{一阶线性非齐次微分方程}{二阶微分方程}
\end{ti}

\begin{ti}
	微分方程 $\bigl(y'\bigr)^4 + 3 \sqrt{y''} + y^2 = \sin x$ 的阶数是 \kuo
	\fourch{1}{2}{3}{4}
\end{ti}

\begin{ti}
	微分方程 $y' = \frac{y}{x}$ 的通解是 \kuo
	\fourch{$y = C\ee^{x}$}{$\ee^{y} = x$}{$\ee^{y} = Cx$}{$y = Cx$}
\end{ti}

\begin{ti}
	方程 $y'' - 2y' - 3y = 0$ 的通解是 \kuo
	\twoch{$y = C_1 \ee^{-x} + C_2 \ee^{3x}$}{$y = \frac{C_1}{x} + C_2 x^2$}{$y = C_1 \ee^{x} + C_2 \ee^{-x}$}{$y = C_1 x + \frac{C_2}{x^3}$}
\end{ti}

\begin{ti}
	已知向量 $\vec a = (0,3,4), \vec b = (2,1,-2)$, 则 $\Prj_{\vec b} \vec a = $ \kuo
	\fourch{5}{$-\frac{1}{3}$}{$-\frac{5}{3}$}{$\frac{1}{3}$}
\end{ti}

\begin{ti}
	平面 $y + z = 1$ 的位置是 \kuo
	\fourch{过 $x$ 轴}{垂直于 $x$ 轴}{平行于 $x$ 轴}{平行于 $yOz$ 面}
\end{ti}

\begin{ti}
	设 $\vec a = (a_x,a_y,a_z), \vec b = (b_x,b_y,b_z)$ 均为非零向量, 且 $a_x b_x + a_y b_y + a_z b_z = 0$, 则 \kuo
	\fourch{$\vec a \parallel \vec b$}{$\vec a + \vec b = \vec 0$}{$\vec a = \lambda \vec b (\lambda \ne 0)$}{$\vec a \perp \vec b$}
\end{ti}

\begin{ti}
	两向量 $\vec a, \vec b$ 平行的充要条件是 \kuo
	\fourch{$\vec a \cdot \vec b = \vec 0$}{$\vec a \cdot \vec b = 0$}{$\vec a \times \vec b = \vec 0$}{$\vec a \times \vec b = 0$}
\end{ti}

\begin{ti}
	设 $\vec a = (2,4,-1), \vec b = (0,-2,2)$, 则同时与 $\vec a$、$\vec b$ 垂直的单位向量 $\vec n = $ \kuo
	\twoch{$6 \vec i - 4 \vec j - 4 \vec k$}
	{$- 6 \vec i + 4 \vec j + 4 \vec k$}
	{$\frac{1}{\bigl| \vec a + \vec b \bigr|} \bigl( 6 \vec i - 4 \vec j - 4 \vec k \bigr)$}
	{$\frac{\pm 1}{\bigl| \vec a \times \vec b \bigr|} \bigl( 6 \vec i - 4 \vec j - 4 \vec k \bigr)$}
\end{ti}

\begin{ti}
	下列平面中通过坐标原点的平面是 \kuo
	\twoch{$x=1$}{$x + 2z + 3y + 4 = 0$}{$3(x-1) - y + (z+3) = 0$}{$x + y + z = 1$}
\end{ti}

\begin{ti}
	两向量 $\vec a, \vec b$ 平行的充要条件是 \kuo
	\fourch{$\vec a \cdot \vec b = \vec 0$}{$\vec a \cdot \vec b = 0$}{$\vec a \times \vec b = \vec 0$}{$\vec a \times \vec b = 0$}
\end{ti}

\begin{ti}
	曲线 $\begin{cases}
		x = t, \\
		y = 2t^2, \\
		z = 3t^3,
	\end{cases}$ 在点 $(1,2,3)$ 处的一个切向量为 \kuo
	\fourch{$(1,2,3)$}{$(2,4,6)$}{$(1,4,9)$}{$(1,4,8)$}
\end{ti}

\begin{ti}
	设 $f(x,y) = x^2 + (y-3) \arctan \frac{x}{y}$, 则 $f_x(2,3) = $ \kuo
	\fourch{1}{2}{3}{4}
\end{ti}

\begin{ti}
	函数 $z = x^2 + y^2$ 在点 $(0,0)$ 处 \kuo
	\twoch{连续, 偏导数存在}{连续, 偏导数不存在}{不连续, 但偏导数存在}{不连续, 且偏导数不存在}
\end{ti}

\begin{ti}
	若 $z = \arctan \frac{x}{y}$, 则 $\dd{z} = $ \kuo
	\fourch{$\frac{x \dd{y} - y \dd{x}}{x^2+y^2}$}{$\frac{x \dd{x} - y \dd{y}}{x^2+y^2}$}{$\frac{y \dd{x} + x \dd{y}}{x^2+y^2}$}{$\frac{y \dd{x} - x \dd{y}}{x^2+y^2}$}
\end{ti}

\begin{ti}
	设 $z = x^2 y + 3x y^2 + x$, 则 $\frac{\partial^2 z}{\partial x \partial y} = $ \kuo
	\fourch{$x+3y$}{$2x+3y$}{$2x+6y$}{$2x+6y+1$}
\end{ti}

\begin{ti}
	设 $z = x^y$, 则 $\frac{\partial^2 z}{\partial y^2} = $ \kuo
	\fourch{$x^y$}{$\frac{x^y}{\ln^2x}$}{$y (y-1) x^{y-2}$}{$x^y \ln^2x$}
\end{ti}

\begin{ti}
	$z = 3x^2y - xy^3 + 8$ 在 $A(1,2)$ 处沿 $A$ 到 $B(3,0)$ 方向的方向导数 $\frac{\partial z}{\partial l} = $ \kuo
	\fourch{$\frac{5\sqrt{2}}{2}$}{$-\frac{5\sqrt{2}}{2}$}{$-\frac{13\sqrt{2}}{2}$}{$\frac{13\sqrt{2}}{2}$}
\end{ti}

\subsubsection{计算题}
\begin{ti}
	求方程 $y'' - 4y' + 4y = 2xe^{2x}$ 的通解
\end{ti}

\begin{ti}
	求圆锥面 $z = \smqty({x^2+y^2})$ 与平面 $z = 2$ 所围立体在 $xOy$ 平面上的投影区域
\end{ti}

\begin{ti}
	求函数 $f(x,y) = x^3 - y^3 - 3xy + 1$ 的极值
\end{ti}

\begin{ti}
	求曲面 $z = y + \ln x - \ln z$ 在点 $M_0(1,1,1)$ 处的切平面和法线方程
\end{ti}

\begin{ti}
	求方程 $\frac{\dd{y}}{\dd{x}} = \frac{y - 3x}{x}$ 的通解
\end{ti}

\begin{ti}
	设 $z = \ee^{u+v}$, 而 $u = xy$, $v = x - y$, 求 $\frac{\partial z}{\partial x}, \frac{\partial z}{\partial y}$
\end{ti}

\subsection{难度与考试近似的题}
\subsubsection{选择题}
\begin{ti}
	下列平面中通过坐标原点的平面是 \kuo
	\twoch{$3(x - 1) - y + (z + 3) = 0$}{$x + y + z = 1$}{$x = 1$}{$x + 2z + 3y + 4 = 0$}
\end{ti}

\begin{ti}
	直线 $L: \frac{x + 3}{-2} = \frac{y + 4}{-7} = \frac{z}{3}$ 与平面 $\pi: 4x - 2y - 2z = 3$ 的关系是 \kuo
	\fourch{$L$ 在 $\pi$ 上}{相交但不垂直}{平行}{垂直相交}
\end{ti}

\begin{ti}
	设 $f(x,y) = xy + (2y - 1) \arccos \frac{x}{y}$, 则 $f_{x}(1,2) = $ \kuo
	\fourch{$2 + 2\sqrt{3}$}{$2 + \sqrt{3}$}{$2 - 2\sqrt{3}$}{$2 - \sqrt{3}$}
\end{ti}

\begin{ti}
	设 $z = y^{x}$, 则 $\frac{\partial^{2}z}{\partial x \partial y} = $ \kuo
	\fourch{$y^{x-1} (1 + x \ln y)$}{$y^{x} \ln^{2}x$}{$x y^{x - 1} \ln x$}{$y^{x - 1} (x + \ln y)$}
\end{ti}

\begin{ti}
	$z = 3x^{2}y - xy^{3} + 8$ 在 $A(1,2)$ 处沿 $A$ 到 $B(3,0)$ 方向的方向导数 $\frac{\partial z}{\partial l} = $ \kuo
	\fourch{$-\frac{13 \sqrt{2}}{2}$}{$\frac{13 \sqrt{2}}{2}$}{$\frac{5 \sqrt{2}}{2}$}{$- \frac{5 \sqrt{2}}{2}$}
\end{ti}

\begin{ti}
	微分方程 $(y'')^{3} + 3y''' + y^{4} = x$ 的阶数是 \kuo
	\fourch{1}{2}{3}{4}
\end{ti}

\begin{ti}
	微分方程 $y' = y$ 的通解是 \kuo
	\fourch{$y = c \ee^{-\frac{x}{2}}$}{$y = c \ee^{-x}$}{$cy = \ee^{x^{2}}$}{$y = c\ee^{x}$}
\end{ti}

\begin{ti}
	方程 $y' + 3xy = 6x^{2}$ 是 \kuo
	\twoch{一阶线性非齐次微分方程}{二阶微分方程}{可分离变量的微分方程}{齐次微分方程}
\end{ti}

\begin{ti}
	设 $\vec{a} = (a_{x},a_{y},a_{z}), \vec{b} = (b_{x},b_{y},b_{z})$ 均为非零向量, 且 $a_{x}b_{x} + a_{y}b_{y} + a_{z}b_{z} = 0$, 则 \kuo
	\fourch{$\vec{a} = \lambda \vec{b} (\lambda \ne 0)$}{$\vec{a} \perp \vec{b}$}{$\vec{a} \parallel \vec{b}$}{$\vec{a} + \vec{b} = \vec{0}$}
\end{ti}

\subsubsection{填空题}
\begin{ti}
	设 $z = \cos(x - y)$, 则 $\frac{\partial z}{\partial y} = $ \hua
\end{ti}

\begin{ti}
	设 $f(x,y) = \arctan \frac{y}{x}$, 则 $f_{x}(1,1) = $ \hua
\end{ti}

\begin{ti}
	设 $z = 2x^{2}y + \sin(xy)$, 则 $\frac{\partial^{2}z}{\partial x^{2}} = $ \hua
\end{ti}

\begin{ti}
	微分方程 $y'' \sin x - y' = \ln x$ 的通解中所含相互独立的任意常数的个数为 \hua
\end{ti}

\begin{ti}
	方程 $\frac{\dd{y}}{\dd{x}} = \frac{\bigl( \frac{y}{x} \bigr)^{2}}{\frac{y}{x} - 1}$ 的通解是 $y = $ \hua
\end{ti}

\begin{ti}
	微分方程 $y'' - 5y' + 4y = 0$ 的通解为 $y = $ \hua
\end{ti}

\begin{ti}
	设 $|\vec{a}| = 3, \bigl| \vec{b} \bigr| = 4$, 且 $\vec{a} \perp \vec{b}$, 则 $\bigl| \bigl( \vec{a} + \vec{b} \bigr) \times \bigl( \vec{a} - \vec{b} \bigr) \bigr| = $ \hua
\end{ti}

\begin{ti}
	直线 $\begin{cases}
		3x - 2y + z + 1 = 0 \\
		2x + y - z - 2 = 0
	\end{cases}$ 的方向向量是 \hua
\end{ti}

\begin{ti}
	曲线 $\begin{cases}
		y^{2} + z^{2} - 2x = 0 \\
		z = 2
	\end{cases}$ 在 $xoy$ 面上的投影曲线方程是 \hua
\end{ti}

\subsubsection{综合题}
\begin{ti}[6 分]
	求方程 $\frac{\dd{y}}{\dd{x}} = y + \cos x$ 的通解
\end{ti}

\begin{ti}[10 分]
	求方程 $s'' - 4s = t + 1$ 的通解
\end{ti}

\begin{ti}[10 分]
	设函数 $z = z(x,y)$ 由方程 $\cos x + 3y - z = \ee^{z}$ 所确定, 求 $\dd{z}$
\end{ti}

\begin{ti}[10 分]
	求曲面 $\ee^{z} - z + x^{2}y = 3$ 在点 $(1,2,0)$ 处的切平面及法线方程
\end{ti}

\begin{ti}[10 分]
	某工厂生产两种产品甲和乙, 出售单价分别为 $10$ 元与 $9$ 元, 生产 $x$ 单位的产品甲与生产 $y$ 单位的产品乙的总费用是 $300 + 3x + 2y + 0.01 (3x^{2} + xy + 3y^{2})$ 元, 求取得最大利润时, 两种产品的产量各为多少单位？
\end{ti}


\section{《高等数学（二）》期末}

\subsection{2017-2018 A}
\subsubsection{选择题（每小题 $3$ 分，共 $24$ 分）}
\begin{ti}
	微分方程 $y''-6 y'+9 y=\left(6 x^{2}+2\right) \ee^{x}$ 的待定特解的一个形式可为 \kuo
	\twoch{$y^{*}=\left(a x^{2}+b x+c\right) \ee^{x}$}{$y^{*}=x\left(a x^{2}+b x+c\right) \ee^{x}$}{$y^{*}=x^{2}\left(a x^{2}+b x+c\right) \ee^{x}$}{$y^{*}=x^{2}\left(x^{2}+1\right) \ee^{x}$}
\end{ti}

\begin{ti}
	设向量 $\vec{a}$ 的三个方向角为 $\alpha$、$\beta$、$\gamma$ , 且已知 $\alpha=60^{\circ}$、$\beta=120^{\circ}$ , 则 $\gamma=$ \kuo
	\fourch{$120^{\circ}$}{$60^{\circ}$}{$45^{\circ}$}{$30^{\circ}$}
\end{ti}

\begin{ti}
	设 $z=\arctan \ee^{x y}$ , 则 $\frac{\partial z}{\partial y}=$ \kuo
	\fourch{$-\frac{x \ee^{x y}}{\sqrt{1-\ee^{2 x y}}}$}{$\frac{x \ee^{x y}}{\sqrt{1-\ee^{2 x y}}}$}{$-\frac{x \ee^{x y}}{1+\ee^{2 x y}}$}{$\frac{x \ee^{x y}}{1+\ee^{2 x y}}$}
\end{ti}

\begin{ti}
	$D$ 为平面区域 $x^{2}+y^{2} \leqslant 4$ , 利用二重积分的性质, $\iint_{D}\left(x^{2}+4 y^{2}+9\right) \dd{x} \dd{y}$ 的最佳估值区间为 \kuo
	\fourch{$[36 \uppi, 52 \uppi]$}{$[36 \uppi, 100 \uppi]$}{$[52 \uppi, 100 \uppi]$}{$[9 \uppi, 25 \uppi]$}
\end{ti}

\begin{ti}
	设 $\Omega=\left\{(x, y, z) | x^{2}+y^{2}+z^{2} \leqslant 2, x \geqslant 0\right\}$ , 则以下等式错误的是 \kuo
	\fourch{$\iiint_{\Omega} x^{2} y \dd{V}=0$}{$\iiint_{\Omega}(x+y) \dd{V}=0$}{$\iiint_{\Omega} z \dd{V}=0$}{$\iiint_{\Omega} x y \dd{V}=0$}
\end{ti}

\begin{ti}
	设 $L$ 为直线 $y=y_0$ 上从点 $A(0,y_0)$ 到点 $B(3,y_0)$ 的有向直线段, 则 $\int_{L} 2 \dd{y}=$ \kuo
	\fourch{$6$}{$6y_0$}{$3y_0$}{$0$}
\end{ti}

\begin{ti}
	$\Sigma$ 为平面 $x+y+z=1$ 与三坐标面所围区域表面的外侧, 则 $\iint_{\Sigma}(2 y+3 z) \dd{y} \dd{z}+(x+2 z) \dd{z} \dd{x}+(y+1) \dd{x} \dd{y}=$ \kuo
	\fourch{$0$}{$\frac{1}{6}$}{$\frac{2}{3}$}{$\frac{5}{3}$}
\end{ti}

\begin{ti}
	交错级数 $\sum_{n=1}^{\infty}(-1)^{n-1} \frac{1}{3^{n-1}}$ \kuo
	\fourch{发散}{条件收敛}{绝对收敛}{无法确定}
\end{ti}

\subsubsection{填空题（每空 $3$ 分，共 $24$ 分）}
\begin{ti}
	以 $y_{1}=\ee^{x}, y_{2}=x \ee^{x}$ 为特解的阶数最低的常系数齐次线性微分方程是 \hua{}
\end{ti}

\begin{ti}
	直线 $L:\begin{cases}
		x=3t-2\\
		y=t+2\\
		z=2t-1
	\end{cases}$ 和平面 $\pi:2 x+3 y+3 z-8=0$ 的交点是 \hua{}
\end{ti}

\begin{ti}
	设 $z=xy^3$ , 则 $\dd{z}=$ \hua{}
\end{ti}

\begin{ti}
	交换二次积分的积分次序后, $\int_{0}^{2} \dd{y} \int_{y^{2}}^{2 y} f(x, y) \dd{x}=$ \hua{}
\end{ti}

\begin{ti}
	设 $\Omega=\{-1 \leqslant x \leqslant 1,-1 \leqslant y \leqslant 3,0 \leqslant z \leqslant 2\}$ , 则 $\iiint_{\Omega} \dd{x} \dd{y} \dd{z}=$ \hua{}
\end{ti}

\begin{ti}
	设 $L$ 为由三点 $(0,0),(3,0),(3,2)$ 围成的平面区域 $D$ 的正向边界曲线, 由格林公式知 $\int_{L}(3 x-y+4) \dd{x}+(5 y+3 x-6) \dd{y}=$ \hua{}
\end{ti}

\begin{ti}
	设 $\Sigma$ 是上半圆锥面 $z=\sqrt{x^{2}+y^{2}}(0 \leqslant z \leqslant 1)$ , 则曲面积分 $\iint_{\Sigma}\left(x^2+y^2\right)\dd{s}=$ \hua{}
\end{ti}

\begin{ti}
	级数 $\sum_{n=1}^{\infty}\left(\frac{1}{n(n+1)}-\frac{1}{2^{n}}\right)$ 的和为 \hua{}
\end{ti}

\subsubsection{综合题（$8$ 小题，共 $52$ 分）}
\begin{ti}[$6$ 分]
	求方程 $\frac{\dd{y}}{\dd{x}}=\frac{x y}{1+x^{2}}$ 的通解
\end{ti}

\begin{ti}[$6$ 分]
	设 $z=\ln \left(x^{2}-y\right)$ , 而 $y=\tan x$ , 求 $\frac{\dd{z}}{\dd{x}}$
\end{ti}

\begin{ti}[$6$ 分]
	计算 $\iint_{D}\left(x^{2}+y^{2}\right) \dd{x} \dd{y}$ , $D$ 为曲线 $x^{2}-2 x+y^{2}=0, y=0$ 围成的在第一象限的闭区域
\end{ti}

\begin{ti}[$6$ 分]
	计算三重积分 $\iiint_{\Omega} z \dd{x} \dd{y} \dd{z}$ , 其中 $\Omega$ 是由圆锥面 $z=\sqrt{x^{2}+y^{2}}$ 与球面 $z=\sqrt{2-x^{2}-y^{2}}$ 围成的区域
\end{ti}

\begin{ti}[$8$ 分]
	用高斯公式计算 $\oiint_{\Sigma}\left(a^{2} x+x^{3}\right) \dd{y} \dd{z}+y^{3} \dd{z} \dd{x}+z^{3} \dd{x} \dd{y}$ , 其中 $\Sigma$ 为球面 $x^{2}+y^{2}+z^{2}=a^{2}$ , 取外侧
\end{ti}

\begin{ti}[$8$ 分]
	用格林公式计算 $\oint_{C} x^{2} y \dd{x}-x y^{2} \dd{y}$ , 其中 $C$ 为圆周 $x^2+y^2=4$ , 取正向
\end{ti}

\begin{ti}[$6$ 分]
	判断级数 $\sum_{n=1}^{\infty} \frac{1}{2^{n-1}(2 n-1)}$ 的敛散性
\end{ti}

\begin{ti}[$6$ 分]
	在区间 $(-1,1)$ 内求幂级数 $\sum_{n=1}^{\infty} \frac{x^{n}}{n}$ 的和函数 $s(x)$
\end{ti}

\subsection{2018-2019}
\subsubsection{选择题}
\begin{ti}
	设 $z = x^y$, 则 $\frac{\partial z}{\partial y} = $ \kuo
	\fourch{$x^y$}{$yx^{y-1}$}{$x^y \ln x$}{$\frac{x^y}{\ln x}$}
\end{ti}

\begin{ti}
	下列级数中收敛的是 \kuo
	\fourch{$\sum_{n=1}^{+\infty} \frac{1}{n}$}{$\sum_{n=1}^{+\infty}(-1)^n$}{$\sum_{n=1}^{+\infty}\frac{n}{2n+1}$}{$\sum_{n=1}^{+\infty} (-1)^n \frac{1}{n}$}
\end{ti}

\begin{ti}
	设有界闭区域 $D$ 由分段光滑曲线 $L$ 所围成, $L$ 取正向, 函数 $P(x,y)$, $Q(x,y)$ 在 $D$ 上具有一阶连续偏导数, 则 $\oint_L P \dd{x} + Q \dd{y} = $ \kuo
	\fourch{
		$\iint_D \Bigl( \frac{\partial Q}{\partial x} - \frac{\partial P}{\partial y} \Bigr) \dd{x} \dd{y}$
	}{
		$\iint_D \Bigl( \frac{\partial P}{\partial x} - \frac{\partial Q}{\partial y} \Bigr) \dd{x} \dd{y}$
	}{
		$\iint_D \Bigl( \frac{\partial P}{\partial y} - \frac{\partial Q}{\partial x} \Bigr) \dd{x} \dd{y}$
	}{
		$\iint_D \Bigl( \frac{\partial Q}{\partial y} - \frac{\partial P}{\partial x} \Bigr) \dd{x} \dd{y}$
	}
\end{ti}

\begin{ti}
	非齐次线性微分方程 $x'' - 2x' + x = (t + 2) \ee^{2t}$ 的一个待定特解形式 $x^* = $ \kuo
	\fourch{$t( At+B ) \ee{2t}$}{$(At + B) \ee{2t}$}{$A t^2 \ee^{2t}$}{$t^2 (At + B) \ee{2t}$}
\end{ti}

\begin{ti}
	设 $D$ 是由 $y=x$, $y=2x$ 及 $y=2$ 围成的, 那么 $\iint_D f(x,y) \dd{x} \dd{y} = $ \kuo
	\fourch{
		$\int_0^2 \dd{y} \int_{\frac{y}{2}}^y f(x,y) \dd{x}$
	}{
		$\int_0^2 \dd{y} \int_y^{\frac{y}{2}} f(x,y) \dd{x}$
	}{
		$\int_0^2 \dd{x} \int_{2x}^2 f(x,y) \dd{y}$
	}{
		$\int_0^2 \dd{x} \int_x^{2x} f(x,y) \dd{y}$
	}
\end{ti}

\begin{ti}
	设 $\Sigma$ 是圆柱面 $x^2 + y^2 = a^2$ 与平面 $z=1$、$z=3$ 围成区域的表面, 取外侧, 则曲面积分 $\oiint_\Sigma \bigl( xy^2 + z \bigr) \dd{x} \dd{y} + \bigl( x^2 z + y \bigr) \dd{x} \dd{z} + \bigl( x + yz^2 \bigr) \dd{y} \dd{z} = $ \kuo
	\fourch{
		$2 \uppi a^2$
	}{
		$6 \uppi a^2$
	}{
		$12 \uppi a^2$
	}{$0$}
\end{ti}

\begin{ti}
	过 $(0,2,4)$ 且与两平面 $x+2z=1$ 和 $y-3z=2$ 平行的直线方程是 \kuo
	\fourch{
		$\frac{x}{1} = \frac{y-2}{0} = \frac{z-4}{2}$
	}{
		$\frac{x}{0} = \frac{y-2}{1} = \frac{z-4}{-3}$
	}{
		$\frac{x}{-2} = \frac{y-2}{-3} = \frac{z-4}{1}$
	}{
		$\frac{x}{-2} = \frac{y-2}{3} = \frac{z-4}{1}$
	}
\end{ti}

\begin{ti}
	设积分区域 $D: 1 \leqslant x^2 + y^2 \leqslant 4$, 则二重积分 $\iint_D \sqrt{x^2 + y^2} \dd{x} \dd{y} = $ \kuo
	\fourch{
		$\int_0^{2\uppi} \dd{\theta} \int_{\rho}^4 \dd{\rho}$
	}{
		$\int_0^{2\uppi} \dd{\theta} \int_1^2 \rho \dd{\rho}$
	}{
		$\int_0^{2\uppi} \dd{\theta} \int_1^2 \rho^2 \dd{\rho}$
	}{
		$\int_0^{2\uppi} \dd{\theta} \int_0^1 \rho^2 \dd{\rho}$
	}
\end{ti}

\subsubsection{填空题}
\begin{ti}
	微分方程 $y'=3y$ 的通解是 $y=$ \hua{}
\end{ti}

\begin{ti}
	已知向量 $\vec{a}$ 与 $\vec{b}$ 方向相反, 且 $\left|\vec{b}\right|=3|\vec{a}|$ , 则 $\vec{b}$ 由 $\vec{a}$ 表示为 $\vec{b}=$ \hua{}
\end{ti}

\begin{ti}
	设 $z=x^3 y+\sin(x+y)$ , 则 $\frac{\partial z}{\partial x}=$ \hua{}
\end{ti}

\begin{ti}
	交换二次积分的积分次序后, $\int_0^1 \dd{x}\int_{x^3}^{x} f(x,y)\dd{y}=$ \hua{}
\end{ti}

\begin{ti}
	设 $\Omega=\left\{ 0\leqslant x\leqslant 3, 0\leqslant y\leqslant \frac{\uppi}{2}, 0\leqslant z\leqslant 2 \right\}$ , 则 $\iiint_{\Omega}xz\sin y\dd{x}\dd{y}\dd{z}=$ \hua{}
\end{ti}

\begin{ti}
	$L$ 为曲线 $y=x^2$ 上从点 $(0,0)$ 到点 $(1,1)$ 的一段弧, 则 $\int_{L}\sqrt{y}\dd{s}=$ \hua{}
\end{ti}

\begin{ti}
	$L$ 为圆周 $x^2+y^2=a^2$ 的正向边界曲线, 由格林公式知 $\oint_{L} \bigl( x^2y\cos x+2xy\sin x-y^2 \ee^x \bigr)\dd{x} + \bigl( x^2\sin x-2y\ee^x + x \bigr)\dd{y}=$ \hua{}
\end{ti}

\begin{ti}
	若 $p$ 满足条件 \hua{}, 则级数 $\sum_{n=1}^{+\infty}\frac{1}{n^{p-3}}$ 一定发散
\end{ti}

\subsubsection{综合题}
\begin{ti}
	求过点 $(3,1,3)$ , 且与直线 $
	\begin{cases}
		x-2y+2z-7=0\\
		x+5y-z+1=0
	\end{cases}
	$ 垂直的平面方程
\end{ti}

\begin{ti}
	已知 $z=\ee^{2x-y}$ , 而 $x=\sin t,y=t^2$, 求 $\frac{\dd{z}}{\dd{t}}$
\end{ti}

\begin{ti}
	计算 $\iint_{D} x\ee^{-y^2}\dd{x}\dd{y}$, $D$ 是曲线 $x=0,y=x^2,y=\sqrt{2}$ 围成的在第一象限的闭区域
\end{ti}

\begin{ti}
	计算以 $xoy$ 面上的圆周 $x^2+y^2=4$ 围成的闭区域为底, 以旋转抛物面 $z=x^2+y^2$ 为顶的曲顶柱体的体积
\end{ti}

\begin{ti}
	计算曲面积分 $\iint_{\Sigma}\left( x^2+y^2 \right)\dd{s}$ , 其中 $\Sigma$ 为锥面 $z=\sqrt{x^2+y^2}$ 被平面 $z=2$ 截下的带锥顶部分
\end{ti}

\begin{ti}
	利用高斯公式计算 $I=\iint_{\Sigma}\left( x^3 z+2x \right)\dd{y}\dd{z}-x^2 yz\dd{z}\dd{x}-x^2 z^2\dd{x}\dd{y}$ , 其中 $\Sigma$ 是由开口向下的旋转抛物面 $z=2-x^2-y^2$ 与平面 $z=1$ 围成立体 $\Omega$ 的表面, 取外侧
\end{ti}

\begin{ti}
	判断级数 $\sum_{n=1}^{+\infty}\frac{2\cdot 4\cdot 6\cdot\cdots\cdot (2n)}{1\cdot 4\cdot 7\cdot\cdots\cdot (3n-2)}$ 的敛散性
\end{ti}

\begin{ti}
	求级数 $\sum_{n=1}^{+\infty}nx^{n+1}(x\in(-1,1))$ 的和函数
\end{ti}

\subsection{2019-2020 11A}
\subsubsection{选择题（每小题 $3$ 分，共 $24$ 分）}
\begin{ti}
	方程 $y''-3 y'+2 y=\ee^{x}$ 的待定特解 $y^*$ 的一个形式是 $y^*=$ \kuo
	\fourch{$\ee^x$}{$ax^2\ee^x$}{$a\ee^x$}{$ax\ee^x$}
\end{ti}

\begin{ti}
	过点 $(3,1,-2)$ 且通过直线 $\frac{x-4}{5}=\frac{y+3}{2}=\frac{z}{1}$ 的平面方程 \kuo
	\twoch{$5x+2y+z-15=0$}{$\frac{x-3}{8}=\frac{y-1}{-9}=\frac{z+2}{-22}$}{$8x-9y-22z-59=0$}{$\frac{x-3}{5}=\frac{y-1}{2}=\frac{z+2}{1}$}
\end{ti}

\begin{ti}
	设 $f(x,y)=\ln\left(x+\frac{y}{2x} \right)$ , 则 $f_{y}(1,0)=$ \kuo
	\fourch{$1$}{$\frac{1}{2}$}{$\frac{1}{3}$}{$0$}
\end{ti}

\begin{ti}
	$D=\{ (x,y)|0\leqslant x\leqslant 1,0\leqslant y\leqslant 2 \}$ , 利用二重积分的性质, $\iint_{D}\frac{1}{\sqrt{x^2+y^2+2xy+16}}\dd{x}\dd{y}$ 的最佳估值区间为 \kuo
	\fourch{$\left[ \frac{2}{5},\frac{1}{2} \right]$}{$\left[ \frac{1}{5},\frac{1}{2} \right]$}{$\left[ \frac{2}{5},1 \right]$}{$\left[ \frac{1}{2},1 \right]$}
\end{ti}

\begin{ti}
	$\Omega$ 由柱面 $x^2+y^2=1$ 、平面 $z=1$ 及三个坐标面围成的在第一卦限内的闭区域, 则 $\iiint_{\Omega}xy\dd{V}=$ \kuo
	\twoch{$\int_{0}^{\uppi}\dd{\theta}\int_{0}^{1}\dd{\rho}\int_{0}^{1}\rho^3\sin\theta\cos\theta\dd{z}$}{$\int_{0}^{2\uppi}\int_{0}^{1}\dd\rho\int_{0}^{1}\rho^2\sin\theta\cos\theta\dd{z}$}{$\int_{0}^{\frac{\uppi}{2}}\dd{\theta}\int_{0}^{1}\dd\rho\int_{0}^{1}\rho^2\sin\theta\cos\theta\dd{z}$}{$\int_{0}^{\frac{\uppi}{2}} \dd{\theta} \int_{0}^{1} \dd{\rho} \int_{0}^{1} \rho^{3} \sin \theta \cos \theta \dd{z}$}
\end{ti}

\begin{ti}
	设 $L$ 是 $xoy$ 平面上的有向曲线, 下列曲线积分中, \kuo 是与路径无关的
	\twoch{$\int_{L} 3 y x^{2} \dd{x}+x^{3} \dd{y}$}{$\int_{L} y \dd{x}-x \dd{y}$}{$\int_{L} 2 x y \dd{x}-x^{2} \dd{y}$}{$\int_{L} 3 y x^{2} \dd{x}+y^{3} \dd{y}$}
\end{ti}

\begin{ti}
	设 $L$ 为圆周 $\begin{cases}
		x=a\cos t\\
		y=a\sin t
	\end{cases}(0\leqslant t\leqslant 2\uppi)$ , 则 $\oint_{L}\left(x^{2}+y^{2}\right) \dd{s}=$ \kuo
	\fourch{$a^3$}{$\uppi a^3$}{$2\uppi a^3$}{$3\uppi a^3$}
\end{ti}

\begin{ti}
	下列级数中收敛的是 \kuo
	\fourch{$\sum_{n=1}^{\infty} \frac{n}{n+1}$}{$\sum_{n=1}^{\infty} \frac{1}{n \sqrt{n+1}}$}{$\sum_{n=1}^{\infty} \frac{1}{2(n+1)}$}{$\sum_{n=1}^{\infty} \frac{1}{\sqrt{n+1}}$}
\end{ti}

\subsubsection{填空题(每空 $3$ 分, 共 $24$ 分)}
\begin{ti}
	微分方程 $\frac{\dd{y}}{\dd{x}}=-3 y+\ee^{2 x}$ 的通解是 $y=$ \hua{}
\end{ti}

\begin{ti}
	平行于 $y$ 轴且通过曲线 $\begin{cases}
		x^{2}+y^{2}+4 z^{2}=1\\
		x^{2}=y^{2}+z^{2}
	\end{cases}$ 的柱面方程是 \hua{}
\end{ti}

\begin{ti}
	设 $z=x^{2} y+x y^{2}$ , 则 $\dd{z}=$ \hua{}
\end{ti}

\begin{ti}
	$\iint_{D} y^{2} \sin ^{3} x \dd{x} \dd{y}=$ \hua{}(区域 $D$ 为: $-4 \leqslant x \leqslant 4,-1 \leqslant y \leqslant 1$)
\end{ti}

\begin{ti}
	设 $D$ 为平面闭区域: $x^{2}+y^{2} \leqslant 1$ , 则 $\iint_{D} \sqrt{x^{2}+y^{2}} \dd{x} \dd{y}$ 化为极坐标系下二次积分的表达式为 \hua{}
\end{ti}

\begin{ti}
	设 $L$ 是任意一条分段光滑的有向闭曲线, 则 $\oint_{L} 2 x y \dd{x}+x^{2} \dd{y}=$ \hua{}
\end{ti}

\begin{ti}
	$I=\iint_{\Sigma}(x+z \sin y) \dd{y} \dd{z}+(y+x \sin z) \dd{z} \dd{x}+z \dd{x} \dd{y}=$ \hua{}, 其中 $\Sigma$ 为球面 $x^{2}+y^{2}+z^{2}=4(z \geqslant 0)$ 与平面 $z=0$ 围成区域的表面, 取外侧
\end{ti}

\begin{ti}
	级数 $\sum_{n=1}^{\infty}(-1)^{n} \frac{1}{n} x^{n}$ 的收敛半径为 \hua{}
\end{ti}

\subsubsection{综合题（请写出求解过程，$8$ 小题，共 $52$ 分）}
\begin{ti}[$6$ 分]
	求过点 $(2,1,1)$ , 且与直线 $\begin{cases}
		x-y+3 z-7=0\\
		3 x+5 y-2 z+1=0
	\end{cases}$ 垂直的平面方程
\end{ti}

\begin{ti}[$6$ 分]
	设 $z=f\left(\ee^{x+y}, \sin (x y)\right)$ , 且 $f$ 具有一阶连续偏导数, 求 $\frac{\partial z}{\partial x}, \frac{\partial z}{\partial y}$
\end{ti}

\begin{ti}[$8$ 分]
	计算 $\iint_{D}\left(x^{2}+y\right) \dd{x} \dd{y}$ , $D$ 是曲线 $y=x^{2}, x=y^{2}$ 围成的闭区域
\end{ti}

\begin{ti}[$6$ 分]
	计算 $\iiint_{\Omega}\left(x^{2}+y^{2}\right) \dd{x} \dd{y} \dd{z}$ , 其中 $\Omega$ 是由圆锥面 $z^{2}=x^{2}+y^{2}$ 及平面 $z=2$ 围成的闭区域
\end{ti}

\begin{ti}[$6$ 分]
	计算 $\int_{\Gamma} x^{3} \dd{x}+3 z y^{2} \dd{y}-x^{2} y \dd{z}$ , 其中 $\Gamma$ 是从点 $A(2,2,1)$ 到原点 $O$ 的直线段 $AO$
\end{ti}

\begin{ti}[$8$ 分]
	空间区域 $\Omega$ 由开口向下的旋转抛物面 $z=1-x^{2}-y^{2}$ 与平面 $z=0$ 所围, $\Omega$ 的表面取外侧为 $\Sigma$ , 利用高斯公式计算 $\oiint_{\Sigma} x^{2} y z^{2} \dd{y} \dd{z}-x y^{2} z^{2} \dd{z} \dd{x}+z(1+x y z) \dd{x} \dd{y}$
\end{ti}

\begin{ti}[$6$ 分]
	判断级数 $\sum_{n=1}^{\infty} \frac{n^{\ee}}{\ee^{n}}$ 的敛散性
\end{ti}

\begin{ti}[$6$ 分]
	求幂级数 $\sum_{n=0}^{\infty}(2 n+1) x^{2 n}(x \in(-1,1))$ 的和函数
\end{ti}

\subsection{2019-2020 5A}
\subsubsection{选择题}
\begin{ti}
	微分方程 $y' = 2y$ 的一个特解是 \kuo
	\fourch{$y = \ee^{x^2}$}{$y = \ee^{2x}$}{$y = \ee^{-\frac{x}{2}}$}{$y = \ee^{-x}$}
\end{ti}

\begin{ti}
	方程 $\frac{x^2}{4} + \frac{y^2}{4} + \frac{z^2}{9} = 1$ 在三维空间中表示 \kuo
	\fourch{椭球面}{柱面}{抛物面}{平面}
\end{ti}

\begin{ti}
	若 $z = x^4 y^3 + 2x$, 则 $\dd{z}|_{(1,2)} =$ \kuo
	\fourch{$\bigl( 4x^3 y^3 + 2 \bigr)\dd{x} + 3x^4 y^2 \dd{y}$}{$12 \dd{x} + 34 \dd{y}$}{$34 \dd{x} + 12 \dd{y}$}{$34 \dd{x} - 12 \dd{y}$}
\end{ti}

\begin{ti}
	设积分区域 $D:$ $x^2+y^2 \leqslant 1$, $y \geqslant 0$, 则二重积分 $\iint_D \ee^{\sqrt{x^2 + y^2}} \dd{x} \dd{y} =$ \kuo
	\fourch{$\int_0^{2\uppi} \dd{\theta} \int_0^1 \rho \ee^{\rho} \dd{\rho}$}{$\int_0^{2\uppi} \dd{\theta} \int_0^1 \ee^{\rho} \dd{\rho}$}{$\int_0^{\uppi} \dd{\theta} \int_0^1 \ee^{\rho} \dd{\rho}$}{$\int_0^{\uppi} \dd{\theta} \int_0^1 \rho \ee^{\rho} \dd{\rho}$}
\end{ti}

\begin{ti}
	设 $\Omega = \bigl\{ (x,y,z) | x^2 + y^2 + z^2 \leqslant 2, x \geqslant 0 \bigr\}$, 则以下等式错误的是 \kuo
	\fourch{$\iiint_{\Omega} x^2 y \dd{v} = 0$}{$\iiint_{\Omega} (x+y) \dd{v} = 0$}{$\iiint_{\Omega} z \dd{v} = 0$}{$\iiint_{\Omega} xy \dd{v} = 0$}
\end{ti}

\begin{ti}
	设 $L$ 是椭圆 $\frac{x^2}{4} + \frac{y^2}{3} = 1$, 其周长记为 $a$, 则 $\oint_L \bigl( 3x^2 + 4y^2 \bigr) \dd{S} =$ \kuo
	\fourch{$6a$}{$3a$}{$12a$}{$0$}
\end{ti}

\begin{ti}
	$L$ 为平面曲线 $x = a \cos t$, $y = b \sin t$ ($0 \leqslant t \leqslant 2\uppi$) 围成区域的边界, 取正向, 则 $\int_L y \dd{x} - x \dd{y} =$ \kuo
	\fourch{$\uppi ab$}{$-\uppi ab$}{$2\uppi ab$}{$-2\uppi ab$}
\end{ti}

\begin{ti}
	$\sum_{n=0}^{\infty} (-1)^n \frac{x^{2n}}{n!}$ 的和函数 $S(x) =$ \kuo
	\fourch{$\ee^{-x^2}$}{$\ee^{x^2}$}{$-\ee^{x^2}$}{$\cos x$}
\end{ti}

\subsubsection{填空题}
\begin{ti}
	以 $\ee^{2x}$, $\ee^{3x}$ 为解的阶数最低的常系数线性齐次微分方程是 \hua{}
\end{ti}

\begin{ti}
	过点 $(1,1,1)$ 与直线 $\begin{cases}
		3x - y + 2z + 2 = 0 \\
		x - 2y + 3z - 5 = 0
	\end{cases}$ 的平面方程 \hua{}
\end{ti}

\begin{ti}
	设 $z = 2\ee^{xy}$, 则 $\frac{\partial^2 z}{\partial x \partial y} =$ \hua{}
\end{ti}

\begin{ti}
	若 $D = \{ 0 \leqslant x \leqslant 2, 0 \leqslant y \leqslant 4 \}$, 则积分 $\iint_D xy \dd{x} \dd{y} =$ \hua{}
\end{ti}

\begin{ti}
	设 $\Omega$ 是圆锥面 $z = \sqrt{x^2+y^2}$ 和球面 $z = \sqrt{4-x^2-y^2}$ 围成, 则 $\iiint_{\Omega} y \sqrt{x^2+y^2} \dd{x} \dd{y} \dd{z} =$ \hua{}
\end{ti}

\begin{ti}
	设 $L$ 为 $x$ 轴上横坐标从 $1$ 到 $2$ 的一段, 则 $\int_L \bigl( 2y - x^2 \bigr) \dd{x} + \bigl( -2x + y^2 \bigr) \dd{y} =$ \hua{}
\end{ti}

\begin{ti}
	$\Sigma$ 为圆锥面 $x^2 + y^2 = z^2$ 与平面 $z = h$ ($h > 0$) 围成区域的表面, 取外侧. 则 $\iint_{\Sigma} (x + y - z) \dd{y} \dd{z} + (z - x) \dd{z} \dd{x} + (x - y) \dd{x} \dd{y} =$ \hua{}
\end{ti}

\begin{ti}
	级数 $\sum_{n=1}^{\infty} \frac{\sin n \alpha}{(n+1)^2}$ 是 \hua{} 收敛（填绝对或条件收敛）
\end{ti}

\subsubsection{综合题}
\begin{ti}
	求过点 $(3,-2,0)$ 且与直线 $\frac{x-1}{1} = \frac{y+1}{-1} = \frac{z-2}{1}$ 平行的直线方程
\end{ti}

\begin{ti}
	设 $z = \ee^{x+y} + \sin xy$, 求 $\frac{\partial z}{\partial x}$, $\frac{\partial^2z}{\partial x \partial y}$
\end{ti}

\begin{ti}
	计算 $\iint_D y \dd{x} \dd{y}$, $D$ 是曲线 $y = x^2$, $y = 1$ 围成的区域
\end{ti}

\begin{ti}
	求平面 $z = 2$ 与球面 $x^2 + y^2 + z^2 = 8$ 所围成在平面 $z = 2$ 之上的部分的体积 $V$
\end{ti}

\begin{ti}
	计算 $\iint_{\varSigma} \Bigl( z + 2x + \frac{4}{3}y \Bigr) \dd{S}$, 其中 $\varSigma$ 为平面 $\frac{x}{2} + \frac{y}{3} + \frac{z}{4} = 1$ 在第一卦限的部分
\end{ti}

\begin{ti}
	利用高斯公式计算 $\iint_{\varSigma} yz \dd{z} \dd{x} + 2 \dd{x} \dd{y}$, $\varSigma$ 是半球球面 $x^2 + y^2 + z^2 = 4$ ($z \geqslant 0$) 与 $z = 0$ 围成区域 $\Omega$ 的表面, 取外侧
\end{ti}

\begin{ti}
	判断级数 $\sum_{n=1}^{\infty} n \Bigl( \frac{3}{5} \Bigr)^n$ 的敛散性
\end{ti}

\begin{ti}
	求级数 $\sum_{n=1}^{\infty} n x^n$ ($x \in (-1,1)$) 的和函数
\end{ti}

\subsection{难度与考试近似的题}
\subsubsection{选择题}
\begin{ti}
	微分方程 $y'=p(x) y$ 的通解是 \kuo
	\fourch{$y=\ee^{\int p(x) \dd{x}}$}{$y=C \ee^{\int-p(x) \dd{x}}$}{$y=C \ee^{\int p(x) \dd{x}}$}{$y=C p(x)$}
\end{ti}

\begin{ti}
	已知曲线 $\begin{cases}
		x^{2}+y^{2}+z^{2}=2\\
		x+y+z=a
	\end{cases}$ 在 $yoz$ 坐标面上的投影曲线为 $\begin{cases}
		y^{2}+y z+z^{2}=1\\
		x=0
	\end{cases}$, 则 $a=$ \kuo
	\fourch{$-1$}{$0$}{$1$}{$2$}
\end{ti}

\begin{ti}
	设 $z=\ee^{y} \tan x$, 则 $\dd{z}=$ \kuo
	\twoch{$\ee^{y} \tan x \dd{x}+\ee^{y} \sec ^{2} x \dd{y}$}{$\frac{\ee^{y}}{1+x^{2}} \dd{x}+\ee^{y} \tan x \dd{y}$}{$\ee^{x} \tan y \dd{x}+\ee^{x} \sec ^{2} y \dd{y}$}{$\ee^{y} \sec ^{2} x \dd{x}+\ee^{y} \tan x \dd{y}$}
\end{ti}

\begin{ti}
	设积分区域 $D : x^{2}+y^{2} \leqslant 4$ , 则二重积分 $\iint_{D} \sqrt{x^{2}+y^{2}} \dd{x} \dd{y}=$ \kuo
	\fourch{$\int_{0}^{2 \uppi} \dd{\theta} \int_{0}^{2} \rho^{2} \dd{\rho}$}{$\int_{0}^{2 \uppi} \dd{\theta} \int_{\rho}^{4} \dd{\rho}$}{$\int_{0}^{2 \uppi} \dd{\theta} \int_{0}^{1} \rho^{2} \dd{\rho}$}{$\int_{0}^{2 \uppi} \dd{\theta} \int_{1}^{2} \rho \dd{\rho}$}
\end{ti}

\begin{ti}
	设 $\Omega$ 由圆锥面 $z=1-\sqrt{x^{2}+y^{2}}$ 与平面 $z=0$ 围成的闭区域, 则 $\iiint_{\Omega} z \dd{V}=$ \kuo
	\twoch{$\int_{0}^{\uppi} \dd{\theta} \int_{0}^{1} \rho \dd{\rho} \int_{0}^{1-\rho} z \dd{z}$}{$\int_{0}^{2 \uppi} \dd{\theta} \int_{0}^{1} \dd{\rho} \int_{0}^{1-\rho} z \dd{z}$}{$\int_{0}^{\uppi} \dd{\theta} \int_{0}^{1} \dd{\rho} \int_{0}^{1-\rho} z \dd{z}$}{$\int_{0}^{2 \uppi} \dd{\theta} \int_{0}^{1} \rho \dd{\rho} \int_{0}^{1-\rho} z \dd{z}$}
\end{ti}

\begin{ti}
	设 $L$ 为圆周 $\begin{cases}
		x=a \cos t\\
		y=a \sin t
	\end{cases}(0\leqslant t\leqslant 2\uppi)$ , 则 $\oint_{L}\left(x^{2}+y^{2}\right) \dd{s}=$ \kuo
	\fourch{$a^3$}{$\uppi a^3$}{$2\uppi a^3$}{$3\uppi a^3$}
\end{ti}

\begin{ti}
	$L$ 为平面闭区域: $-1 \leqslant x \leqslant 1,0 \leqslant y \leqslant 1$ 的正向边界, 则 $\int_{L}\left(\frac{1}{2} y+3 x \ee^{x}\right) \dd{x}-\left(\frac{1}{2} x-y \sin y\right) \dd{y}=$ \kuo
	\fourch{$-2$}{$2$}{$-1$}{$1$}
\end{ti}

\begin{ti}
	设幂级数 $\sum_{n=1}^{\infty} a_{n} x^{n}$ 的收敛半径为 $R(0<R<+\infty)$ , 则幂级数 $\sum_{n=1}^{\infty} a_{n}\left(\frac{x}{2}\right)^{n}$ 的收敛半径为 \kuo
	\fourch{$\frac{R}{2}$}{$2R$}{$R$}{$\frac{2}{R}$}
\end{ti}

\begin{ti}
	微分方程 $\frac{\dd^{2} y}{\dd{x}^{2}}-3 \frac{\dd{y}}{\dd{x}}+2 y=x \ee^{3 x}$ 的待定特解 $y^{*}$ 的一个形式是 \kuo
	\twoch{$y^{*}=(a x+b)+c \ee^{3 x}$}{$y^{*}=(a x+b)+c x \ee^{3 x}$}{$y^{*}=(a x+b) \ee^{3 x}$}{$y^{*}=(a x+b) x \ee^{3 x}$}
\end{ti}

\begin{ti}
	过点 $(3,2,-7)$  且在三坐标轴上的截距相等, 则此平面方程是 \kuo
	\fourch{$x+y+z+2=0$}{$z+y+z-2=0$}{$x-y+z-2=0$}{$x-y-z-2=0$}
\end{ti}

\begin{ti}
	设 $L$ 是平面有向曲线, 下列曲线积分中, \kuo 是与路径无关的
	\twoch{$\int_{L}\left(y \ee^{x}+x^{2}-y\right) \dd{x}+\left(x+\ee^{x}-2 y^{2}\right) \dd{y}$}{$\int_{L}(\cos x+y) \dd{x}+(x+\cos y) \dd{y}$}{$\int_{L}(\cos x-y) \dd{x}+(x+\cos y) \dd{y}$}{$\int_{L}\left(\frac{1}{2} y+3 x \ee^{x}\right) \dd{x}-\left(\frac{1}{2} x-y \sin y\right) \dd{y}$}
\end{ti}

\begin{ti}
	设 $\Sigma$ 是平面 $x=1, y=1, z=1$ 与三个坐标面围成区域的表面, 取外侧, 则曲面积分 $\iint_{\Sigma} 2 x \dd{y} \dd{z}+2 z \dd{z} \dd{x}+3 y \dd{x} \dd{y}=$ \kuo
	\fourch{$0$}{$2$}{$4$}{$7$}
\end{ti}

\begin{ti}
	级数 $1+\left(\frac{1}{2}\right)^{2}+\left(\frac{1}{3}\right)^{2}+\cdots+\left(\frac{1}{n}\right)^{2}+\cdots$ 是 \kuo
	\fourch{幂级数}{调和级数}{$p$ 级数}{等比级数}
\end{ti}

\begin{ti}
	方程 $\left(3 x^{2}+y \cos x\right) \dd{x}+\left(\sin x-4 y^{3}\right) \dd{y}=0$ 是 \kuo
	\twoch{可分离变量微分方程}{一阶线性方程}{全微分方程}{\textsf{A}、\textsf{B}、\textsf{C} 均不对}
\end{ti}

\begin{ti}
	$z=f(x, y)$ 在 $\left(x_{0}, y_{0}\right)$ 可微, 则 $\frac{\partial z}{\partial x}, \frac{\partial z}{\partial y}$ 在 $\left(x_{0}, y_{0}\right)$ \kuo
	\fourch{连续}{不连续}{不一定存在}{一定存在}
\end{ti}

\begin{ti}
	级数 $\sum_{n=2}^{\infty}\left(\frac{1}{\sqrt{n}-1}-\frac{1}{\sqrt{n}+1}\right)$ 是 \kuo
	\fourch{发散}{收敛}{条件收敛}{绝对收敛}
\end{ti}

\begin{ti}
	曲面 $z=\sqrt{x^{2}+y^{2}}$ 与平面 $z=1$ 所围立体的体积为 \kuo
	\twoch{$\iiint_{\Omega}\left(x^{2}+y^{2}\right) \dd{V}$}{$\int_{0}^{2 \pi} \dd{\theta} \int_{0}^{1} r \dd{r} \int_{r}^{1} \dd{z}$}{$\int_{-1}^{1} \dd{x} \int_{-\sqrt{1-x^{2}}}^{\sqrt{1-x^{2}}} \dd{y} \int_{0}^{x^{2}+y^{2}} \dd{z}$}{$\int_{0}^{2 \pi} \dd{\theta} \int_{0}^{1} r \dd{r} \int_{0}^{1} \dd{z}$}
\end{ti}

\begin{ti}
	方程 $y''-3 y'+2 y=3 x-\ee^{x}$ 的特解形式为 \kuo
	\fourch{$(a x+b) \ee^{x}$}{$a x+b+c x \ee^{x}$}{$a x+b+c \ee^{x}$}{$(a x+b) x \ee^{x}$}
\end{ti}

\begin{ti}
	设 $\overrightarrow{AB}$ 与 $u$ 轴的夹角为 $\frac{\uppi}{3}$ , 则 $\overrightarrow{AB}$ 在 $u$ 轴上的投影是 \kuo
	\fourch{$\overrightarrow{AB}\cos\frac{\uppi}{3}$}{$\overrightarrow{AB}\sin\frac{\uppi}{3}$}{$\left|\overrightarrow{AB}\right|\cos\frac{\uppi}{3}$}{$\left|\overrightarrow{AB}\right|\sin\frac{\uppi}{3}$}
\end{ti}

\begin{ti}
	过点 $M_{1}(3,-2,1), M_{2}(-1,0,2)$ 的直线方程是 \kuo
	\twoch{$-4(x-3)+2(y+2)+(z-1)=0$}{$\frac{x-3}{4}=\frac{y+2}{2}=\frac{z-1}{1}$}{$\frac{x+1}{4}=\frac{y}{2}=\frac{z-2}{1}$}{$\frac{x-3}{4}=\frac{y+2}{-2}=\frac{z-1}{-1}$}
\end{ti}

\begin{ti}
	直线 $\begin{cases}
		x+y+3 z=0\\
		x-y-z=0
	\end{cases}$ 与平面 $x-y-z+1=0$ 的夹角是 \kuo
	\fourch{$60^\circ$}{$0^\circ$}{$30^\circ$}{$90^\circ$}
\end{ti}

\begin{ti}
	设 $f(x,y)=\begin{cases}
		\frac{1}{xy}\sin\left(x^2y\right), & \text{当}xy\ne0,\\
		0, & \text{当}xy=0,
	\end{cases}$ 则当 $y\ne0$ 时, $f_{x}(0,y)=$ \kuo
	\fourch{$0$}{$1$}{$2$}{不存在}
\end{ti}

\begin{ti}
	曲线 $\begin{cases}
		z=\frac{x^2}{2}+\frac{y^2}{4},\\
		y=2.
	\end{cases}$ 在点 $\left( 1,2,\frac{3}{2} \right)$ 处的切线与 $x$ 轴的正向所成的倾角是 \kuo
	\fourch{$\arctan 1$}{$30^\circ$}{$60^\circ$}{$90^\circ$}
\end{ti}

\begin{ti}
	$D$ 是矩形闭区域 $0 \leqslant x \leqslant 1, \quad 0 \leqslant y \leqslant 2 \quad, I=\iint_{D}(x+y+1) \dd{x} \dd{y}$ , 利用二重积分的性质, $I$ 的最佳估计区间为 \kuo
	\fourch{$[0,1]$}{$[0,2]$}{$[1,3]$}{$2,8$}
\end{ti}

\begin{ti}
	$L$ 为 $y=x^2$ 上从 $A(1,1)$ 到 $B(0,0)$ 的一段弧, 则 $\int_{L}x\dd{y}=$ \kuo
	\fourch{$\int_{0}^{1} 2 x^{2} \dd{x}$}{$\int_{1}^{0} x \dd{y}$}{$\int_{1}^{0} 2 x^{2} \dd{x}$}{$\int_{0}^{1} \sqrt{y} \dd{y}$}
\end{ti}

\begin{ti}
	当 $\sum_{n=1}^{\infty}\left(a_{n}+b_{n}\right)$ 收敛时, $\sum_{n=1}^{\infty} a_{n}$ 与 $\sum_{n=1}^{\infty} b_{n}$ \kuo
	\fourch{可能不同时收敛}{不可能同时收敛}{必同时收敛}{必同时发散}
\end{ti}

\begin{ti}
	非齐次线性微分方程 $x''-2 x'+5 x=t \ee^{t} \sin 2 t$ 的特解形式 $x^*=$ \kuo
	\twoch{$(A t+B) \ee^{t} \sin 2 t$}{$\ee^{t}[(A t+B) \cos 2 t+(C t+D) \sin 2 t]$}{$t(A t+B) \ee^{t} \sin 2 t$}{$t \ee^{t}[(A t+B) \cos 2 t+(C t+D) \sin 2 t]$}
\end{ti}

\begin{ti}
	设向量 $\vec{a}=(1,2,3)$ 、 $\vec{b}=(2,0,1)$ , 则向量 $\vec{a}\times\vec{b}$ 在 $y$ 轴上的分向量为 \kuo
	\fourch{$5$}{$5\vec{j}$}{$-5$}{$-5\vec{j}$}
\end{ti}

\begin{ti}
	两向量 $\vec{a}$、$\vec{b}$ 平行的充要条件是 \kuo
	\fourch{$\vec{a}\cdot\vec{b}=0$}{$\vec{a}\times\vec{b}=0$}{$\vec{a}\cdot\vec{b}=\vec{0}$}{$\vec{a}\times\vec{b}=\vec{0}$}
\end{ti}

\begin{ti}
	$f(x,y)$ 在点 $(x_0,y_0)$ 处两个偏导数存在是 $f(x,y)$ 在 $(x_0,y_0)$ 处可微的 \kuo
	\fourch{必要条件}{充分条件}{充分必要条件}{以上都不是}
\end{ti}

\begin{ti}
	设上半球 $V=\left\{(x, y, z) | x^{2}+y^{2}+z^{2} \leqslant 1, z \geqslant 0\right\}$, 则以下等式错误的是 \kuo
	\fourch{$\iiint_{V} x \dd{V}=0$}{$\iiint_{V} y \dd{V}=0$}{$\iiint_{V} z \dd{V}=0$}{$\iiint_{V} xy \dd{V}=0$}
\end{ti}

\begin{ti}
	设 $f(x)=\begin{cases}
		x, & x\in[-\uppi,0)\\
		1, & x\in[0,\uppi)
	\end{cases}$ 的傅里叶级数的和函数为 $S(x)$, 则 $S(0)=$ \kuo
	\fourch{$0$}{$1$}{$\frac{1}{2}$}{$-\frac{1}{2}$}
\end{ti}

\begin{ti}
	级数 $\sum_{n=1}^{\infty}(-1)^{n} \frac{1}{n}$ \kuo
	\fourch{发散}{条件收敛}{绝对收敛}{以上都不对}
\end{ti}

\subsubsection{填空题}
\begin{ti}
	以 $\ee^x,x\ee^x$ 为解的阶数最低的常系数线性齐次微分方程是 \hua{}
\end{ti}

\begin{ti}
	过点 $A(1,-2,1)$ 且以 $\vec{n}=(1,2,3)$ 为法向量的平面方程是 \hua{}
\end{ti}

\begin{ti}
	设 $z=\sin \left(x^{2}+y\right)$ , 则 $\frac{\partial^{2} z}{\partial x \partial y}=$ \hua{}
\end{ti}

\begin{ti}
	设 $D$ 是圆环形闭区域 $1 \leqslant x^{2}+y^{2} \leqslant 4$ , 那么 $\iint_{D} \sqrt{x^{2}+y^{2}} \dd{\sigma}=$ \hua{}
\end{ti}

\begin{ti}
	设 $\Omega$ 为球体: $x^{2}+y^{2}+z^{2} \leqslant 4$ , 则 $\iiint_{\Omega} x^{2} \sin (y z) \dd{x} \dd{y} \dd{z}=$ \hua{}
\end{ti}

\begin{ti}
	$L$ 为抛物线 $x=y^2$ 上从点 $(1,-1)$ 到 $(1,1)$ 的一段弧, 则 $\int_{L} x y \dd{y}=$ \hua{}
\end{ti}

\begin{ti}
	$\oiint_{\Sigma}(x y+z) \dd{x} \dd{y}+(x z+y) \dd{x} \dd{z}+(x+y z) \dd{y} \dd{z}=$ \hua{}, 其中 $\Sigma$ 是由六张平面 $x=1,x=2,y=1,y=2,z=1,z=3$ 围成的六面体的表面, 取内侧
\end{ti}

\begin{ti}
	级数 $\frac{1}{3}+\frac{1}{\sqrt{3}}+\frac{1}{\sqrt[3]{3}}+\cdots+\frac{1}{\sqrt[n]{3}}+\cdots$ 是 \hua{}(填收敛或发散)
\end{ti}

\begin{ti}
	微分方程 $y'=p(x) y$ 的通解是 $y=$ \hua{}
\end{ti}

\begin{ti}
	设 $\vec{a}$ 与轴 $\vec{l}$ 的夹角为 $\frac{\uppi}{6}$, 且 $|\vec{a}|=4$, 则 $\Prj_{\vec{l}} \vec{a}=$ \hua{}
\end{ti}

\begin{ti}
	设 $f(x, y)=\tan \left(x y^{2}\right)$ , 则 $f_{x}(0,2)=$ \hua{}
\end{ti}

\begin{ti}
	交换二次积分次序的积分次序后, $\int_{1}^{2} \dd{x} \int_{2-x}^{\sqrt{2 x-x^{2}}} f(x, y) \dd{y}=$ \hua{}
\end{ti}

\begin{ti}
	已知 $\Omega$ 是由旋转抛物面 $z=x^{2}+y^{2}$ 与上半球面 $z=\sqrt{2-x^{2}-y^{2}}$ 围成的区域, 则 $\iiint_{\Omega} x y z \dd{x} \dd{y} \dd{z}=$ \hua{}
\end{ti}

\begin{ti}
	设 $\Sigma$ 是球面 $x^{2}+y^{2}+z^{2}=1$, 则 $\iint_{\Sigma}\left(x^{2}+y^{2}+z^{2}\right) \dd{s}=$ \hua{}
\end{ti}

\begin{ti}
	积分 $\oint_{L}\left(x^{2}-y\right) \dd{x}+\left(y^{2}+x\right) \dd{y}=$ \hua{}, 其中 $L$ 为圆周 $(x-1)^{2}+y^{2}=a^{2}$ 的正向
\end{ti}

\begin{ti}
	级数 $\sum_{n=1}^{\infty}(-1)^{n-1} \frac{1}{\sqrt{n}}$ 是 \hua{}收敛(填条件收敛或绝对收敛)
\end{ti}

\begin{ti}
	设 $z=x^{y}$ , 则 $\frac{\partial z}{\partial y}=$ \hua{}
\end{ti}

\begin{ti}
	积分 $\iint_{D} x y \dd{x} \dd{y}=$ \hua{}, 其中 $D$ 为 $0 \leqslant x \leqslant 2, 0 \leqslant y \leqslant 4$
\end{ti}

\begin{ti}
	$L$ 为 $y=x^2$ 点 $(0,0)$ 到 $(1,1)$ 的一段弧, 则 $\int_{L} \sqrt{y} \dd{s}=$ \hua{}
\end{ti}

\begin{ti}
	级数 $\sum_{n=1}^{\infty} \frac{(-1)^{n}}{n^{p}}$ 当 $p$ 满足 \hua{}时条件收敛
\end{ti}

\begin{ti}
	方程 $y \ee^{x} \dd{x}-\left(1+\ee^{x}\right) \dd{y}=0$ 的通解为 \hua{}
\end{ti}

\begin{ti}
	设 $z=\ln \sqrt{1+x^{2}+y^{2}}$, 则 $\dd\left.z\right|_{(1,1)}=$ \hua{}
\end{ti}

\begin{ti}
	函数 $z=x^{2}+y^{2}$ 在点 $P(1,2)$ 沿从点 $(1,2)$ 到点 $(2,2+\sqrt{3})$ 的方向上的方向导数为 \hua{}
\end{ti}

\begin{ti}
	改换二次积分的积分次序: $\int_{0}^{1} \dd{y} \int_{0}^{y} f(x, y) \dd{x}=$ \hua{}
\end{ti}

\begin{ti}
	平面 $x+y+z=1$ 含在圆柱面 $x^{2}+y^{2}=2 x$ 内部的那部分平面面积为 \hua{}
\end{ti}

\begin{ti}
	$L$ 为圆周 $x^{2}+y^{2}=1$ , 则 $\int_{L}\left(x^{2}+y^{2}\right) \dd{s}=$ \hua{}
\end{ti}

\begin{ti}
	$\Sigma$ 是 $xoy$ 平面上的圆域: $x^{2}+y^{2} \leqslant 1$, 取下侧, 则 $\iint_{\Sigma} \dd{x} \dd{y}=$ \hua{}
\end{ti}

\begin{ti}
	级数 $\sum_{n=1}^{\infty} \frac{3^{n}+4^{n}}{7^{n}}$ 的和为 \hua{}
\end{ti}

\begin{ti}
	$\ee^{x^{2}}$ 的 $x$ 的幂级数展开式为 \hua{}
\end{ti}

\begin{ti}
	微分方程 $x'''-2x''-x'+2x=0$ 的通解是 \hua{}
\end{ti}

\begin{ti}
	过点 $(1,0,1)$ 及以 $(1,2,3)$ 为方向向量的直线的对称式方程为 \hua{}
\end{ti}

\begin{ti}
	函数 $z=x^y$ 的全微分 $\dd{z}=$ \hua{}
\end{ti}

\begin{ti}
	二元函数 $u=x^{2}-x y+y^{2}$ 在点 $(-1,1)$ 处沿方向 \hua{}的方向导数最大
\end{ti}

\begin{ti}
	交换二次积分的次序 $\int_0^1 \dd{y}\int_{-1}^{-y} f(x,y)\dd{x}=$ \hua{}
\end{ti}

\begin{ti}
	若 $L$ 为抛物线 $y^2=2x$ 上介于 $(2,-2)$ 与 $(2,2)$ 两点间的曲线段, 则 $\int_{L} y\dd{s}=$ \hua{}
\end{ti}

\begin{ti}
	若 $\Sigma$ 是曲面 $x^2+y^2+z^2=1$, 则 $\iint_{\Sigma} \dd{s}=$ \hua{}
\end{ti}

\begin{ti}
	函数 $f(x)=3^x$ 的幂级数展开式为 \hua{}
\end{ti}

\subsubsection{综合题}
\begin{ti}[$6$ 分]
	求过点 $(2,0,-3)$ , 且过直线 $\begin{cases}
		x-2 y+4 z-7=0\\
		3 x+5 y-2 z+1=0
		\end{cases}$ 垂直的平面方程
\end{ti}

\begin{ti}[$6$ 分]
	设 $z=x^{y}(x>0)$ , 求 $\frac{\partial z}{\partial x}, \frac{\partial^{2} z}{\partial x \partial y}$
\end{ti}

\begin{ti}[$6$ 分]
	计算 $\iint_{D} x^{2} y^{2} \dd{x} \dd{y}$ , 其中 $D=\{(x, y) | 0 \leqslant x \leqslant 1,0 \leqslant y \leqslant 1\}$
\end{ti}

\begin{ti}[$6$ 分]
	计算 $I=\iiint_{\Omega}\left(x^{2}+y^{2}\right) \dd{V}$, 其中 $\Omega$ 为旋转抛物面 $z=x^{2}+y^{2}$ 与平面 $z=4$ 所围成的区域
\end{ti}

\begin{ti}[$8$ 分]
	$L$ 是圆环区域 $D : 1 \leqslant x^{2}+y^{2} \leqslant 4$ 的正向边界曲线, 计算曲线积分 $\oint_{L} \sqrt{x^{2}+y^{2}} \dd{x}+\Bigl[x y^{2}+y \ln \Bigl(x+\sqrt{x^{2}+y^{2}}\Bigr)\Bigr] \dd{y}$
\end{ti}

\begin{ti}[$8$ 分]
	计算 $\iint_{\Sigma} \frac{2}{z} \dd{s}$, 其中 $\Sigma$ 是球面 $x^{2}+y^{2}+z^{2}=1$ 在平面 $z=\frac{1}{2}$ 上方的部分
\end{ti}

\begin{ti}[$6$ 分]
	判断级数 $\sum_{n=1}^{\infty} \frac{3^{n}}{n \cdot 2^{n}}$ 的敛散性
\end{ti}

\begin{ti}[$6$ 分]
	求幂级数 $\sum_{n=0}^{\infty}(n+1) x^{n}$ 在收敛域 $(-1,1)$ 的和函数 $s(x)$
\end{ti}

\begin{ti}[$6$ 分]
	求过点 $(3,-2,1)$, 且与直线 $\frac{x-1}{1}=\frac{y+1}{1}=\frac{z-2}{3}$ 平行的直线方程
\end{ti}

\begin{ti}[$6$ 分]
	设 $ z = e^{xy} + \cos(x + y)$, 求 $\dd{z}$
\end{ti}

\begin{ti}[$6$ 分]
	计算 $\iint_{D}\frac{y}{x}\dd{x}\dd{y}$, $D$ 是由直线 $ y = 2x,y = x, x = 2, x = 4$ 围成的闭区域
\end{ti}

\begin{ti}[$6$ 分]
	计算 $\iiint_{\Omega}z\dd{x}\dd{y}\dd{z}$, 其中 $\Omega$ 由平面 $z = 3$ 与旋转抛物面 $x^2 + y^2 = 3z$ 围成的区域
\end{ti}

\begin{ti}[$6$ 分]
	计算 $\int_{L} 2 x y \dd{x}+x^{2} \dd{y}$, $L$ 为抛物线 $y=x^{2}$ 上从 $O(0,0)$ 到 $B(1,1)$ 的一段弧
\end{ti}

\begin{ti}[$8$ 分]
	利用高斯公式计算 $\oiint_{\Sigma} 2 x z \dd{y} \dd{z}+y z \dd{z} \dd{x}-z^{2} \dd{x} \dd{y}$, 其中 $\Sigma$ 为由上半圆锥面 $z=\sqrt{x^{2}+y^{2}}$ 与上半球面 $z=\sqrt{2-x^{2}-y^{2}}$ 所围立体 $\Omega$ 的表面, 取外侧
\end{ti}

\begin{ti}[$6$ 分]
	判断级数 $\sum_{n=1}^{\infty} n 2^{n}$ 的敛散性
\end{ti}

\begin{ti}[$6$ 分]
	求幂级数 $\sum_{n=0}^{\infty}(n+1) x^{n}$ 在收敛域 $(-1,1)$ 的和函数 $s(x)$
\end{ti}

\begin{ti}[$8$ 分]
	$z=f\left(y^{2}-x^{2}\right)$, 其中 $f(u)$ 有连续的二阶偏导数, 求 $\frac{\partial^{2} z}{\partial x^{2}}$
\end{ti}

\begin{ti}[$8$ 分]
	计算 $\int_{L}\left(\ee^{x} \sin y-2 y\right) \dd{x}+\left(\ee^{x} \cos y-2\right) \dd{y}$, $L$ 为由点 $A(1,0)$ 到 $B(0,1)$, 再到 $C(-1,0)$ 的有向折线
\end{ti}

\begin{ti}[$8$ 分]
	计算 $\oiint_{\Sigma} x y^{2} \dd{y} \dd{z}+y z^{2} \dd{z} \dd{x}+z x^{2} \dd{x} \dd{y}$, 其中 $\Sigma$ 为球体 $x^{2}+y^{2}+z^{2} \leqslant 4$ 及锥体 $z=\sqrt{x^{2}+y^{2}}$ 的公共部分的外表面
\end{ti}

\begin{ti}[$8$ 分]
	求级数 $\sum_{n=2}^{\infty} 2 n x^{n}$ 的收敛域及和函数
\end{ti}

\begin{ti}[$8$ 分]
	计算曲面积分 $\iint_{\Sigma}\left(x^{2}+y^{2}\right) \dd{s}$, 其中 $\Sigma$ 为锥面 $z=\sqrt{3\left(x^2+y^2\right)}$ 被平面 $z=3$ 截下的带锥顶的部分
\end{ti}

\begin{ti}[$7$ 分]
	求函数 $z=x^2+y^2$ 在适合条件 $\frac{x}{2}+\frac{y}{3}=1$ 下的极小值
\end{ti}

\begin{ti}[$8$ 分]
	求方程 $y^{\prime \prime}-3 y^{\prime}+2 y=3 \ee^{x}$ 的通解
\end{ti}

\begin{ti}[$7$ 分]
	把 $f(x)=x,(0<x<\uppi)$ 展开为余弦级数
\end{ti}

\begin{ti}[$8$ 分]
	已知曲线积分 $\int_{(0,0)}^{(x, y)}\left[\ee^{x}(x+1)^{n}+\frac{n}{x+1} f(x)\right] y \dd{x}+f(x) \dd{y}$ 与路径无关, 其中 $f(x)$ 可微, $f(0)=0$, 试确定 $f(x)$ , 并计算曲线积分的值
\end{ti}

\begin{ti}[$5$ 分]
	求过点 $(2,5,-3)$ 且与直线 $\begin{cases}
		x=5-2t\\
		y=1+t\\
		z=7
	\end{cases}$ 垂直的平面方程
\end{ti}

\begin{ti}[$5$ 分]
	由 $\ee^{x}-x y z=0$ 确定了函数 $z=z(x,y)$, 求 $\frac{\partial z}{\partial x}$
\end{ti}

\begin{ti}[$5$ 分]
	计算 $I=\iint_{D}\left(x^{2}+y^{2}\right) \dd{x} \dd{y}$, 其中 $D=\left\{(x, y) | 1 \leqslant x^{2}+y^{2} \leqslant 4\right\}$
\end{ti}

\begin{ti}[$8$ 分]
	利用格林公式, 计算 $\oint_{L}\left(2 x^{2} y-2 y\right) \dd{x}+\left(\frac{1}{3} x^{3}-2 x\right) \dd{y}$, 其中 $L$ 为以 $y=x,y=x^{2}$, 围成区域的正向边界
\end{ti}

\begin{ti}[$8$ 分]
	设 $\Sigma$ 是由旋转抛物面 $z=x^{2}+y^{2}$ 与平面 $z=2$ 所围成的封闭曲面, 取外侧. 用高斯公式计算 $\iint_{\Sigma} 4\left(1-y^{2}\right) \dd{z} \dd{x}+z(8 y+1) \dd{x} \dd{y}$
\end{ti}

\begin{ti}[$8$ 分]
	求幂级数 $\sum_{n=0}^{\infty} \frac{x^{2 n+1}}{2 n+1}$ 在收敛域 $(-1,1)$ 内的和函数
\end{ti}

\begin{ti}[$8$ 分]
	求微分方程 $y''-2 y'+y=\ee^{x}$ 的通解
\end{ti}

\begin{ti}[$5$ 分]
	设函数 $f(x)$ 在 $[a,b]$ 上连续且 $f(x)>0$, 证明 $\int_{a}^{b} f(x) \dd{x} \int_{a}^{b} \frac{1}{f(x)} \dd{x} \geq(b-a)^{2}$
\end{ti}

\begin{ti}[$7$ 分]
	设 $u=f(x,xy)$, $f$ 具有二阶连续偏导数, 求 $\frac{\partial u}{\partial x},\frac{\partial^2u}{\partial x\partial y}$
\end{ti}

\begin{ti}[$7$ 分]
	求曲面 $\ee^{z}-z+x y=3$ 在点 $(2,1,0)$ 处的切平面及法线方程
\end{ti}

\begin{ti}[$10$ 分]
	设 $\Omega$ 是曲面 $\Sigma_1:z=\sqrt{x^2+y^2}$ 与 $\Sigma_2:2-x^2-y^2$ 所围成的立体, 求 $\Omega$ 的体积 $V$ 与表面积 $S$
\end{ti}

\begin{ti}[$10$ 分]
	计算 $\iint_{\Sigma}\left(z+x y^{2}\right) \dd{y} \dd{z}+\left(y z^{2}-x z\right) \dd{z} \dd{x}+\left(x^{2} z+x^{3}\right) \dd{x} \dd{y}$, 其中 $\Sigma$ 为 $x^{2}+y^{2}+z^{2}=4(z \leqslant 0)$, 取下侧
\end{ti}

\begin{ti}[$10$ 分]
	计算 $\int_{L}\left(2 x y^{3}-y^{2} \cos x\right) \dd{x}+\left(1-2 y \sin x+3 x^{2} y^{2}\right) \dd{y}$, 其中 $L$ 为抛物线 $2x=\uppi y^2$ 从点 $O(0,0)$ 到点 $A\left( \frac{\uppi}{2},1 \right)$ 的一段弧
\end{ti}

\begin{ti}[$8$ 分]
	求幂级数 $\sum_{n=1}^{\infty}nx^{n-1}$ 的收敛域与和函数
\end{ti}
	\chapter{《线性代数》试卷汇总}
\section{秋}
\subsection{2018-2019 B}
\subsubsection{选择题}
\begin{ti}
	行列式 $\begin{vmatrix}
		0 & a & 0 & 0 \\
		b & c & 0 & 0 \\
		0 & 0 & e & d \\
		0 & 0 & 0 & f
	\end{vmatrix}$ 的值为 \kuo
	\fourch{$abcdef$}{$abcdf$}{$-abef$}{$cdf$}
\end{ti}

\begin{ti}
	设 $\bA$ 为 $n$ 阶非对称方阵, 则下列不是对称矩阵的是 \kuo
	\fourch{$\bA \bA^{\TT}$}{$\bA - \bA^{\TT}$}{$\bA + \bA^{\TT}$}{$\bA \bA^{\TT}$}
\end{ti}

\begin{ti}
	若 $n$ 个 $m$ 维向量线性无关, 则 \kuo
	\twoch{去掉一个向量后仍线性无关}{再增加一个向量后也线性无关}{其中只有一个向量不能被其他向量线性表示}{$n > m$}
\end{ti}

\begin{ti}
	设 $\bA = (a_{ij})_{m \times n}$, 且 $\bA$ 的行列式 $|\bA| = 0$, 但 $\bA$ 中存在某元素的代数余子式非零, 则齐次线性方程组 $\bA \bX = \symbfit 0$ 的基础解系中解向量个数为 \kuo
	\fourch{$k$}{$l$}{$1$}{$n$}
\end{ti}

\begin{ti}
	设 $\bA$ 为 $n$ 阶矩阵, 以下结论中, \kuo{} 成立
	\onech{$\bA$ 的特征向量的线性组合仍为特征向量}{$\bA$ 的特征向量即为齐次方程 $(\lambda \bE - \bA) \bX = \symbfit 0$ 的全部解}{$\bA$ 与 $\bA^{\TT}$ 有相同特征向量}{矩阵 $\bA$ 的属于特征值 $\lambda$ 的特征向量也是矩阵 $\bA^{2}$ 的属于特征值 $\lambda^{2}$ 的特征向量}
\end{ti}

\begin{ti}
	下列结论错误的是 \kuo
	\onech{同一矩阵相同特征值对应特征向量必定线性无关}{若几阶方阵有几个互不相同的特征值, 则该矩阵必可对角化}{相似矩阵的特征值相同, 特征多项式也相同}{若几阶方阵有几个线性无关的特征向量, 则该矩阵必可对角化}
\end{ti}

\subsubsection{填空题}
\begin{ti}
	设 $\bA$ 为 $4$ 阶方阵, 且 $|\bA| = 3$, 则 $|-2 \bA| = $ \hua
\end{ti}

\begin{ti}
	行列式 $\begin{vmatrix}
		2 & 1 & -4 \\
		1 & 3 & 2 \\
		-3 & 1 & 5
	\end{vmatrix}$ 的元素 $a_{21} = 1$ 的代数余子式 $\bA_{21} = $ \hua
\end{ti}

\begin{ti}
	设向量 $(2,-3,4)$ 与向量 $(-4,6,a)$ 线性相关, 则 $a = $ \hua
\end{ti}

\begin{ti}
	设 $\bA$ 为 $3$ 阶方阵, 且 $|\bA| = 4$, 则 $\bigl|\bA^{\TT} \bA\bigr| = $ \hua
\end{ti}

\begin{ti}
	设 $\balpha = (1,1,1,1)^{\TT}$, 则向量 $\balpha$ 的长度 $\| \balpha \| = $ \hua
\end{ti}

\begin{ti}
	设 $\bA$ 为 $m \times n$ 矩阵, $\symbfit B$ 为 $\bA \bX = \symbfit b$ 所对应的增广矩阵, 则 $\bA \bX = \symbfit b$ 有唯一解的充分必要条件是 \hua
\end{ti}

\begin{ti}
	已知方阵 $n$ 个特征值 $\lambda_{1},\lambda_{2},\cdots,\lambda_{n}$, 则 $\bigl|\bA^{2}\bigr| = $ \hua
\end{ti}

\subsubsection{大题}
\begin{ti}
	计算 $\begin{vmatrix}
		1 & 2 & 3 & 4 \\
		3 & 2 & 6 & -2 \\
		4 & 3 & 0 & -1 \\
		3 & 0 & 0 & 0
	\end{vmatrix}$
\end{ti}

\begin{ti}
	用分块法求矩阵 $\bA = \begin{pmatrix}
		1 & 0 & 0 & 0 \\
		0 & 1 & 4 & 0 \\
		0 & 0 & -1 & 0 \\
		0 & 0 & 0 & 8
	\end{pmatrix}$ 的逆矩阵 $\bA^{-1}$
\end{ti}

\begin{ti}
	求向量组 $\balpha_{1} = (1,1,4,3)^{\TT},\balpha_{2} = (1,-1,-2,5)^{\TT},\balpha_{3} = (-3,2,3,-14)^{\TT},\balpha_{4} = (1,3,10,1)^{\TT}$ 的秩及最大无关组
\end{ti}

\begin{ti}
	求 $\bA = \begin{bmatrix}
		2 & -2 & 0 \\
		-2 & 1 & -2 \\
		0 & -2 & 0
	\end{bmatrix}$ 的特征值和特征向量
\end{ti}

\begin{ti}
	求线性方程组 $\begin{cases}
		2x_{1} + 3x_{2} + x_{3} = 1 \\
		x_{1} + x_{2} - 2x_{3} = 1 \\
		x_{1} + 3x_{2} + 8x_{3} = -1
	\end{cases}$ 的通解
\end{ti}

\begin{ti}
	设 $b_{1}$、$b_{2}$、$b_{3}$ 是互不相同的数, 且 $\balpha_{1} = \bigl(1,b_{1},b_{1}^{2}\bigr)^{\TT},\balpha_{2} = \bigl(1,b_{2},b_{2}^{2}\bigr)^{\TT},\balpha_{3} = \bigl(1,b_{3},b_{3}^{2}\bigr)^{\TT}$, 证明：$\balpha_{1}$、$\balpha_{2}$、$\balpha_{3}$ 线性无关
\end{ti}

\subsection{2019-2020 9A}
\subsubsection{选择题}
\begin{ti}
	行列式 $\begin{vmatrix} 1 & x & 0 \\ -1 & 1-x & y \\ 0 & -1 & 1-y \end{vmatrix}$ 的值为 \kuo
	\fourch{$xy$}{1}{$x$}{$y$}
\end{ti}

\begin{ti}
	下述论断正确的是 \kuo
	\onech{如果仅当 $k_1 = k_2 = \cdots = k_r = 0$, 有 $k_1 \alpha_1 + k_2 \alpha_2 + \cdots + k_r \alpha_r = 0$, 则 $\alpha_1$, $\alpha_2$, $\cdots$, $\alpha_r$ 线性无关}
	{若 $\alpha_1$, $\alpha_2$, $\cdots$, $\alpha_r$ 线性相关, 则存在不全为 0 的数 $k_1$, $k_2$, $k_r$, 使得 $k_1 \alpha_1 + k_2 \alpha_2 + \cdots + k_r \alpha_r = 0$}
	{若 $\alpha_1$, $\alpha_2$, $\cdots$, $\alpha_r$ 线性无关, $\beta_1$, $\beta_2$, $\cdots$, $\beta_s$ 线性无关, 则 $\alpha_1$, $\alpha_2$, $\cdots$, $\alpha_r$, $\beta_1$, $\beta_2$, $\cdots$, $\beta_s$ 线性无关}
	{若 $\alpha_1$, $\alpha_2$, $\cdots$, $\alpha_r$ 线性无关, 则 $\alpha_1$, $\alpha_2$, $\cdots$, $\alpha_{r-1}$ 线性相关}
\end{ti}

\begin{ti}
	设 $|A|,|B|$ 均为 $n$ ($n>2$) 阶行列式, 则 \kuo
	\fourch{$|A+B| = |A| + |B|$}{$|A-B| = |A| - |B|$}{$|A|\cdot|B| = |A\cdot B|$}{$\begin{vmatrix}
		O & A \\ B & O
	\end{vmatrix} = |A| \cdot |B|$}
\end{ti}

\begin{ti}
	已知 $\alpha_1,\alpha_2,\alpha_3$ 是非齐次线性方程组 $Ax=b$ 的三个不同的解, 则向量组 $\alpha_1-\alpha_2,\alpha_1+\alpha_2-2\alpha_3,\frac{2}{3}(\alpha_2-\alpha_1),\allowbreak\alpha_1-3\alpha_2+2\alpha_3$ 是对应的齐次线性方程组 $Ax=0$ 的解向量共有 \kuo{} 个
	\fourch{4 个}{3 个}{2 个}{1 个}
\end{ti}

\begin{ti}
	已知三阶方阵 $A$ 的三个特征值分别为 $1,2,3$, 则 $\bigl|A^2\bigr| = $ \kuo
	\fourch{$-6$}{6}{$-36$}{36}
\end{ti}

\begin{ti}
	已知 $A = \begin{vmatrix}
		2 & 0 & 0 \\ 0 & 0 & 1 \\ 0 & 1 & x
	\end{vmatrix}$ 和 $B = \begin{vmatrix}
		2 & 0 & 0 \\ 0 & y & 0 \\ 0 & 0 & 1
	\end{vmatrix}$ 相似, 则 \kuo
	\fourch{$x=0,y=0$}{$x=0,y=1$}{$x=1,y=1$}{$x=1,y=2$}
\end{ti}

\subsubsection{填空题}
\begin{ti}
	设 $\alpha_1,\alpha_2,\alpha_3,\beta_1,\beta_2$ 都是 4 维列向量, 且 4 阶行列式 $|(\alpha_1,\alpha_2,\alpha_3,\beta_1)| = m$, $|(\alpha_1,\alpha_2,\beta_2,\alpha_3)| = n$, 则 4 阶行列式 $|(\alpha_1,\alpha_2,\alpha_3,\beta_1+\beta_2)| = $ \hua
\end{ti}

\begin{ti}
	已知 $D$ 为 4 阶行列式, $D$ 的第二行元素依次为 $a_{21} = -2,a_{22} = 1, a_{23} = 0,a_{24} = -2$, 对应的余子式依次为 $M_{21} = 1, M_{22} = 3, M_{23} = 5, M_{24} = 3$, 则 $D = $ \hua
\end{ti}

\begin{ti}
	若向量 $\alpha_1 = (a,0,b), \alpha_2 = (1,1,1), \alpha_3 = (1,3,2)$ 线性相关, 则常数 $a,b$ 满足关系式 \hua
\end{ti}

\begin{ti}
	矩阵 $A = \begin{bmatrix}
		1 & 0 \\ 0 & 2
	\end{bmatrix}$, 则 $A^x + E = $ \hua
\end{ti}

\begin{ti}
	$n$ 元齐次方程组 $Ax = 0$ 的系数矩阵的秩为 $n-3$, 则基础解系所含解向量的个数为 \hua
\end{ti}

\begin{ti}
	已知向量 $\alpha = (1,2,2,3)', \beta = (3,1,5,2)'$, 则它们的内积 $[\alpha,\beta] = $ \hua
\end{ti}

\begin{ti}
	设二阶方阵 $A,B$ 相似, $A$ 的特征值为 $2,3$, 则 $B^{-1}$ 的特征值为 \hua
\end{ti}

\subsubsection{计算题}
\begin{ti}
	计算 4 阶行列式 $\begin{vmatrix}
		1 & 0 & -1 & 1 \\
		2 & 5 & 3	 & 1 \\
		-1 & 7 & 0 & 1 \\
		1 & 8 & 1 & 1
	\end{vmatrix}$
\end{ti}

\begin{ti}
	求向量组 $\alpha_1 = (-1,1,7,10)^\TT$, $\alpha_2 = (2,1,1,1)^\TT$, $\alpha_3 = (3,1,-1,-2)^\TT$, $\alpha_4 = (8,5,9,11)^\TT$ 的秩以及它的一个最大无关组
\end{ti}

\begin{ti}
	求矩阵 $X$ 满足 $\begin{bmatrix}
		1 & 1 & -1 \\
		-2 & 1 & 1 \\
		1 & 1 & 1
	\end{bmatrix} X = \begin{bmatrix}
		2 \\ 3 \\ 6
	\end{bmatrix}$
\end{ti}

\begin{ti}
	求线性方程组 $\begin{cases}
		2x_1 + x_2 - x_3 = 1 \\
		x_1 - x_2 + x_3 = 2 \\
		4x_1 + 3x_2 - 3x_3 = -1
	\end{cases}$ 的通解
\end{ti}

\begin{ti}
	求矩阵 $A = \begin{bmatrix}
		-2 & 0 & 0 \\
		0 & 1 & 0 \\
		1 & 1 & 4
	\end{bmatrix}$ 的特征值和特征向量
\end{ti}

\subsubsection{证明题}
\begin{ti}
	已知 $\alpha_1,\alpha_2,\alpha_3$ 是 $Ax=0$ 的一个基础解系, 证明：$\beta_1 = \alpha_1,\beta_2 = \alpha_1 + 2\alpha_2, \beta_3 = 2\alpha_3 + \alpha_2$ 也是该方程组的一个基础解系
\end{ti}

\section{春}
\subsection{2017-2018}
\subsubsection{选择题}
\begin{ti}
	行列式 $\begin{vmatrix}
		2 & 1 & 1 & 1 \\
		1 & 2 & 1 & 1 \\
		1 & 1 & 2 & 1 \\
		1 & 1 & 1 & 2
	\end{vmatrix}$ 的值为 \kuo
	\fourch{$-5$}{$5$}{$9$}{$-9$}
\end{ti}

\begin{ti}
	齐次线性方程组 $\begin{cases}
		2x_{1} + 3x_{2} + 4x_{3} - 5x_{4} = 0 \\
		x_{2} - 2x_{3} - 3x_{4} = 0
	\end{cases}$ 的基础解系中含有解向量的个数是 \kuo
	\fourch{$2$}{$1$}{$3$}{$4$}
\end{ti}

\begin{ti}
	设 $\bA,\bB,\bC$ 均为 $n$ 阶方阵, 若由 $\bA \bB = \bA \bC$ 能推出 $\bB = \bC$, 则 $\bA$ 应满足下列条件中 \kuo
	\fourch{$\bA = \symbfit 0$}{$\bA \ne \symbfit 0$}{$\bA$ 不可逆}{$\bA$ 可逆}
\end{ti}

\begin{ti}
	若向量组 $\balpha_{1} = (1,2,0),\balpha_{2} = (2,1,3),\balpha_{3} = (t,2,-1)$ 线性相关, 则 $t$ 的值为 \kuo
	\fourch{$\frac{1}{2}$}{$-\frac{1}{2}$}{$1$}{$-1$}
\end{ti}

\begin{ti}
	设 $0$ 是矩阵 $\bA = \begin{pmatrix}
		1 & 0 & 2 \\
		0 & 3 & 2 \\
		1 & 0 & a
	\end{pmatrix}$ 的特征值, 则 $a = $ \kuo
	\fourch{$1$}{$-1$}{$2$}{$-2$}
\end{ti}

\begin{ti}
	设三阶方阵 $\bA$ 的三个特征值为 $1,2,3$, 则 $\bA^{\astt}$ 的三个特征值为 \kuo
	\fourch{$6,12,36$}{$6,3,2$}{$1,\frac{1}{2},\frac{1}{3}$}{$36,18,12$}
\end{ti}

\subsubsection{填空题}
\begin{ti}
	已知下三角行列式 $\bD$ 的主对角线上有一个元素为零, 则 $\bD = $ \hua
\end{ti}

\begin{ti}
	已知 $\bD$ 为四阶行列式, $\bD$ 的第 $4$ 行元素依次为 $a_{41} = -3, a_{42} = 0, a_{43} = -2, a_{44} = 2$, 对应的余子式依次为 $M_{41} = 1, M_{42} = 7, M_{43} = 1, M_{44} = 3$, 则 $\bD = $ \hua
\end{ti}

\begin{ti}
	当 $m$ \hua{} (填“$>$”或“$<$”) $n$ 时, $n$ 个 $m$ 维向量组成的向量组一定线性相关
\end{ti}

\begin{ti}
	设 $\bA = \begin{pmatrix}
		1 & -1 & 1 \\
		1 & 1 & -1
	\end{pmatrix}, \bB = \begin{pmatrix}
		1 & 2 & 3 \\
		-1 & -2 & 4
	\end{pmatrix}$, 则 $\bA + 3 \bB = $ \hua
\end{ti}

\begin{ti}
	若 $n$ 元齐次线性方程组 $\bA \symbfit x = \symbfit 0$ 有唯一零解的充要条件为 \hua
\end{ti}

\begin{ti}
	向量 $(0,0,1,1)^{\TT}$ 与 $(1,0,-1,1)^{\TT}$ 的夹角为 \hua
\end{ti}

\begin{ti}
	$(\lambda \bE - \bA) \symbfit x = \symbfit 0$ 的非零解向量都是 $\bA$ 的属于 \hua{} 的特征向量
\end{ti}

\subsubsection{计算题}
\begin{ti}[10 分]
	计算 $4$ 阶行列式 $\begin{bmatrix}
		2 & -2 & 0 & 0 \\
		0 & 2 & -2 & 0 \\
		0 & 0 & 2 & -2 \\
		-2 & 0 & 0 & 2
	\end{bmatrix}$
\end{ti}

\begin{ti}[10 分]
	已知矩阵 $\bA = \begin{pmatrix}
		2 & 3 & 0 & 0 \\
		-1 & 2 & 0 & 0 \\
		0 & 0 & 7 & 2 \\
		0 & 0 & 3 & 1
	\end{pmatrix}$, 利用分块法求 $\bA$ 的逆矩阵
\end{ti}

\begin{ti}[10 分]
	求向量组 $\balpha_{1} = (1,2,3,0)^{\TT}, \balpha_{2} = (-1,-2,0,3)^{\TT}, \balpha_{3} = (2,4,6,0)^{\TT}, \balpha_{4} = (1,-2,-1,0)^{\TT}, \balpha_{5} = (0,2,\allowbreak 4,2)^{\TT}$ 的秩及向量组的一个最大无关组
\end{ti}

\begin{ti}[13 分]
	求矩阵 $\bA = \begin{pmatrix}
		3 & 2 & 0 \\
		2 & 3 & 0 \\
		0 & 0 & 2
	\end{pmatrix}$ 的特征值和特征向量
\end{ti}

\begin{ti}[12 分]
	求线性方程组 $\begin{cases}
		2x_{1} + 5x_{2} + 2x_{3} = 2 \\
		x_{1} + 2x_{2} - x_{3} = 3 \\
		-x_{1} - x_{2} + 5x_{3} = -7
	\end{cases}$ 的通解（要求用它的一个特解和对应的齐次线性方程组的基础解系表示）
\end{ti}

\subsubsection{证明题（6 分）}
已知 $\balpha_{1},\balpha_{2},\balpha_{3},\balpha_{4}$ 是 $\bA \symbfit x = \symbfit 0$ 的一个基础解系, 证明 $\bbeta_{1} = \balpha_{1}, \bbeta_{2} = \balpha_{1} + 2\balpha_{2}, \bbeta_{3} = \balpha_{1} + 2\balpha_{3}, \bbeta_{4} = \balpha_{1} + 2\balpha_{4}$ 也是该方程组的一个基础解系

\section{难度与考试近似的题}
\subsubsection{填空题}
\begin{ti}
	已知 $n$ 阶矩阵 $\bA$ 各列元素之和为 $0$, 则 $|\bA| = $ \hua
\end{ti}

\begin{ti}
	三阶方阵 $\bA$ 的行列式为 $2$, 则 $|\bA^{\astt}| = $ \hua{}；$\bigl|\bA^{-1} + \bA^{\astt}\bigr| = $ \hua
\end{ti}

\begin{ti}
	齐次线性方程组 $\begin{cases}
		x_{1} + 2x_{3} = 0 \\
		2x_{2} - x_{4} = 0 \\
		x_{3} - x_{4} = 0
	\end{cases}$ 基础解系为 \hua
\end{ti}

\begin{ti}
	已知三阶方阵 $\bA$ 的特征值为 $1,2,3$, 则 $\bB = 3 \bE - 2 \bA$ 的特征值为 \hua
\end{ti}

\begin{ti}
	已知 $\bA = \begin{pmatrix}
		5 & 1 & 0 \\
		4 & 1 & 0 \\
		0 & 0 & 4
	\end{pmatrix}$, 则 $\bA^{-1} = $ \hua
\end{ti}

\begin{ti}
	若向量组 (1) 与向量组 (2) 可以相互线性表示, 则称向量组 (1) 与向量组 (2) \hua
\end{ti}

\begin{ti}
	设 $\bA,\bB$ 均为 $n$ 阶矩阵, 且 $\bB$ 为可逆矩阵, 若 $\bA \bB = \bB$, 则 $\bA = $ \hua
\end{ti}

\subsubsection{选择题}
\begin{ti}
	设方阵 $\bA = \begin{pmatrix}
		a & b \\
		c & d
	\end{pmatrix} (ad - bc \ne 0)$, 则 $\bA^{-1}$ 为 \kuo
	\fourch{$\frac{1}{ad - bc}\begin{pmatrix}
		d & -b \\
		-c & a
	\end{pmatrix}$}{$\frac{1}{bc - ad}\begin{pmatrix}
		d & -b \\
		-c & a
	\end{pmatrix}$}{$\frac{1}{ad - bc}\begin{pmatrix}
		-c & b \\
		d & -a
	\end{pmatrix}$}{$\frac{1}{bc - ad}\begin{pmatrix}
		d & -a \\
		-b & c
	\end{pmatrix}$}
\end{ti}

\begin{ti}
	向量组 $\balpha_{1} = (0,0,1,2,-1)^{\TT}, \balpha_{2} = (1,3,-2,2,-1)^{\TT}, \balpha_{3} = (2,6,-4,5,7)^{\TT}, \balpha_{4} = (-1,-3,4,0,-19)^{\TT}$ 的秩为 \kuo
	\fourch{1}{2}{3}{4}
\end{ti}

\begin{ti}
	设 $\bA$、$\bB$ 均为 $n$ 阶非零矩阵, 且 $\bA \bB = \symbfit 0$, 则 $\bA$ 和 $\bB$ 的秩 \kuo
	\twoch{必有一个等于 $0$}{都小于 $n$}{一个小于 $n$, 一个等于 $n$}{都等于 $n$}
\end{ti}

\begin{ti}
	对方程组 $\begin{cases}
		x_{1} - 2x_{2} + x_{3} = 1 \\
		-2x_{1} + 4x_{2} - 2x_{3} = 2
	\end{cases}$, \kuo{} 成立
	\fourch{有唯一解}{有两个解}{有无穷解}{无解}
\end{ti}

\begin{ti}
	设 $\bA$ 为正交阵, 则 $\bA^{-1}$ 为 \kuo
	\fourch{对称阵}{单位阵}{正交阵}{对角阵}
\end{ti}

\subsubsection{计算题}
\begin{ti}[8 分]
	计算行列式 $D = \begin{vmatrix}
		1 & 1 & 1 & 0 \\
		2 & 1 & -1 & 1 \\
		1 & 2 & -1 & 1 \\
		0 & 1 & 2 & 3
	\end{vmatrix}$
\end{ti}

\begin{ti}[10 分]
	求解矩阵方程 $\bX \bA = \bB$, 其中 $\bA = \begin{pmatrix}
		2 & 1 & -1 \\
		2 & 1 & 0 \\
		1 & -1 & 1
	\end{pmatrix}, B = \begin{pmatrix}
		1 & -1 & 3 \\
		4 & 3 & 2
	\end{pmatrix}$
\end{ti}

\begin{ti}[10 分]
	求解非齐次线性方程组 $\begin{cases}
		x_{1} - x_{2} + x_{3} = 1 \\
		2x_{1} - 2x_{2} + 5x_{3} = 2 \\
		2x_{1} + x_{2} + 12x_{3} = 0
	\end{cases}$
\end{ti}

\begin{ti}[12 分]
	求行向量组 $\balpha_{1} = (1,3,4,-2), \balpha_{2} = (2,1,3,-1), \balpha_{3} = (3,-1,2,0), \balpha_{4} = (4,-3,1,1)$ 的秩及一个极大无关组, 且求出其余向量由这一极大无关组线性表示的表达式
\end{ti}

\begin{ti}[13 分]
	设 $\bA = \begin{pmatrix}
		3 & -4 & -4 \\
		0 & 2 & 0 \\
		2 & -2 & -3
	\end{pmatrix}$, 求一可逆矩阵 $\symbfit P$, 使 $\symbfit P^{-1} \bA \symbfit P = \symbfit \varLambda$ 为对角阵
\end{ti}

\subsubsection{证明题（本题 8 分）}
设 $\lambda_{1}$、$\lambda_{2}$ 是方阵 $\bA$ 的特征值, 对应的特征向量分别为 $\symbfit p_{1}$、$\symbfit p_{2}$, 证明：若 $\lambda_{1} \ne \lambda_{2}$, 则 $\symbfit p_{1} + \symbfit p_{2}$ 不会是 $\bA$ 的特征向量

\newpage
\guanggao
	\chapter{《概率论与数理统计》试卷汇总}

\section{秋}
\subsection{2018-2019 B14}
\subsubsection{选择题（每题 $3$ 分, 共 $21$ 分）}
\begin{ti}
	从 $0,1,2,\cdots,9$ 中任意选出 $3$ 个不同的数字, 三个数字中不含 $0$ 与 $5$ 的概率是 \kuo{}
	\fourch{$\frac{1}{15}$}
	{$\frac{2}{15}$}
	{$\frac{14}{15}$}
	{$\frac{7}{15}$}
\end{ti}

\begin{ti}
	某人射击中靶的概率为 $\frac{3}{4}$ . 若射击直到中靶为止, 则射击次数为 $3$ 的概率为 \kuo{}
	\fourch{$\left(\frac{3}{4}\right)^3$}
	{$\left(\frac{1}{4}\right)^2\times\frac{3}{4}$}
	{$\left(\frac{1}{4}\right)^3$}
	{$\left(\frac{3}{4}\right)^2\times\frac{1}{4}$}
\end{ti}

\begin{ti}
	设随机变量 $X$ 的概率密度 $f(x)$ 满足 $f(-x)=f(x)$ , $F(x)$ 是分布函数, 则 \kuo{}
	\twoch{$F(-a)=1-F(a)$}
	{$F(-a)=\frac{1}{2}F(a)$}
	{$F(-a)=F(a)$}
	{$F(-a)=\frac{1}{2}-F(a)$}
\end{ti}

\begin{ti}
	设二维随机变量 $(X,Y)$ 的分布律为 $P\left\{X=i,Y=j\right\}=c\cdot i\cdot j,i=1,2,3,j=1,2,3$, 则 $c=$ \kuo{}
	\fourch{$\frac{1}{12}$}
	{$\frac{1}{3}$}
	{$\frac{1}{36}$}
	{$\frac{1}{2}$}
\end{ti}

\begin{ti}
	设随机变量 $X$ 服从均匀分布, 其概率密度为 $f(x)=
	\begin{cases}
	\frac{1}{2}, & 1<x<3\\
	0, & \text{其他}
	\end{cases}
	$, 则 $D(X)=$ \kuo{}
	\fourch{$3$}
	{$\frac{1}{3}$}
	{$\frac{1}{2}$}
	{$2$}
\end{ti}

\begin{ti}
	设总体 $X\sim N\left(0,\sigma^2\right)$, $X_1,X_2,\cdots,X_n$ 是总体 $X$ 的一个样本, $\overline{X},S^2$ 分别为样本均值和样本方差, 则下列样本函数中, 服从 $\chi^2(n)$ 分布的是 \kuo{}
	\fourch{$\sum_{i=1}^{n}X_i^2$}
	{$\frac{\overline{X}}{S/\sqrt{n-1}}$}
	{$\frac{(n-1)S^2}{\sigma^2}$}
	{$\frac{1}{\sigma^2}\sum_{i=1}^{n}X_i^2$}
\end{ti}

\begin{ti}
	设 $X_1,X_2,\cdots,X_n$ 是来自正态总体 $N\left(\mu,\sigma^2\right)$ 的一个样本, $\sigma^2$ 未知, $\overline{X}$是样本均值, $S^2=\frac{1}{n-1}\sum_{i=1}^{n}\Bigl(X_i-\overline{X}\Bigr)^2$, 如果 $\overline{X}-k\frac{S}{\sqrt{n}}$ 是 $\mu$ 的置信度为 $1-\alpha$ 的单侧置信下限, 则 $k$ 应取 \kuo{}
	\fourch{$t_{1-\alpha}(n)$}
	{$t_{\alpha}(n)$}
	{$t_{\alpha}(n-1)$}
	{$t_{\alpha/2}(n-1)$}
\end{ti}

\subsubsection{填空题（每题 $3$ 分, 共 $21$ 分）}
\begin{ti}
	设 $A,B$ 为随机事件, $P(A)=0.8$ , $P(A-B)=0.3$ , 则 $P\left(\overline{AB}\right)=$ \hua{}
\end{ti}

\begin{ti}
	设随机变量 $X$ 的分布律为 $P\left\{x=k\right\}=c(0.5)^k,k=1,2,3,\cdots$, 则常数 $c=$ \hua{}
\end{ti}

\begin{ti}
	设随机变量 $X$ 的概率密度为 $f(x)=
	\begin{cases}
	3x^2, & 0<x<1\\
	0, & \text{其他}
	\end{cases}
	$, 则 $P\left\{\left|X\right|<0.2\right\}=$ \hua{}
\end{ti}

\begin{ti}
	设随机变量 $X$ 的概率密度为 $f(x)=
	\begin{cases}
	\frac{1}{c}, & 0<x<c\\
	0, & \text{其他}
	\end{cases}
	$, 则 $\EE(X)=$ \hua{}
\end{ti}

\begin{ti}
	设二维随机变量 $(X,Y)$ 的概率密度为
	\begin{equation*}
		f(x,y)=
		\begin{cases}
		\sin x\cdot\cos y, & 0<x<\uppi/2,\ 0<y<\uppi/2\\
		0, & \text{其他}
		\end{cases},
	\end{equation*}
	则 $P\left\{0<X<\uppi/4,\uppi/4<Y<\uppi/2\right\}=$ \hua{}
\end{ti}

\begin{ti}
	设随机变量 $X$ 的数学期望 $\EE(X)=\mu$, 方差 $D(X)=\sigma^2$, 则由切比雪夫不等式有 $P\{|X-\mu|\geq3\sigma\}\leq$ \hua{}
\end{ti}

\begin{ti}
	设 $X_1,X_2$ 是取自正态总体 $X\sim N\left(\mu,\sigma^2\right)$ 的一个容量为 $2$ 的样本, 则 $\mu$ 的无偏估计量 $\hat\mu_1=\frac{1}{2}X_1+\frac{1}{2}X_2$, $\hat\mu_2=\frac{2}{3}X_1+\frac{1}{3}X_2$, $\hat\mu_3=\frac{1}{4}X_1+\frac{3}{4}X_2$ 中最有效的是 \hua{}
\end{ti}

\subsubsection{解答题（共 $58$ 分）}
\begin{ti}[$10$ 分]
	车间里有甲、乙、丙 $3$ 台机床生产同一种产品, 已知它们的次品率依次是 $0.05$、 $0.1$、 $0.2$, 产品所占份额依次是 $20\%$、 $30\%$、 $50\%$. 现从产品中任取 $1$ 件, 发现它是次品, 求次品来自机床乙的概率
\end{ti}

\begin{ti}[$10$ 分]
	设随机变量 $X$ 的分布函数为 $F(x)=
	\begin{cases}
	k-k\ee^{-x^3}, & x>0\\
	0, & x\leq0
	\end{cases}
	$, 试求:
	\begin{enumerate}
		\item 常数 $k$;
		\item $X$ 的概率密度 $f(x)$
	\end{enumerate}
\end{ti}

\begin{ti}[$10$ 分]
	设二维随机变量 $(X,Y)$ 的概率密度为:
	\begin{equation*}
		f(x,y)=
		\begin{cases}
		\frac{1}{4}, & 2\leqslant x\leq4,1\leqslant y\leq3\\
		0, & \text{其他}
		\end{cases},
	\end{equation*}
	试求 $(X,Y)$ 关于 $X$ 与 $Y$ 的边缘概率密度 $f_X(x)$ 与 $f_Y(y)$, 并判断 $X$ 与 $Y$ 是否相互独立
\end{ti}

\begin{ti}[$10$ 分]
	已知红黄两种番茄杂交的第二代结红果的植株与结黄果的植株的比率为 $3:1$, 现种植杂交种 $400$ 株, 试用中心极限定理近似计算, 结红果的植株介于 $285$ 与 $315$ 之间的概率. $\Bigl(\varPhi\Bigl(\sqrt{3}\Bigr)=0.9582,\allowbreak\varPhi\Bigl(\sqrt{2}\Bigr)=0.9207\Bigr)$
\end{ti}

\begin{ti}[$8$ 分]
	设二维随机变量 $(X,Y)$ 的分布律为
	\begin{center}
		\begin{tabularx}{0.8\textwidth}{ZZZZ}
			\hline
			 & \multicolumn{3}{c}{$Y$}\\
			\cline{2-4}
			$X$ & $-1$ & $0$ & $1$\\
			\hline
			$-1$ & $\frac{1}{8}$ & $\frac{1}{8}$ & $\frac{1}{8}$\\
			$0$ & $\frac{1}{8}$ & $0$ & $\frac{1}{8}$\\
			$1$ & $\frac{1}{8}$ & $\frac{1}{8}$ & $\frac{1}{8}$\\
			\hline
		\end{tabularx}
	\end{center}
	求 $\Cov(X,Y)$
\end{ti}

\begin{ti}[$10$ 分]
	设 $X_1,X_2,\cdots,X_n$ 为总体 $X$ 的一个样本, 总体 $X$ 的概率密度为:
	\begin{equation*}
		f(x)=
		\begin{cases}
		(\alpha+1)x^\alpha, & 0<x<1\\
		0, & \text{其他}
		\end{cases},
	\end{equation*}
	求未知参数 $\alpha$ 的矩估计
\end{ti}

\subsection{2018-2019 B14 答案}
\subsubsection{选择题（每题 $3$ 分, 共 $21$ 分）}
\begin{ti}
	从 $0,1,2,\cdots,9$ 中任意选出 $3$ 个不同的数字, 三个数字中不含 $0$ 与 $5$ 的概率是 \kuoD{}
	\fourch{$\frac{1}{15}$}
	{$\frac{2}{15}$}
	{$\frac{14}{15}$}
	{$\frac{7}{15}$}
\end{ti}

\begin{ti}
	某人射击中靶的概率为 $\frac{3}{4}$ . 若射击直到中靶为止, 则射击次数为 $3$ 的概率为 \kuoB{}
	\fourch{$\left(\frac{3}{4}\right)^3$}
	{$\left(\frac{1}{4}\right)^2\times\frac{3}{4}$}
	{$\left(\frac{1}{4}\right)^3$}
	{$\left(\frac{3}{4}\right)^2\times\frac{1}{4}$}
\end{ti}

\begin{ti}
	设随机变量 $X$ 的概率密度 $f(x)$ 满足 $f(-x)=f(x)$ , $F(x)$ 是分布函数, 则 \kuoA{}
	\twoch{$F(-a)=1-F(a)$}
	{$F(-a)=\frac{1}{2}F(a)$}
	{$F(-a)=F(a)$}
	{$F(-a)=\frac{1}{2}-F(a)$}
\end{ti}

\begin{ti}
	设二维随机变量 $(X,Y)$ 的分布律为 $P\left\{X=i,Y=j\right\}=c\cdot i\cdot j,i=1,2,3,j=1,2,3$, 则 $c=$ \kuoC{}
	\fourch{$\frac{1}{12}$}
	{$\frac{1}{3}$}
	{$\frac{1}{36}$}
	{$\frac{1}{2}$}
\end{ti}

\begin{ti}
	设随机变量 $X$ 服从均匀分布, 其概率密度为 $f(x)=
	\begin{cases}
	\frac{1}{2}, & 1<x<3\\
	0, & \text{其他}
	\end{cases}
	$, 则 $D(X)=$ \kuoB{}
	\fourch{$3$}
	{$\frac{1}{3}$}
	{$\frac{1}{2}$}
	{$2$}
\end{ti}

\begin{ti}
	设总体 $X\sim N\left(0,\sigma^2\right)$, $X_1,X_2,\cdots,X_n$ 是总体 $X$ 的一个样本, $\overline{X},S^2$ 分别为样本均值和样本方差, 则下列样本函数中, 服从 $\chi^2(n)$ 分布的是 \kuoD{}
	\fourch{$\sum_{i=1}^{n}X_i^2$}
	{$\frac{\overline{X}}{S/\sqrt{n-1}}$}
	{$\frac{(n-1)S^2}{\sigma^2}$}
	{$\frac{1}{\sigma^2}\sum_{i=1}^{n}X_i^2$}
\end{ti}

\begin{ti}
	设 $X_1,X_2,\cdots,X_n$ 是来自正态总体 $N\left(\mu,\sigma^2\right)$ 的一个样本, $\sigma^2$ 未知, $\overline{X}$是样本均值, $S^2=\frac{1}{n-1}\sum_{i=1}^{n}\Bigl(X_i-\overline{X}\Bigr)^2$, 如果 $\overline{X}-k\frac{S}{\sqrt{n}}$ 是 $\mu$ 的置信度为 $1-\alpha$ 的单侧置信下限, 则 $k$ 应取 \kuoC{}
	\fourch{$t_{1-\alpha}(n)$}
	{$t_{\alpha}(n)$}
	{$t_{\alpha}(n-1)$}
	{$t_{\alpha/2}(n-1)$}
\end{ti}

\subsubsection{填空题（每题 $3$ 分, 共 $21$ 分）}
\begin{ti}
	设 $A,B$ 为随机事件, $P(A)=0.8$ , $P(A-B)=0.3$ , 则 $P\left(\overline{AB}\right)=$ \huaa{$0.5$}
\end{ti}

\begin{ti}
	设随机变量 $X$ 的分布律为 $P\left\{x=k\right\}=c(0.5)^k,k=1,2,3,\cdots$, 则常数 $c=$ \huaa{$1$}
\end{ti}

\begin{ti}
	设随机变量 $X$ 的概率密度为 $f(x)=
	\begin{cases}
	3x^2, & 0<x<1\\
	0, & \text{其他}
	\end{cases}
	$, 则 $P\left\{\left|X\right|<0.2\right\}=$ \huaa{$\frac{1}{125}$}
\end{ti}

\begin{ti}
	设随机变量 $X$ 的概率密度为 $f(x)=
	\begin{cases}
	\frac{1}{c}, & 0<x<c\\
	0, & \text{其他}
	\end{cases}
	$, 则 $\EE(X)=$ \huaa{$\frac{c}{2}$}
\end{ti}

\begin{ti}
	设二维随机变量 $(X,Y)$ 的概率密度为
	\begin{equation*}
		f(x,y)=
		\begin{cases}
		\sin x\cdot\cos y, & 0<x<\uppi/2,\ 0<y<\uppi/2\\
		0, & \text{其他}
		\end{cases},
	\end{equation*}
	则 $P\left\{0<X<\uppi/4,\uppi/4<Y<\uppi/2\right\}=$ \huaa{$\left( \frac{2-\sqrt{2}}{2} \right)^2$}
\end{ti}

\begin{ti}
	设随机变量 $X$ 的数学期望 $\EE(X)=\mu$, 方差 $D(X)=\sigma^2$, 则由切比雪夫不等式有 $P\{|X-\mu|\geq3\sigma\}\leq$ \huaa{$\frac{1}{9}$}
\end{ti}

\begin{ti}
	设 $X_1,X_2$ 是取自正态总体 $X\sim N\left(\mu,\sigma^2\right)$ 的一个容量为 $2$ 的样本, 则 $\mu$ 的无偏估计量 $\hat\mu_1=\frac{1}{2}X_1+\frac{1}{2}X_2$, $\hat\mu_2=\frac{2}{3}X_1+\frac{1}{3}X_2$, $\hat\mu_3=\frac{1}{4}X_1+\frac{3}{4}X_2$ 中最有效的是 \huaa{$\hat\mu_1$}
\end{ti}


\subsubsection{解答题（共 $58$ 分）}
\begin{ti}[$10$ 分]
	车间里有甲、乙、丙 $3$ 台机床生产同一种产品, 已知它们的次品率依次是 $0.05$、 $0.1$、 $0.2$, 产品所占份额依次是 $20\%$、 $30\%$、 $50\%$. 现从产品中任取 $1$ 件, 发现它是次品, 求次品来自机床乙的概率
	\begin{solution}
		设抽取的产品为次品的事件为 $A$, 抽取的次品来自机床甲的事件为 $B_1$, 抽取的次品来自机床乙的事件为 $B_2$, 抽取的次品来自机床丙的事件为 $B_3$
		
		根据全概率公式
		\begin{equation*}
			\begin{aligned}
			P(A)&=P(A|B_1)P(B_1)+P(A|B_2)P(B_2)+P(A|B_3)P(B_3)\\
			&=0.05\times0.2+0.1\times0.3+0.2\times0.5=0.14
			\end{aligned}
		\end{equation*}
		根据贝叶斯公式
		\begin{equation*}
			P(B_2|A)=\frac{P(AB_2)}{P(A)}=\frac{P(A|B_2)P(B_2)}{P(A)}=\frac{0.1\times0.3}{0.14}=\frac{3}{14}
		\end{equation*}
	\end{solution}
\end{ti}

\begin{ti}[$10$ 分]
	设随机变量 $X$ 的分布函数为 $F(x)=
	\begin{cases}
	k-k\ee^{-x^3}, & x>0\\
	0, & x\leq0
	\end{cases}
	$, 试求:
	\begin{enumerate}
		\item 常数 $k$;
		\item $X$ 的概率密度 $f(x)$
	\end{enumerate}
	\begin{solution}
		\begin{enumerate}
		  \item 根据分布函数的性质
		  \[
			  \lim_{x\to+\infty}F(x)=k=1
		  \]
		  \item 
		  \[
			  F(x)=
		  	\begin{cases}
		  	1-\ee^{-x^3}, & x>0\\
		  	0, & x\leq0 
		  	\end{cases},
		  \]
		  则
		  \[
			  f(x)=F'(x)=
		  	\begin{cases}
		  	3x^2\ee^{-x^3}, & x>0\\
		  	0, & x\leq0
		  \end{cases}
		  \]
		\end{enumerate}
	  \end{solution}
\end{ti}

\begin{ti}[$10$ 分]
	设二维随机变量 $(X,Y)$ 的概率密度为:
	\begin{equation*}
		f(x,y)=
		\begin{cases}
		\frac{1}{4}, & 2\leqslant x\leq4,1\leqslant y\leq3\\
		0, & \text{其他}
		\end{cases},
	\end{equation*}
	试求 $(X,Y)$ 关于 $X$ 与 $Y$ 的边缘概率密度 $f_X(x)$ 与 $f_Y(y)$, 并判断 $X$ 与 $Y$ 是否相互独立
	\begin{solution}
		\[
		f_X(x)=\int_{-\infty}^{+\infty}f(x,y)\dd y=
		\begin{cases}
		\int_{1}^{3}\frac{1}{4}\dd y, & 2\leqslant x\leq4\\
		0, & \text{其它}
		\end{cases}=\begin{cases}
		\frac{1}{2}, & 2\leqslant x\leq4\\
		0, & \text{其它}
		\end{cases}
		\]
		同理
		\[
		f_Y(y)=
		\begin{cases}
		\frac{1}{2}, & 1\leqslant y\leq3\\
		0, & \text{其它}
		\end{cases},
		f_X(x)f_Y(y)=
		\begin{cases}
		\frac{1}{4}, & 2\leqslant x\leq4,1\leqslant y\leqslant 3\\
		0, & \text{其它}
		\end{cases}=f(x,y)
		\]
		因此 $X$ 与 $Y$ 相互独立
	\end{solution}
\end{ti}

\begin{ti}[$10$ 分]
	已知红黄两种番茄杂交的第二代结红果的植株与结黄果的植株的比率为 $3:1$, 现种植杂交种 $400$ 株, 试用中心极限定理近似计算, 结红果的植株介于 $285$ 与 $315$ 之间的概率. $\Bigl(\varPhi\Bigl(\sqrt{3}\Bigr)=0.9582,\allowbreak\varPhi\Bigl(\sqrt{2}\Bigr)=0.9207\Bigr)$
	\begin{solution}
		设结红果的植株的株数为 $X$, $X\sim B(400,3/4)$, 则 $\EE(X)=300$, $D(X)=75$, 根据中心极限定理
			\begin{align*}
			 P(285\leqslant X\leqslant 315)&=P\left(\frac{-15}{\sqrt{75}}\leq\frac{X-300}{\sqrt{75}}\leq\frac{15}{\sqrt{75}}\right)=\varPhi\left(\sqrt{3}\right)-\varPhi\left(-\sqrt{3}\right)\\
			&=2\varPhi\left(\sqrt{3}\right)-1=0.9164
			\end{align*}
	\end{solution}
\end{ti}

\begin{ti}[$8$ 分]
	设二维随机变量 $(X,Y)$ 的分布律为
	\begin{center}
		\begin{tabularx}{0.8\textwidth}{ZZZZ}
			\hline
			 & \multicolumn{3}{c}{$Y$}\\
			\cline{2-4}
			$X$ & $-1$ & $0$ & $1$\\
			\hline
			$-1$ & $\frac{1}{8}$ & $\frac{1}{8}$ & $\frac{1}{8}$\\
			$0$ & $\frac{1}{8}$ & $0$ & $\frac{1}{8}$\\
			$1$ & $\frac{1}{8}$ & $\frac{1}{8}$ & $\frac{1}{8}$\\
			\hline
		\end{tabularx}
	\end{center}
	求 $\Cov(X,Y)$
	\begin{solution}
		$\EE(X)=-1\times\frac{3}{8}+1\times\frac{3}{8}=0$ , 同理通过计算得 $\EE(Y)=0$ , $\EE(XY)=0$, 因此
		\[
			\Cov(X,Y)=\EE(XY)-\EE(X)\EE(Y)=0
		\]
	\end{solution}
\end{ti}

\begin{ti}[$10$ 分]
	设 $X_1,X_2,\cdots,X_n$ 为总体 $X$ 的一个样本, 总体 $X$ 的概率密度为:
	\begin{equation*}
		f(x)=
		\begin{cases}
		(\alpha+1)x^\alpha, & 0<x<1\\
		0, & \text{其他}
		\end{cases},
	\end{equation*}
	求未知参数 $\alpha$ 的矩估计
	\begin{solution}
		\[
			\EE(X)=\int_{0}^{1}(\alpha+1)x^{\alpha+1}\dd x=\frac{\alpha+1}{\alpha+2}, \mu_1=\overline{X}=\sum_{i=1}^{n}\frac{X_i}{n},
		\]
		因此
		\[
			\alpha=\frac{2\overline{X}-1}{1-\overline{X}}
		\]
	\end{solution}
\end{ti}

\subsection{2019-2020 A20}
\subsubsection{选择题}
\begin{ti}
	已知 $P(A) = \frac{1}{2}$, $P(B|A) = \frac{1}{4}$, 则 $P(AB) = $ \kuo
	\fourch{$\frac{1}{2}$}{$\frac{1}{4}$}{$\frac{1}{6}$}{$\frac{1}{8}$}
\end{ti}

\begin{ti}
	某类灯泡使用时数在 $500$ 小时以上的概率为 $0.5$, 从中取 $3$ 个灯泡使用, 则在使用 $500$ 小时后无一损坏的概率为 \kuo
	\fourch{$\frac{1}{8}$}{$\frac{3}{8}$}{$\frac{5}{8}$}{$\frac{7}{8}$}
\end{ti}

\begin{ti}
	如果抛一枚硬币 $4$ 次, 设 $X$ 表示正面朝上的次数, 则 $P\{ X = 3 \} = $ \kuo
	\fourch{$\frac{1}{2}$}{$\frac{1}{8}$}{$\frac{1}{4}$}{$\frac{1}{3}$}
\end{ti}

\begin{ti}
	设随机变量 $(X,Y)$ 概率密度为 $f(x,y) = \begin{cases}
		x^{2} + \frac{1}{3} xy, & 0 \leqslant x \leqslant 1, 0 \leqslant y \leqslant 2\\
		0, & \text{其他}.
	\end{cases}$ 则 $P\{ 0 \leqslant X \leqslant 1, 0 \leqslant Y \leqslant 1 \} = $ \kuo
	\fourch{$\frac{1}{12}$}{$\frac{3}{12}$}{$\frac{5}{12}$}{$\frac{7}{12}$}
\end{ti}

\begin{ti}
	设随机变量 $X \sim N(2,4)$, 且 $aX - 1 \sim N(0,1)$, 则 \kuo
	\fourch{$a = -2$}{$a = -0.5$}{$a = 0.5$}{$a = 2$}
\end{ti}

\begin{ti}
	设随机变量 $X$ 的数学期望 $\EE(X) = \mu$, 方差 $D(X) = \sigma^{2}$, 则由切比雪夫不等式有 $P\{ |X - \mu| < 2 \sigma \} \geqslant $ \kuo
	\fourch{$\frac{1}{2}$}{$\frac{1}{3}$}{$\frac{2}{3}$}{$\frac{3}{4}$}
\end{ti}

\begin{ti}
	设 $\hat{\theta}$ 为未知参数 $\theta$ 的估计量, 若 $E \bigl( \hat{\theta} \bigr)$ 存在, 且有 $E \bigl( \hat{\theta} \bigr) = \theta$, 则称 $\hat{\theta}$ 为参数 $\theta$ 的 \kuo
	\fourch{有偏估计量}{无偏估计量}{一致估计量}{有效估计量}
\end{ti}

\subsubsection{填空题}
\begin{ti}
	袋中装有标号 $1,2,3,4,5$ 的 $5$ 个球, 从中任取一个, 则取到标号大于 $3$ 的球的概率为 \hua
\end{ti}

\begin{ti}
	设 $A,B$ 为两个事件, $P(A) = 1$, $P(B) = 0.7$, $P\bigl( \overline{B}\bigl|A \bigr) = 0.3$, 则 $P(AB) = $ \hua
\end{ti}

\begin{ti}
	设随机变量 $X$ 的概率密度为 $f(x) = \begin{cases}
		x^{\alpha}, & 0 < x < 1 \\
		0 & \text{其他}
	\end{cases}$, 则 $\alpha = $ \hua
\end{ti}

\begin{ti}
	设二维随机变量 $(X,Y)$ 的概率密度为 $f(x,y) = \begin{cases}
		C x^{2} \sin y, & 0 < x < 1, 0 < y < \frac{\uppi}{2}, \\
		0, & \text{其他}.
	\end{cases}$ 则 $C = $ \hua
\end{ti}

\begin{ti}
	设随机变量 $X$ 的概率密度为 $f(x) = \begin{cases}
		\frac{1}{b - a}, & a < x < b, \\
		0, & \text{其他}.
	\end{cases}$ 则 $D(X) = $ \hua
\end{ti}

\begin{ti}
	设总体 $X$ 服从 0-1 分布 $B(1,p)$, 即 $P\{ X = 1 \} = p$ (其中 $p$ 为未知参数)。$X_{1},X_{2},\cdots,X_{n}$ 是来自 $X$ 的一个样本, 则样本函数 $X_{1} + X_{2}$, $\max_{1 \leqslant i \leqslant n} X_{i}$, $X_{n} + 2p$, $(X_{n} - X_{1})^{2}$ 中不是统计量的是 \hua
\end{ti}

\begin{ti}
	设 $\theta$ 和 $X_{1},X_{2},\cdots,X_{n}$ 分别是总体 $X$ 的未知参数和一个样本, $\theta_{1},\theta_{2}$ 是由样本确定的两个统计量, 满足 $P\{ \theta_{1} < \theta < \theta_{2} \} = 1 - \alpha$, 则称随机区间 $(\theta_{1},\theta_{2})$ 为 $\theta$ 的置信度为 \hua{} 的置信区间
\end{ti}

\subsubsection{解答题}
\begin{ti}
	有两批种子, 第一批发芽率为 $0.9$, 第二批发芽率为 $0.96$, 现将数量相同两批种子混合在一起, 从中任取一粒, 求该粒种子发芽的概率
\end{ti}

\begin{ti}
	设随机变量 $X$ 服从 $[0,2]$ 的均匀分布, 求随机变量 $Y = X^{2}$ 的概率密度
\end{ti}

\begin{ti}
	设二维随机变量 $(X,Y)$ 分布律如表~\ref{tab:1}, 求 $(X,Y)$ 关于 $X$ 与 $Y$ 的边缘分布律, 并判断 $X$ 与 $Y$ 是否相互独立
	\begin{table}[htbp]
		\centering
		\caption{}\label{tab:1}
		\begin{tabularx}{0.8\textwidth}{ZZZZ}
			\hline
			 & \multicolumn{3}{c}{$Y$}\\
			\cline{2-4}
			$X$ & 0 & 1 & 2 \\
			\hline
			0 & $\frac{1}{4}$ & $\frac{1}{4}$ & $\frac{1}{16}$\\
			1 & $\frac{1}{4}$ & $\frac{1}{8}$ & 0\\
			2 & $\frac{1}{16}$ & 0 & 0\\
			\hline
		\end{tabularx}
	\end{table}
\end{ti}

\begin{ti}
	设 $X$ 与 $Y$ 两个随机变量 $\EE(X) = 2$, $\EE(X^{2}) = 20$, $\EE(Y) = 3$, $\EE(Y^{2}) = 34$, $\Cov(X,Y) = 10$, 求 $\EE(X - Y)$, $D(X + Y)$
\end{ti}

\begin{ti}
	某工厂有 $200$ 台同类型机器, 由于工艺等原因, 每台机器工作时间只有 \SI{80}{\percent}, 各台机器相互独立, 求任一时刻有 $156$ 至 $164$ 台机器正在工作的概率($\varPhi(0.71) = 0.7611, \varPhi(1.98) = 0.9762$)
\end{ti}

\begin{ti}
	设 $X_{1},X_{2},\cdots,X_{n}$ 为总体 $X$ 的一个样本, 总体 $X$ 概率密度为 $f(x) = \begin{cases}
		\frac{2}{\theta} \ee^{- \frac{2x}{\theta}}, & x > 0, \\
		0, & x \leqslant 0.
	\end{cases}$ 求参数 $\theta(\theta > 0)$ 的极大似然估计
\end{ti}

\section{春}
\subsection{2018-2019 B}
\subsubsection{选择题（每题 $3$ 分，共 $21$ 分）}
\begin{ti}
	设随机变量 $X$ 分布律为 $P\{ X=k \}=pq^{k-1}(p>0,q>0,k=1,2,\cdots)$ , 则 \kuo{}
	\fourch{$p+q=1$}
	{$p=\frac{1}{q}-1$}
	{$p+q=2$}
	{$pq=1$}
\end{ti}

\begin{ti}
	从 $0,1,2,\cdots,9$ 这 $10$ 个数字中随机抽取 $1$ 个数字, 则取到奇数的概率为 \kuo{}
	\fourch{$\frac{1}{2}$}
	{$\frac{3}{5}$}
	{$\frac{4}{9}$}
	{$\frac{5}{9}$}
\end{ti}

\begin{ti}
	一批产品, 优质品占 \SI{20}{\percent}, 进行重复抽样检查, 共取 $5$ 件, 则恰好 $3$ 件是优质品的概率为 \kuo{}
	\fourch{$10\times 0.2^3$}
	{$0.2^3$}
	{$10\times 0.2^3\times 0.8^2$}
	{$0.2^3\times 0.8^2$}
\end{ti}

\begin{ti}
	同时掷骰子, 以 $X$ 和 $Y$ 分别代表第 $1$ 和第 $2$ 个骰子的点数, 则 $P\{ X+Y=5 \}=$ \kuo{}
	\fourch{$\frac{1}{12}$}
	{$\frac{1}{9}$}
	{$\frac{1}{36}$}
	{$\frac{1}{18}$}
\end{ti}

\begin{ti}
	设 $X$ 为一个随机变量, 且 $\EE(X)=2$, $\EE\left( X^2 \right)=5$, 则 $D(2X)=$ \kuo{}
	\fourch{$6$}
	{$12$}
	{$2$}
	{$4$}
\end{ti}

\begin{ti}
	设总体 $X\sim N\left( \mu,\sigma^2 \right)$, $X_1,\cdots,X_n$ 为样本, $\overline{X}$, $S^2$ 分别为样本均值和样本方差, 则下列选项服从 $\chi^2(n-1)$ 分布的是 \kuo{}
	\fourch{$\frac{(n-1)S^2}{\sigma^2}$}
	{$\frac{1}{\sigma^2}\sum_{i=1}^n \left(X_i-\mu\right)^2$}
	{$\frac{\overline{X}-\mu}{\sigma/\sqrt{n}}$}
	{$\frac{\overline{X}-\mu}{S/\sqrt{n}}$}
\end{ti}

\begin{ti}
	设 $\hat{\theta}_1$ 和 $\hat{\theta}_2$ 都是由样本确定的两个统计量, 满足 \kuo{}, 则随机区间 $\left(\hat{\theta}_1,\hat{\theta}_2\right)$ 称为 $\theta$ 的置信水平为 \SI{95}{\percent} 的置信区间
	\twoch{$P\left\{ \hat{\theta}_1<\theta<\hat{\theta}_2 \right\}=0.05$}
	{$P\left\{ \hat{\theta}_1<\theta<\hat{\theta}_2 \right\}=0.95$}
	{$P\left\{ \hat{\theta}_1<\theta \right\}=0.95$}
	{$P\left\{ \theta<\hat{\theta}_2 \right\}=0.95$}
\end{ti}

\subsubsection{填空题（每题 $3$ 分，共 $21$ 分）}
\begin{ti}
	设随机变量 $X$ 的概率密度为 $f(x)=
	\begin{cases}
		\frac{1}{x^2}, & a<x<3a,\\
		0, & \text{其他}.
	\end{cases}$, 则 $a=$\hua{}
\end{ti}

\begin{ti}
	设随机变量 $X$ 的数学期望 $\EE(X)=\mu$, 方差 $D(X)=\sigma^2$, 则由切比雪夫不等式有 $P\{ |X-\mu|<2\sigma \} \geq$\hua{}
\end{ti}

\begin{ti}
	设随机事件 $A$, $B$ 互不相容, 且 $P(A)=0.3$, $P\left( \overline{B} \right)=0.5$, 则 $P\left( B\middle|\overline{A} \right)=$\hua{}
\end{ti}

\begin{ti}
	设 $\hat{\theta}=\hat{\theta}(X_1,X_2,\cdots,X_n)$ 是 $\theta$ 的一个估计量, 若 $\EE\left( \hat{\theta} \right)$ 存在, 且有\hua{}, 则 $\hat{\theta}$ 是 $\theta$ 的无偏估计量
\end{ti}

\begin{ti}
	事件 $A$ 在一次实验中发生的概率为 $\frac{1}{3}$, $X$ 表示在 $3$ 次重复独立试验中发生的次数, 则 $P\{ X\leqslant 2 \}=$\hua{}
\end{ti}

\begin{ti}
	设二维随机变量 $(X,Y)$ 的概率密度为: $f(x,y)=
	\begin{cases}
		4xy, & 0<x<1,0<y<1,\\
		0, & \text{其他}.
	\end{cases}
	$, 则 $P\bigl\{ 0<X<\frac{1}{2},\frac{1}{4}<Y<1 \bigr\}=$ \hua{}
\end{ti}

\begin{ti}
	设随机变量 $X$ 的概率密度为 $f(x)=
	\begin{cases}
		\frac{1}{\beta-\alpha}, & \alpha<x<\beta,\\
		0, & \text{其他}.
	\end{cases}
	$, 则 $\EE(X)=$ \hua{}
\end{ti}

\subsubsection{解答题（共 $58$ 分）}
\begin{ti}[$8$ 分]
	设随机变量 $X$ 的分布律为
	\begin{center}
		\begin{tabularx}{0.8\textwidth}{Z|ZZZZZ}
			$X$ & $-1$ & $0$ & $1$ & $4$ & $7$\\
			\hline
			$P$ & $0.1$ & $0.2$ & $0.4$ & $0.16$ & $0.14$
		\end{tabularx}
	\end{center}
	求:
	\begin{enumerate}
		\item $P\{ |X-2|\leqslant 1 \}$
		\item $Y=X^2$ 的分布律
	\end{enumerate}
\end{ti}

\begin{ti}[$10$ 分]
	设某区域内肥胖者占 \SI{10}{\percent}, 不胖不瘦者占 \SI{82}{\percent}, 瘦者占 \SI{8}{\percent}, 肥胖者患高血压的概率为 \SI{30}{\percent}, 不胖不瘦者患高血压的概率为 \SI{10}{\percent}, 瘦者患高血压的概率为 \SI{5}{\percent}, 求:
	\begin{enumerate}
		\item 该地区的居民患高血压的概率
		\item 若在该地区任选一人, 发现有高血压, 属于肥胖者的概率
	\end{enumerate}
\end{ti}

\begin{ti}[$10$ 分]
	设二维随机变量 $(X,Y)$ 的概率密度为 $f(x,y)=
	\begin{cases}
		\ee^{-x}, & x>y,y>0,\\
		0, & \text{其他}.
	\end{cases}
	$, 试求: $(X,Y)$ 关于 $X$ 与 $Y$ 的边缘概率密度 $f_X(x)$ 与 $f_Y(y)$
\end{ti}

\begin{ti}[$10$ 分]
	设 $X\sim N(0,2)$, $Y\sim N(0,1)$, 且相互独立, $U=X+Y+1$, $V=X-Y+1$, 求 $\Cov(U,V)$
\end{ti}

\begin{ti}[$10$ 分]
	某通信系统有 $60$ 台相互独立起作用的交换机, 每台交换机能清晰接受信号的概率为 $0.90$. 系统工作时, 要求能清晰接受信号的交换机至少 $54$ 台, 求该通信系统能正常工作的概率(结果取最接近的值, 其中 $\varPhi(2.58)=0.9951$, $\varPhi(1.36)=0.9131$)
\end{ti}

\begin{ti}[$10$ 分]
	设 $X_1,X_2,\cdots,X_n$ 为总体 $X$ 的一个样本, 总体 $X$ 服从参数为 $p(0<p<1)$ 的两点分布, 其分布律为 $P\{ X=x \}=p^x (1-p)^{1-x}(x=0,1)$, 求未知参数 $p$ 的矩估计
\end{ti}

\subsection{2019-2020 A}
\subsubsection{选择题}
\begin{ti}
	设袋中有 4 只白球, 2 只黑球, 从袋中任取 2 球（不放回）, 取得 2 只白球的概率 \kuo{}
	\fourch{$\frac{3}{5}$}{$\frac{1}{5}$}{$\frac{2}{5}$}{$\frac{4}{5}$}
\end{ti}

\begin{ti}
	设随机变量 $X$ 服从正态分布 $N\bigl( 10,\sigma^2 \bigr)$, 且 $P\{ X < a \} = 0.5$, 则 $a =$ \kuo{}
	\fourch{$10$}{$10+\sigma$}{$\sigma$}{$10-\sigma$}
\end{ti}

\begin{ti}
	设二维随机变量 $(X,Y)$ 的分布律为
	\begin{center}
		\begin{tabularx}{0.8\textwidth}{Z*{2}{Z}}
			\hline
			 & \multicolumn{2}{c}{$Y$} \\
			\cline{2-3}
			$X$ & 1 & 2 \\
			\hline
			1 & $a$ & $\frac{1}{4}$ \\
			2 & $\frac{1}{4}$ & $b$ \\
			\hline
		\end{tabularx}
	\end{center}
	若 $X$ 与 $Y$ 相互独立, 则 \kuo{}
	\fourch{$a = \frac{1}{8}$, $b = \frac{3}{8}$}{$a = \frac{3}{8}$, $b = \frac{1}{8}$}{$a = \frac{1}{4}$, $b = \frac{1}{4}$}{$a = -\frac{1}{4}$, $b = \frac{3}{4}$}
\end{ti}

\begin{ti}
	设随机变量 $X$ 与 $Y$ 相互独立, 且 $D(X) = 3$, $D(Y) = 2$, 则 $D(2Y-X) =$ \kuo{}
	\fourch{1}{5}{7}{11}
\end{ti}

\begin{ti}
	设 $X$, $Y$ 为任意两个随机变量, 且 $\Cov(X,Y) = 0$, 则下述结论正确的是 \kuo{}
	\fourch{$D(2X-Y) = 2D(X) - D(Y)$}
	{$D(2X-Y) = 2D(X) + D(Y)$}
	{$D(2X-Y) = 4D(X) - D(Y)$}
	{$D(2X-Y) = 4D(X) + D(Y)$}
\end{ti}

\begin{ti}
	若从总体 $X$ 中抽出容量为 5 的一个样本, 样本观察值为 1, 3, 4, 4, 8, 则样本平均值为 \kuo{}
	\fourch{3}{4}{$4.2$}{$4.5$}
\end{ti}

\begin{ti}
	设总体 $X \sim B(1,p)$, 其中未知参数 $0<p<1$, $X_1$, $X_2$, $\cdots$, $X_n$ 是 $X$ 的一个样本, 则 $P$ 的矩估计为 \kuo{}
	\fourch{$\overline{X}$}{$\frac{1}{2} \overline{X}$}{$\frac{\overline{X}}{\overline{X}+1}$}{$1 - \overline{X}$}
\end{ti}

\subsubsection{填空题}
\begin{ti}
	已知 $A \subset B$, $P(A) = 0.2$, $P(B) = 0.3$, 则 $P\bigl( B \overline{A} \bigr) =$ \hua
\end{ti}

\begin{ti}
	甲、乙两人独立地进行射击, 甲击中的概率为 $0.9$, 乙击中的概率为 $0.8$, 则甲乙都击中的概率为 \hua
\end{ti}

\begin{ti}
	设随机变量 $X$ 的概率密度为 $f(x) = \begin{cases}
		\lambda \ee^{-3x}, & x > 0 \\
		0, & \text{其他}
	\end{cases}$, 则 $P\{ X \leqslant 0.1 \} =$ \hua
\end{ti}

\begin{ti}
	设二维随机变量 $(X,Y)$ 的概率密度 $f(x,y) = \begin{cases}
		Ax \ee^{-y}, & 0 < x < 2, y > 0 \\
		0, & \text{其他}
	\end{cases}$, 则 $A =$ \hua
\end{ti}

\begin{ti}
	设随机变量 $X$ 的概率密度为 $f(x) = \begin{cases}
		\ee^{-x}, & x > 0 \\
		0, & \text{其他}
	\end{cases}$, 则 $\EE(X) =$ \hua
\end{ti}

\begin{ti}
	设随机变量 $X$ 的数学期望 $\EE(X) = 4$, 方差 $D(X) = 5$, 则由切比雪夫不等式得 $P\{ |X - 4| < 3 \} \geqslant$ \hua
\end{ti}

\begin{ti}
	设总体 $X \sim N(\mu,1)$, $X_1$, $X_2$, $\cdots$, $X_n$ 为总体 $X$ 的一个样本, 则当 $C_1$, $\cdots$, $C_n$ 满足 \hua{} 时, $\sum_{i=1}^n C_i X_i$ 是 $\mu$ 的一个无偏估计量
\end{ti}

\subsubsection{计算题}
\begin{ti}
	两个箱子, 第一个箱子中有 4 个黑球, 一个白球, 第二个箱子中有 3 个黑球, 3 个白球. 先随机地取一个箱子, 再从这个箱子中取出一个球, 若已知取出的是白球, 问此球是第 2 个箱子的概率
\end{ti}

\begin{ti}
	设随机变量 $X$ 的分布律为
	\begin{center}
		\begin{tabular}{c|*{4}{c}}
			$X$ & 1 & 2 & $\cdots$ & 10 \\
			\hline
			$P$ & $2a$ & $4a$ & $\cdots$ & $20a$
		\end{tabular}
	\end{center}
	求:
	\begin{enumerate}
		\item 常数 $a$;
		\item $P\{ X>3 \}$
	\end{enumerate}
\end{ti}

\begin{ti}
	设二维随机变量 $(X,Y)$ 的概率密度为: $f(x,y) = \begin{cases}
		k (x+y)^2, & |x| \leqslant 1, |y| \leqslant 1 \\
		0, & \text{其他}
	\end{cases}$, 
	求:
	\begin{enumerate}
		\item $k$ 的值;
		\item $(X,Y)$ 关于 $X$ 与 $Y$ 的边缘概率密度 $f_X(x)$ 与 $f_Y(y)$
	\end{enumerate}
\end{ti}

\begin{ti}
	设二维随机变量 $(X,Y)$ 的概率密度为 $f(x,y) = \begin{cases}
		1, & 0 < x < 1, 0 < y < 1, \\
		0, & \text{其他}.
	\end{cases}$ 求 $D(X)$ 及 $D(Y)$
\end{ti}

\begin{ti}
	某保险公司第 $i$ 个月收到的保险费（万元）是随机变量 $X_i$, 且 $\EE(X_i) = 10$, $D(X_i) = 1$, 试用中心极限定理确定 100 个月收到的保险费超过 1010 万元的概率. （结果取最接近的值, 其中 $\varPhi(1) = 0.8413$, $\varPhi(1.5) = 0.9332$）
\end{ti}

\begin{ti}
	某种清漆的干燥时间 $X$（单位: 小时）服从正态分布 $N(\mu,\sigma^2)$, 随机抽取容量为 9 的一个样本, 测得其干燥时间的样本均值 $\overline{X} = 6$. 样本标准差 $S = 0.3$, 若 $\sigma^2$ 未知, 求 $\mu$ 的置信度为 \SI{95}{\percent} 的置信区间. （结果取最接近的值, 其中 $t_{0.025}(8) = 2.3060$, $t_{0.025}(9) = 2.2622$）
\end{ti}

\subsection{2020-2021 A}
只有重要的数据部分。
\subsubsection{选择题}
\begin{ti}
	$X \sim N(1,4)$, $D(2X-3) = $ \kuo
	\fourch{5}{8}{19}{16}
\end{ti}

\begin{ti}
	$f(x)$ 连续型 \kuo
	\fourch{
		$\int_{-\infty}^{+\infty} f(x) \dd{x} = 1$
	}{
		$\lim_{x \to +\infty} f(x) = 1$
	}{
		$0 \leqslant f(x) \leqslant 1$
	}{
		单调不减
	}
\end{ti}

\begin{ti}
	$X_1, \cdots, X_{16}$, $\cdots$ , $X \sim N(2,5^2)$, $\overline{X}$, 则 $\frac{\overline{X}^{-2}}{ \sigma/4 } \sim $ \kuo
	\fourch{$t(15)$}{$t(16)$}{$X^2(15)$}{$N(0,1)$}
\end{ti}

\begin{ti}
	$X \sim U(0, \theta)$, $\theta$ 的矩估计 \kuo
	\fourch{$\overline{X}$}{$2\overline{X}$}{$2\overline{X}-1$}{$2\overline{X}+1$}
\end{ti}

\begin{ti}
	6 只晶体管有一只次品, 取 2 只取得一次一正的概率 \kuo
	\fourch{$\frac{7}{15}$}{$\frac{1}{3}$}{$\frac{2}{3}$}{$\frac{3}{4}$}
\end{ti}

\begin{ti}
	$F(x,y)$, $X$ 的边缘分布函数 $F_X(x)$ \kuo
	\fourch{
		$F_X(x) = F(x, +\infty)$
	}{
		$F_X(x) = F(+\infty, y)$
	}{
		$F_X(x) = 1 - F(x, +\infty)$
	}{
		$F_X(x) = 1 - F(+\infty, y)$
	}
\end{ti}

\subsubsection{填空题}
\begin{ti}
	$P(A) + P(B) = 0.8$, $P(AB) = 0.3$, 则 $P \bigl( \overline{A \cup B} \bigr) = $ \hua
\end{ti}

\begin{ti}
	$f(x) = \begin{cases}
		\lambda \ee^{-3x}, 	& x \geqslant 0 \\
		0, 									& x < 0
	\end{cases}$, $\lambda = $ \hua
\end{ti}

\begin{ti}
	$f(x,y) = \begin{cases}
		1,	& 0<x<1,0<y<1	\\
		0,
	\end{cases}$, $P\bigl\{ 0 < x < \frac{1}{2}, -1 < y < 1 \bigr\} = $ \hua
\end{ti}

\begin{ti}
	$f(x) = \begin{cases}
		\frac{1}{b}, & -2<x<4, \\
		0,					 & \text{其它},
	\end{cases}$, $D(x) = $ \hua
\end{ti}

\begin{ti}
	$\EE(X) = 0.1$, $\EE(Y) = 10$, $X$, $Y$ 独立, $\EE(Y-4XY) = $ \hua
\end{ti}

\begin{ti}
	$\EE(X) = M$, $D(X) = \sigma^2$, $P\bigl\{ |x - \mu| < a \bigr\} \geqslant \frac{3}{4}$, $a = $ \hua
\end{ti}

\begin{ti}
	$\theta_1$, $\theta_2$ 无偏估计量, 若 \hua{} 则 $\theta_1$ 较 $\theta_2$ 有效
\end{ti}

\subsubsection{计算题}
\begin{ti}
	一仓库有10箱同种规格的产品, 其中由甲, 乙, 丙三厂生产的分别为5箱, 3箱, 2箱, 三厂产品的次品率依次为0.1, 0.2, 0.3, 从这10箱产品中任取一箱, 再从这箱中任取一件, 求取得正品的概率? 若确实取得正品, 求正品由甲厂生产的概率
\end{ti}

\begin{ti}
	$f(x) = \begin{cases}
		0,	& x < 0 \\
		a \sin x,	& 0 \leqslant x \leqslant \frac{\uppi}{2}	\\
		1,	&	x > \frac{\uppi}{2}
	\end{cases}$, 求
	\begin{enumerate}
		\item $a$
		\item $P \bigl\{ |x| < \frac{\uppi}{6} \bigr\}$
	\end{enumerate}
\end{ti}

\begin{ti}
	$f(x,y) = \begin{cases}
		k(6-x-y), & 0<x<2, 2<y<3,	\\
		0,	&	\text{其它},
	\end{cases}$
	\begin{enumerate}
		\item $k$
		\item 边缘概率 $f_X(x), f_Y(y)$
	\end{enumerate}
\end{ti}

\begin{ti}
	\mbox{}
	\begin{center}
		\begin{tabular}{cccccc}
			$x$	&	1	&	2	&	3	&	4	&	5	\\
			$P$	&	$0.05$	&	0.20	&	0.35	&	0.30	&	0.1
		\end{tabular}
	\end{center}
	求 $\EE(X)$, $D(X)$
\end{ti}

\begin{ti}
	木柱 100 根, 其中 \SI{80}{\percent} 长度不小于 \SI{3}{m}, 问其中至少 30 根小于 \SI{3}{m} 的概率 ($\Phi(2.5) = 0.9938$)
\end{ti}

\begin{ti}
	磁棒 $X$ 服从 $N(10,\sigma^2)$, $\overline{X} = 15$, $s^2 = 0.05$, $\sigma^2$ 未知. $\mu$ 置信 \SI{95}{\percent} 的置信区间. ($t_{0.025}(4) = 2.7764$)
\end{ti}

\section{拼凑卷}
此复习题非一份完整的考试卷, 而是多个不同的试卷拼凑（老师给的）。

\subsubsection{选择题（每题 $3$ 分）}
\begin{ti}
	已知事件 $A,B$ 满足 $P(AB)=P\left( \overline{A}\bigcap\overline{B} \right)$, 且 $P(A)=0.4$, 则 $P(B)=$ \kuo{}
	\fourch{0.4}
	{0.5}
	{0.6}
	{0.7}
\end{ti}

\begin{ti}
	有 $\gamma$ 个球, 随机地放在 $n$ 个盒子中 $(\gamma \leqslant n)$, 则某指定的 $\gamma$ 个盒子中各有一球的概率为 \kuo{}
	\fourch{$\frac{\gamma !}{n^{\gamma}}$}
	{$\mathrm{C}_n^r\frac{\gamma !}{n^\gamma}$}
	{$\frac{n!}{\gamma^n}$}
	{$\mathrm{C}_\gamma^n\frac{n!}{\gamma^n}$}
\end{ti}

\begin{ti}
	设随机变量 $X$ 的概率密度为 $f(x)=c\ee^{-\left|x\right|}$, 则 $c=$ \kuo{}
	\fourch{$-\frac{1}{2}$}{$0$}{$\frac{1}{2}$}{$1$}
\end{ti}

\begin{ti}
	掷一颗骰子 $600$ 次, 求“一点”出现次数的均值为 \kuo{}
	\fourch{$50$}
	{$100$}
	{$120$}
	{$150$}
\end{ti}

\begin{ti}
	设每次试验成功的概率为 $p (0<p<1)$, 重复进行试验直到第 $n$ 次才取得 $r (1\leqslant r\leqslant n)$ 次成功的概率为 \kuo{}
	\fourch{$\mathrm{C}_{n-1}^{r-1}p^r(1-p)^{n-r}$}
	{$\mathrm{C}_{n}^{r}p^r(1-p)^{n-r}$}
	{$\mathrm{C}_{n-1}^{r-1}p^{r-1}(1-p)^{n-r+1}$}
	{$p^r(1-p)^{n-r}$}
\end{ti}

\begin{ti}
	离散型随机变量 $X$ 的分布函数为 $F(x)$, 则 $P(X=x_k)=$ \kuo{}
	\fourch{$P(x_{k-1}\leqslant X\leqslant x_k)$}
	{$F(x_{k+1})-F(x_{k-1})$}
	{$P(x_{k-1}<X<x_{k+1})$}
	{$F(x_{k})-F(x_{k-1})$}
\end{ti}

\begin{ti}
	设随机变量 $X, Y$ 是相互独立的两个随机变量, 其分布函数分别为 $F_X(x), F_Y(y)$, 则随机变量 $Z=\max (X,\allowbreak Y)$ 的分布函数为 \kuo{}
	\twoch{$F_Z(z)=\max\left\{F_X(x),F_Y(y)\right\}$}
	{$F_Z(z)=\max\left\{\left|F_X(x)\right|,\left|F_Y(y)\right|\right\}$}
	{$F_Z(z)=F_X(x)F_Y(y)$}
	{$F_Z(z)=F_X(z)F_Y(z)$}
\end{ti}

\begin{ti}
	设随机变量 $(X, Y)$ 的方差 $D(X)=4, D(Y)=1$, 相关系数 $\rho_{XY}=0.6$, 则方差 $D(3X-2Y)=$ \kuo{}
	\fourch{$40$}
	{$34$}
	{$25.6$}
	{$17.6$}
\end{ti}

\begin{ti}
	设 $(X_1,X_2,\cdots,X_n)$ 为总体 $N\left(1,2^2\right)$ 的一个样本, $\overline X$ 为样本均值, 则下列结论中正确的是 \kuo{}
	\twoch{$\frac{\overline{X}-1}{2/\sqrt{n}}\sim t(n)$}
	{$\frac{1}{4}\sum_{i=1}^{n}(X_i-1)^2\sim F(n,1)$}
	{$\frac{\overline{X}-1}{\sqrt{2}/\sqrt{n}}\sim N(0,1)$}
	{$\frac{1}{4}\sum_{i=1}^{n}(X_i-1)^2\sim\chi^2(n)$}
\end{ti}

\begin{ti}
	设总体 $X$ 在 $(\mu-\rho,\mu+\rho)$ 上服从均匀分布, 则参数 $\mu$ 的矩估计量为 \kuo{}
	\fourch{$\frac{1}{\overline{X}}$}
	{$\frac{1}{n-1}\sum_{i=1}^{n}X_i$}
	{$\frac{1}{n-1}\sum_{i=1}^{n}X_i^2$}
	{$\overline{X}$}
\end{ti}

\begin{ti}
	设二维随机变量 $(X,Y)$ 的分布律为:
	\begin{center}
		\begin{tabularx}{0.8\textwidth}{ZZZ}
			\hline
			 & \multicolumn{2}{c}{$Y$}\\
			\cline{2-3}
			$X$ & 1 & 2\\
			\hline
			1 & $a$ & $\frac{2}{9}$\\
			\hline
			2 & $b$ & $\frac{4}{9}$\\
			\hline
		\end{tabularx}
	\end{center}
	若 $X$ 与 $Y$ 相互独立, 则 \kuo{}
	\fourch{$a=\frac{4}{9}, b=\frac{1}{9}$}
	{$a=\frac{1}{9}, b=\frac{4}{9}$}
	{$a=\frac{2}{9}, b=\frac{1}{9}$}
	{$a=\frac{1}{9}, b=\frac{2}{9}$}
\end{ti}

\begin{ti}
	设随机变量 $X$ 的概率密度为 $f(x)=
	\begin{cases}
	ax, & 0\leqslant x\leqslant 2\\
	0, & \text{其他}
	\end{cases}
	$, 则 $a=$ \kuo{}
	\fourch{$1$}
	{$\frac{1}{4}$}
	{$\frac{1}{2}$}
	{$\frac{1}{3}$}
\end{ti}

\begin{ti}
	设 $\EE(X)=2$, $\EE(Y)=3$, 则 $\EE(3X-4Y+5)=$ \kuo{}
	\fourch{1}
	{6}
	{$-6$}
	{$-1$}
\end{ti}

\begin{ti}
	设随机变量 $X\sim N(1,4)$, 则下列随机变量中服从标准正态分布的是 \kuo{}
	\fourch{$Y=\frac{X-1}{4}$}
	{$Y=\frac{X-1}{2}$}
	{$Y=\frac{X+1}{4}$}
	{$Y=\frac{X+1}{2}$}
\end{ti}

\begin{ti}
	设盒中有 $4$ 支铅笔, $2$ 支钢笔, 从盒中任取 $2$ 支笔(不放回抽样), 则取得 $1$ 支铅笔和 $1$ 支钢笔的概率是 \kuo{}
	\fourch{$\frac{8}{15}$}
	{$\frac{4}{5}$}
	{$\frac{3}{5}$}
	{$\frac{7}{15}$}
\end{ti}

\begin{ti}
	设总体 $X$ 服从正态分布 $N(\mu,1)$, $X_1$, $X_2$ 是来自总体 $X$ 的一个样本, $\hat\mu_1=\frac{2}{3}X_1+\frac{1}{3}X_2$, 
	$\hat\mu_2=\frac{1}{4}X_1+\frac{3}{4}X_2$, $\hat\mu_3=\frac{1}{2}X_1+\frac{1}{2}X_2$ 都是 $\mu$ 的无偏估计量, 则其中方差最小的是 \kuo{}
	\fourch{$\hat\mu_3$}
	{$\hat\mu_2$}
	{$\hat\mu_1$}
	{一样大}
\end{ti}

\begin{ti}
	设随机变量 $X\sim\chi^2(m_1)$, $Y\sim\chi^2(m_2)$, 且 $X$ 与 $Y$ 相互独立, 则下列结果中不正确的是 \kuo{}
	\twoch{$X+Y\sim\chi^2(m_1+m_2)$}
	{$D(Y)=m_2$}
	{$D(X)=2m_1$}
	{$\EE(X)=m_1$}
\end{ti}

\subsubsection{填空题（每题 3 分）}
\begin{ti}
	已知 $P(B)=0.3$, $P\left(\overline{A}\bigcup B\right)=0.7$, 且 $A$ 与 $B$ 相互独立, 则 $P(A)=$ \hua{}
\end{ti}

\begin{ti}
	设随机变量 $X$ 服从参数为 $\lambda$ 的泊松分布, 且 $P\left\{X=0\right\}=\frac{1}{3}$, 则 $\lambda=$ \hua{}
\end{ti}

\begin{ti}
	设 $X\sim N\left(2,\sigma^2\right)$, 且 $P\left\{2<X<4\right\}=0.2$, 则 $P\left\{X<0\right\}=$ \hua{}
\end{ti}

\begin{ti}
	已知 $D(X)=2$, $D(Y)=1$, 且 $X$ 和 $Y$ 相互独立, 则 $D(X-2Y)=$ \hua{}
\end{ti}

\begin{ti}
	设 $S^2$ 是从 $N(0,1)$ 中抽取容量为 $16$ 的样本方差, 则 $D\left(S^2\right)=$ \hua{}
\end{ti}

\begin{ti}
	一批电子元件共有 $100$ 个, 次品数为 $5$. 连续两次不放回地从中任取一个, 则第二次才取得正品的概率为 \hua{}
\end{ti}

\begin{ti}
	设事件 $A$ 与 $B$ 相互独立, $P(A)=0.2$, $P(B)=0.3$, 则 $P\left(A\bigcup B\right)=$ \hua{}
\end{ti}

\begin{ti}
	设 $\overline{X}$ 为总体 $X\sim N(3,4)$ 中抽取的样本 $(X_1,X_2,X_3,X_4)$ 的均值, 则 $P\left(-1<\overline{X}<5\right)=$ \hua{}
\end{ti}

\begin{ti}
	甲、乙两人独立地进行射击, 甲击中的概率为 $0.9$, 乙击中的概率为 $0.8$, 则甲中乙不中的概率等于 \hua{}
\end{ti}

\begin{ti}
	设 $P(A)=P(B)=P(C)=\frac{1}{3}$, $ P(AB)=P(AC)=0$, $ P(BC)=\frac{1}{5}$. 则 $A, B, C$ 中至少有一个发生的概率为 \hua{}
\end{ti}

\begin{ti}
	设二维随机变量 $(X,Y)$ 的概率密度为 $f(x,y)=
	\begin{cases}
	 C\ee^{-3x}\sin y, & x>0, 0<y<\frac{\uppi}{2}\\
	0, & \text{其他}
	\end{cases}
	$, 则 $C=$ \hua{}
\end{ti}

\begin{ti}
	设随机变量 $X$ 服从参数为 $3$ 的指数分布, 则 $P\left\{X\geqslant  1\right\}=$ \hua{}
\end{ti}

\begin{ti}
	总体 $Y\sim N\left(\mu,\sigma^2\right)$, $\overline{Y}=\frac{1}{n}\sum_{i=1}^{n}Y_i$ 为样本均值, $S$ 为样本标准差, 当 $\sigma$ 为未知时, $\mu$ 的置信度为 $1-\alpha (0<\alpha<1)$ 的双侧置信区间为 \hua{}
\end{ti}

\begin{ti}
	设 $X,Y$ 为两个随机变量, 已知 $\EE(X)=1$, $\EE(Y)=2$, $\EE(XY)=5$, 则 $\Cov(X,Y-4)=$ \hua{}
\end{ti}

\begin{ti}
	若 $P(A)=0.5$, $P\left(B\overline{A}\right)=0.2$, 则 $P(A+B)=$ \hua{}
\end{ti}

\begin{ti}
	已知随机变量 $X\sim
	\begin{bmatrix}
	-1 & 0 & 2 & 5\\
	0.25 & 0.25 & 0.25 & 0.25
	\end{bmatrix}
	$, 那么 $\EE(X)=$ \hua{}
\end{ti}

\begin{ti}
	设 $\hat\theta$ 是未知参数 $\theta$ 的一个无偏估计量, 则 $\EE(\hat\theta)=$ \hua{}
\end{ti}

\begin{ti}
	随机变量 $X$ 服从均匀分布 $U(1,3)$, 则 $P(X>2)=$ \hua{}
\end{ti}

\begin{ti}
	设随机变量 $X\sim B(100,0.15)$, 则 $\EE(X)=$ \hua{}
\end{ti}

\begin{ti}
	设随机变量 $X\sim N(3,4)$, 已知 $\varPhi(1)=0.8413$, 则 $P(X<1)=$ \hua{}
\end{ti}

\begin{ti}
	设随机变量 $X$ 的概率密度函数为 $f(x)=
	\begin{cases}
	3x^2, & 0\leqslant x\leqslant 1\\
	0, & \text{其它}
	\end{cases}
	$, 则$ P\left(X<\frac{1}{2}\right)=$ \hua{}
\end{ti}

\begin{ti}
	设随机变量 $X\sim
	\begin{bmatrix}
	0 & 1 & 3\\
	0.5 & 0.35 & 0.15
	\end{bmatrix}
	$, 则 $P(X<2)=$ \hua{}
\end{ti}

\begin{ti}
	设随机变量 $X$ 的期望存在, 则 $\EE(X-\EE(X))=$ \hua{}
\end{ti}

\begin{ti}
	设 $X$ 为随机变量, 已知 $D(X)=2$, 那么 $D(3X-5)=$ \hua{}
\end{ti}

\subsubsection{计算题}
\begin{ti}[10 分]
	设随机变量 $X$ 与 $Y$ 具有概率密度: $f(x,y)=
	\begin{cases}
	\frac{1}{8}(x+y) & 0\leqslant x\leqslant 2, 0\leqslant y\leqslant 2\\
	0 & \text{其它}
	\end{cases}
	$. 试求: $D(X)$, $D(Y)$ 与 $D(2X-3Y)$
\end{ti}

\begin{ti}[10 分]
	某电子计算机主机有 $100$ 个终端, 每个终端有 \SI{80}{\percent} 的时间被使用. 若各个终端是否被使用是相互独立的, 试求至少有 $15$ 个终端空闲的概率. $(\varPhi(1.25)=0.8944, \varPhi(0.31)=0.6217$
\end{ti}

\begin{ti}[10 分]
	试求正态总体 $N\left(\mu,0.5^2\right)$ 的容量分别为 $10,15$ 的两独立样本均值差的绝对值大于 $0.4$ 的概率. $(\varPhi(1.96)=0.975)$
\end{ti}

\begin{ti}[10 分]
	设总体 $X$ 的密度函数为 $f(x)=
	\begin{cases}
	\frac{2}{\theta^2}(\theta-x), & 0<x<\theta\\
	0, & \text{其它}
	\end{cases}
	$, $\theta>0$, $X_1,X_2,\cdots,X_{10}$ 为来自总体 $X$ 的样本, 试求当样本观察值分别为 $0.5$, $1.3$, $0.6$, $1.7$, $2.2$, $1.2$, $0.8$, $1.5$, $2.0$, $1.6$ 时未知参数 $\theta$ 的矩估计值
\end{ti}

\begin{ti}[10 分]
	某商店拥有某产品共计 $12$ 件, 其中 $4$ 件次品, 已经售出 $2$ 件, 现从剩下的 $10$ 件产品中任取一件, 求这件是正品的概率
\end{ti}

\begin{ti}[10 分]
	设某种电子元件的寿命服从正态分布 $N(40,100)$, 随机地取 $5$ 个元件, 求恰有两个元件寿命小于 $50$ 的概率. $(\varPhi(1)=0.8413, \varPhi(2)=0.9772)$
\end{ti}

\begin{ti}[12 分]
	设总体 $X$ 的分布律为 $P\left\{X=k\right\}=(1-p)^{k-1}p, k=1,2,\cdots.$ ($p$ 为未知参数), $X_1,X_2,\cdots,X_n$ 是总体 $X$ 的一个样本, 求 $p$ 的极大似然估计量
\end{ti}

\begin{ti}[10 分]
	两台车床加工同样的零件, 第一台出现不合格品的概率是 $0.03$, 第二台出现不合格品的概率是 $0.06$, 加工出来的零件放在一起, 并且已知第一台加工的零件数比第二台加工的零件数多一倍.
	\begin{enumerate}
		\item 求任取一个零件是合格品的概率.
		\item 如果取出的零件是不合格品, 求它是由第二台车床加工的概率
	\end{enumerate}
\end{ti}

\begin{ti}[10 分]
	某仪器装了 $3$ 个独立工作的同型号电子元件, 其寿命(单位: 小时)都服从同一指数分布, 密度函数为 $f(x)=
	\begin{cases}
	\frac{1}{600}\ee^{-\frac{x}{600}}, & x>0\\
	0, & \text{其它}
	\end{cases}
	$, 试求此仪器在最初使用的 $200$ 小时内, 至少有一个此种电子元件损坏的概率
\end{ti}
	\chapter{复变函数试卷汇总}

\section{复习题 1}
\subsubsection{选择题(每小题 $3$ 分, 共 $15$ 分)}
\begin{enumerate}
	\item $\frac{(\sqrt{3}-\ii)^{4}}{(1-\ii)^{8}}=$ (\hspace{1pc})
	\twoch{$-\frac{1}{2}+\frac{\sqrt{3}}{2}\ii$}{$-\frac{1}{8}\left(1+\sqrt{3}\ii\right)$}{$\frac{1}{8}\left(-1+\sqrt{3} \ii\right)$}{$-\frac{1}{2}-\frac{\sqrt{3}}{2} \ii$}
	
	\item 设 $f(z)=2 x^{3}+3 y^{3} \ii$ , 则 $f(z)$ (\hspace{1pc})
	\twoch{处处不可导}{仅在 $6x^2=9y^2$ 上可导, 处处不解析}{处处解析}{仅在 $(0,0)$ 点可导}
	
	\item 下列等式正确的是 (\hspace{1pc})
	\twoch{$\Ln \mathrm{i}=\left(2 k \uppi-\frac{\uppi}{2}\right) \ii, \ln \ii=\frac{\uppi}{2} \ii$}{$\Ln \ii=\left( 2k\uppi+\frac{\uppi}{2}\right)\ii,\ln\ii=-\frac{\uppi}{2}\ii $}{$\Ln \ii=\left(2 k \uppi+\frac{\uppi}{2}\right) \ii, \ln \ii=\frac{\uppi}{2} \ii$}{$\Ln \ii=\left(2 k \uppi-\frac{\uppi}{2}\right) \ii, \ln \ii=-\frac{\uppi}{2} \ii$}
	
	\item $z=0$ 是函数 $\frac{1-\cos z}{z-\sin z}$ 的 (\hspace{1pc})
	\fourch{本性奇点}{可去奇点}{二级极点}{一级极点}
	
	\item 设 $\mathrm{C}$ 为 $z=(1-\ii)t$ , $t$ 从 $1$ 到 $0$ 的一段, 则 $\int_{\mathrm{C}} \overline{z} \dd z=$ (\hspace{1pc})
	\fourch{$-1$}{$1$}{$-\ii$}{$\ii$}
\end{enumerate}

\subsubsection{填空题(每小题 $3$ 分, 共 $15$ 分)}
\begin{enumerate}
	\item 若 $z+|z|=2+\ii$ , 则 $z=$\underline{\hspace{8pc}}
	
	\item 若 $\mathrm{C}$ 为正向圆周 $|z|=\frac{1}{2}$ , 则 $\oint_{\mathrm{C}} \frac{1}{z-2} \dd z=$\underline{\hspace{8pc}}
	
	\item 若 $z=2-\uppi\ii$ , 则 $\ee^{z}=$\underline{\hspace{8pc}}
	
	\item 若 $f(z)=\cos z^2$ , 则 $f(z)$ 在 $z=0$ 处泰勒展开式中 $z^4$ 项的系数 $a_4=$\underline{\hspace{8pc}}
	
	\item 函数 $f(t)=\sin t$ 的拉普拉斯变换 $F(s)=$\underline{\hspace{8pc}}
\end{enumerate}

\subsubsection{计算题(70分)}
\begin{enumerate}
	\item 设 $u(x,y)=x-2xy$ 且 $f(0)=0$ , 求解析函数 $f(z)=u+\ii v$ . ( $10$ 分)
	
	\item 计算积分 $\oint_{\mathrm{C}}\frac{2\ee^x}{z^5}\dd z$ 的值, 其中 $\mathrm{C}$ 为正向圆周 $|z|=1$ . ( $7$ 分)
	
	\item 计算积分 $\oint_{\mathrm{C}}\frac{3z+5}{z^2-z}\dd z$ 的值, 其中 $\mathrm{C}$ 为正向圆周 $|z|=\frac{1}{2}$ . ( $7$ 分)
	
	\item 求函数 $\frac{1-\cos z}{z^3}$ 在有限奇点处的留数. ( $7$ 分)
	
	\item 求函数 $\frac{2z^2+1}{z^2+2z}$ 在有限奇点处的留数. ( $7$ 分)
	
	\item 将 $f(z)=\frac{z}{(z-2)(z-6)}$ 在 $2<|z|<6$ 内展开为洛朗级数. ( $10$ 分)
	
	\item 若函数 $f(z)=a y^{3}+b x^{2} y+\ii\left(x^{3}+c x y^{2}\right)$ 是复平面上的解析函数, 求 $a,b,c$ 的值. ( $12$ 分)
	
	\item 利用拉普拉斯变换解常微分方程初值问题: $\begin{cases}
	x''(t)+6x'(t)+9x(t)=\ee^{-3t}\\
	x(0)=0, x'(0)=0
	\end{cases}$ . ( $10$ 分)
\end{enumerate}



\section{复习题 1 答案}
\subsubsection{选择题(每小题 $3$ 分, 共 $15$ 分)}
\begin{enumerate}
	\item $\frac{(\sqrt{3}-\ii)^{4}}{(1-\ii)^{8}}=$ (\hspace{0.25pc}D\hspace{0.25pc})
	\twoch{$-\frac{1}{2}+\frac{\sqrt{3}}{2}\ii$}{$-\frac{1}{8}\left(1+\sqrt{3}\ii\right)$}{$\frac{1}{8}\left(-1+\sqrt{3} \ii\right)$}{$-\frac{1}{2}-\frac{\sqrt{3}}{2} \ii$}
	
	\item 设 $f(z)=2 x^{3}+3 y^{3} \ii$ , 则 $f(z)$ (\hspace{0.25pc}B\hspace{0.25pc})
	\twoch{处处不可导}{仅在 $6x^2=9y^2$ 上可导, 处处不解析}{处处解析}{仅在 $(0,0)$ 点可导}
	
	\item 下列等式正确的是 (\hspace{0.25pc}C\hspace{0.25pc})
	\twoch{$\Ln \mathrm{i}=\left(2 k \uppi-\frac{\uppi}{2}\right) \ii, \ln \ii=\frac{\uppi}{2} \ii$}{$\Ln \ii=\left( 2k\uppi+\frac{\uppi}{2}\right)\ii,\ln\ii=-\frac{\uppi}{2}\ii $}{$\Ln \ii=\left(2 k \uppi+\frac{\uppi}{2}\right) \ii, \ln \ii=\frac{\uppi}{2} \ii$}{$\Ln \ii=\left(2 k \uppi-\frac{\uppi}{2}\right) \ii, \ln \ii=-\frac{\uppi}{2} \ii$}
	
	\item $z=0$ 是函数 $\frac{1-\cos z}{z-\sin z}$ 的 (\hspace{0.25pc}D\hspace{0.25pc})
	\fourch{本性奇点}{可去奇点}{二级极点}{一级极点}
	
	\item 设 $\mathrm{C}$ 为 $z=(1-\ii)t$ , $t$ 从 $1$ 到 $0$ 的一段, 则 $\int_{\mathrm{C}} \overline{z} \dd z=$ (\hspace{0.25pc}A\hspace{0.25pc})
	\fourch{$-1$}{$1$}{$-\ii$}{$\ii$}
\end{enumerate}

\subsubsection{填空题(每小题 $3$ 分, 共 $15$ 分)}
\begin{enumerate}
	\item 若 $z+|z|=2+\ii$ , 则 $z=$\underline{\hspace{1pc}$\frac{3}{4}+\ii$\hspace{1pc}}
	
	\item 若 $\mathrm{C}$ 为正向圆周 $|z|=\frac{1}{2}$ , 则 $\oint_{\mathrm{C}} \frac{1}{z-2} \dd z=$\underline{\hspace{1pc}$0$\hspace{1pc}}
	
	\item 若 $z=2-\uppi\ii$ , 则 $\ee^{z}=$\underline{\hspace{1pc}$-\ee^2$\hspace{1pc}}
	
	\item 若 $f(z)=\cos z^2$ , 则 $f(z)$ 在 $z=0$ 处泰勒展开式中 $z^4$ 项的系数 $a_4=$\underline{\hspace{1pc}$-\frac{1}{2}$\hspace{1pc}}
	
	\item 函数 $f(t)=\sin t$ 的拉普拉斯变换 $F(s)=$\underline{\hspace{1pc}$\frac{1}{s^2+1}$\hspace{1pc}}
\end{enumerate}

\subsubsection{计算题(70分)}
\begin{enumerate}
	\item 设 $u(x,y)=x-2xy$ 且 $f(0)=0$ , 求解析函数 $f(z)=u+\ii v$ . ( $10$ 分)
	\begin{solution}
		解析函数的 $u,v$ 必定满足 $\mathrm{C}.-\mathrm{R}.$ 方程, 即
		\begin{equation*}
			\begin{cases}
			\frac{\partial u}{\partial x}=\frac{\partial v}{\partial y}\\
			\frac{\partial u}{\partial y}=-\frac{\partial v}{\partial x}
			\end{cases}
		\end{equation*}
		$\frac{\partial v}{\partial y}=\frac{\partial u}{\partial x}=1-2 y$ , $\frac{\partial v}{\partial y}$ 对 $y$ 积分得 $v=y-y^{2}+\varphi(x)$
		
		$\frac{\partial u}{\partial y}=-2 x=-\frac{\partial v}{\partial x}=-\varphi^{\prime}(x)$ , 可以得出 $\varphi(x)=x^{2}+C$
		
		由于 $f(0)=0$ , 因此 $C=0$ ,即 $f(z)=x-2 x y+\ii\left(y-y^{2}+x^{2}\right)$
	\end{solution}
	
	\item 计算积分 $\oint_{\mathrm{C}}\frac{2\ee^x}{z^5}\dd z$ 的值, 其中 $\mathrm{C}$ 为正向圆周 $|z|=1$ . ( $7$ 分)
	\begin{solution}
		根据高阶导数公式 $f^{(n)}(z_0)=\frac{n!}{2\uppi\ii}\oint_{\mathrm{C}}\frac{f(z)}{(z-z_0)^{n+1}}\dd z$ , 那么
		\begin{equation*}
			\oint_{\mathrm{C}} \frac{2 \ee^{z}}{(z-0)^{5}} \dd z=\frac{2 \uppi \ii}{4 !}\left.\left(2 \ee^{z}\right)^{(4)}\right|_{z=0}=\frac{\uppi \ii}{6}
		\end{equation*}
	\end{solution}
	
	\item 计算积分 $\oint_{\mathrm{C}}\frac{3z+5}{z^2-z}\dd z$ 的值, 其中 $\mathrm{C}$ 为正向圆周 $|z|=\frac{1}{2}$ . ( $7$ 分)
	\begin{solution}
		\begin{equation*}
			\oint_{\mathrm{C}}\frac{3z+5}{z^2-z}\dd z=2\uppi\ii\underset{z=0}{\Res}\frac{3z+5}{z(z-1)}=2\uppi\ii\left.\frac{3z+5}{z-1}\right|_{z=0}=-10\uppi\ii
		\end{equation*}
	\end{solution}
	
	\item 求函数 $\frac{1-\cos z}{z^3}$ 在有限奇点处的留数. ( $7$ 分)
	\begin{solution}
		对 $\cos z$ 进行洛朗展开, $\cos z=1+\sum_{n=1}^{\infty}(-1)^n\frac{z^{2n}}{(2n)!}$ , 那么 $1-\cos z=\sum_{n=1}^{\infty}(-1)^{n+1}\frac{z^{2n}}{(2n)!}$
		
		那么 $\frac{1-\cos z}{z^3}=\sum_{n=1}^{\infty}(-1)^{n+1}\frac{z^{2n-3}}{(2n)!}$ , 根据洛朗系数公式, $\underset{z=0}{\Res}\frac{1-\cos z}{z^3}=c_{-1}=\frac{1}{2}$
	\end{solution}
	
	\item 求函数 $\frac{2z^2+1}{z^2+2z}$ 在有限奇点处的留数. ( $7$ 分)
	\begin{solution}
		\begin{equation*}
			\underset{z=0}{\Res}\frac{2z^2+1}{z^2+2z}=\left.\frac{2z^2+1}{z+2} \right|_{z=0}=\frac{1}{2} , 
			\underset{z=-2}{\Res}\frac{2z^2+1}{z^2+2z}=\left.\frac{2z^2+1}{z}\right|_{z=-2}=-\frac{9}{2}
		\end{equation*}
		
	\end{solution}
	
	\item 将 $f(z)=\frac{z}{(z-2)(z-6)}$ 在 $2<|z|<6$ 内展开为洛朗级数. ( $10$ 分)
	\begin{solution}
		\begin{align*}
			f(z)&=\frac{z}{4}\left( \frac{1}{z-6}-\frac{1}{z-2}\right) =\frac{z}{4}\left( -\frac{1}{6}\frac{1}{1-z/6}-\frac{1}{z}\frac{1}{1-2/z}\right) \\
			&=\frac{z}{4}\left( -\frac{1}{6}\sum_{n=0}^{\infty}(z/6)^n-\frac{1}{z}\sum_{n=0}^{\infty}(2/z)^n\right)\\
			&=-\frac{1}{4}\left( \sum_{n=0}^{\infty}(z/6)^{n+1}+\sum_{n=0}^{\infty}(2/z)^n\right)  
		\end{align*}
	\end{solution}
	
	\item 若函数 $f(z)=a y^{3}+b x^{2} y+\ii\left(x^{3}+c x y^{2}\right)$ 是复平面上的解析函数, 求 $a,b,c$ 的值. ( $12$ 分)
	\begin{solution}
		若 $f(z)$ 为解析函数, 则其实部、虚部满足 $\mathrm{C}.-\mathrm{R}.$ 方程, 设 $u=ay^3+bx^2y$ , $v=x^3+cxy^2$ , 则有
		\begin{equation*}
			\begin{cases}
			\frac{\partial u}{\partial x}=2 b x y=2 c x y=\frac{\partial v}{\partial y}\\
			\frac{\partial u}{\partial y}=3 a y^{2}+b x^{2}=-3 x^{2}-c y^{2}=-\frac{\partial v}{\partial x}
			\end{cases}
		\end{equation*}
		解得\begin{equation*}
			\begin{cases}
			a=1\\
			b=c=-3
			\end{cases}
		\end{equation*}
		
	\end{solution}
	
	\item 利用拉普拉斯变换解常微分方程初值问题: $\begin{cases}
	x''(t)+6x'(t)+9x(t)=\ee^{-3t}\\
	x(0)=0, x'(0)=0
	\end{cases}$ . ( $10$ 分)
	\begin{solution}
		设 $\LL[x]=X(s)$ , 对等式两边作拉普拉斯变换
		\begin{align*}
			\LL[x''+6x'+9x]&=s^2X(s)-sx(0)-x'(0)+6sX(s)-6x(0)+9X(s)\\
			&=s^2X(s)+6sX(s)+9X(s)=\frac{1}{s+3}
		\end{align*}
		那么有 $X(s)=\frac{1}{(s+3)^3}$ , 根据拉普拉斯变换的微分性质 $F''(s)=\LL[t^2f(t)]$
		\begin{equation*}
			\frac{1}{(s+3)^3}=\frac{1}{2}\left(\frac{1}{s+3} \right)''=\frac{\LL[t^2\ee^{-3t}]}{2}
		\end{equation*}
		那么 $x(t)=\frac{t^2\ee^{-3t}}{2}$
	\end{solution}
\end{enumerate}
	\backmatter
	\chapter{后记}
首先向一路披荆斩棘做到这里的读者表示祝贺，至少在精神上你已经成
为一名合格的江理学子。从此你生是江理的人，死是江理的鬼。

面对后面纷繁复杂的课程，守望“三实”的初心。同学们要带着“为人诚实、基础扎实、工作踏实”的精神气质，真诚对待每一门课，仔细倾听老师的讲课，尊重课程就是尊重自己，爱学习是最大的爱自己。愿“三实”品质伴随同学们漫漫课程路，创造理想的成绩。
\begin{flushright}
	二〇二〇年六月于赣州
\end{flushright}
\end{document}